\chapter{Algebraic curves, function fields and Riemann surfaces}\label{ch:curves}

The paper's objects live on curves: the modular curve $X_0(N)$, the elliptic curve $E$, the modular parametrisation $\varphi_E\colon X_0(N)\to E$ and its degree, the
cohomology $H^1(X_0(N),\Zl)$ with its cup product, and the Jacobian $J_0(N)$ with its torsion and Tate modules. This chapter supplies the theory of curves that these
words presuppose, in two languages: algebraic (varieties, function fields, divisors, differentials, Riemann-Roch, morphisms and their degree, Riemann-Hurwitz), in the
classical style of Silverman's book, and analytic (compact Riemann surfaces, holomorphic forms, periods, $H_1$, $H^1$ and the cup product, the Jacobian as a complex
torus), joined by Riemann's existence theorem. It ends with correspondences and the theorem that the adjoint of a correspondence for the cup product is its
transpose, which is what makes the Hecke operators self-adjoint for the pairing of Section 4.1 of the paper. It uses \Cref{ch:groups} (inverse limits, $\Zl$, Moebius
maps), \Cref{ch:rings} (polynomial rings, finite fields), \Cref{ch:modules} (Nakayama, discrete valuation rings, lattices, saturation, perfect and alternating pairings,
adjoints, idempotents), \Cref{ch:reps} (contragredients), \Cref{ch:galois} (separability, Hilbert 90) and \Cref{ch:ant} (the identity $\sum e_if_i=n$), and the complex
analysis of \Cref{ch:prereq}.

\section*{Core results at a glance}

Throughout, $k$ is a perfect field, $C$ a smooth projective curve over $k$ of genus $g$, $X$, $Y$ compact Riemann surfaces, $\varphi\colon X\to Y$ a nonconstant morphism of
degree $n$, $\langle\ ,\ \rangle$ the cup product, $A^\dagger$ the adjoint of an operator, and $\mathcal C=(C,\alpha,\beta)$ a correspondence with transpose $\mathcal C^t=(C,\beta,\alpha)$.

\begin{enumerate}[label=\textbf{\thechapter.\arabic*}]
\item \textbf{Conventions.} $G_k$ acts on points coordinatewise; $\varphi^*$ is pullback, $\varphi_*$ pushforward, $A^\dagger$ an adjoint; cup product $\int_X\theta\wedge\eta$ with the complex
orientation; $\pi_q$ the Frobenius morphism; $l(D)$ the dimension of $\mathcal L(D)$. \Cref{conv:curves:conventions}.
\item \textbf{Affine varieties} (definition): zero sets $V(S)\subseteq\A^n(\bar k)$ with $I(V)$ prime; coordinate ring, function field, dimension $=$ transcendence degree; defined over $k$.
\Cref{def:curves:affine-variety}.
\item \textbf{Nullstellensatz.} $I(V(\mf a))=\sqrt{\mf a}$; for $k$ perfect, $V$ is defined over $k$ iff $\sigma(V)=V$ for all $\sigma\in G_k$. \emph{(stated; reference)}
\Cref{thm:curves:nullstellensatz}.
\item \textbf{Projective varieties} (definition): $\PP^n=(\bar k^{n+1}\setminus0)/\bar k^\times$, homogeneous ideals, charts $U_i\simeq\A^n$, projective closure, function field of quotients of forms
of equal degree. \Cref{def:curves:projective-space}.
\item \textbf{Points over subfields.} $\PP^n(\bar k)^{G_L}=\PP^n(L)$; $\PP^n=\A^n\sqcup\PP^{n-1}$; the function field of a projective variety is that of any affine chart. \emph{(proved
here)} \Cref{prop:curves:projective-points}.
\item \textbf{Curves, rational maps, morphisms} (definition): a curve is a projective variety of dimension one defined over $k$; morphisms are rational maps regular everywhere; defined
over $k$ iff Galois-equivariant. \Cref{def:curves:curve-morphism}.
\item \textbf{Plane curves} (definition): $V(F)$, $F$ homogeneous and irreducible over $\bar k$; $P$ singular iff $F_X=F_Y=F_Z=0$ at $P$. \Cref{def:curves:plane-curve}.
\item \textbf{Affine criterion.} $XF_X+YF_Y+ZF_Z=dF$; an affine point is singular iff $f=f_x=f_y=0$; tangent line $F_X(P)X+F_Y(P)Y+F_Z(P)Z=0$. \emph{(proved here)}
\Cref{prop:curves:euler}.
\item \textbf{Smooth points are discrete valuations.} $P$ smooth iff $\dim\mf m_P/\mf m_P^2=1$ iff $\bar k[C]_P$ is a discrete valuation ring; a line other than the tangent gives a uniformiser;
$\ord_P$ is the order of vanishing. \emph{(proved here for plane models; general case stated)} \Cref{thm:curves:smooth-dvr}.
\item \textbf{Bezout.} Plane curves of degrees $m$, $n$ without common component meet in $mn$ points with multiplicity. \emph{(stated; reference)} \Cref{thm:curves:bezout}.
\item \textbf{Smooth models.} Every curve is birational over $k$ to a unique smooth projective curve; its points are the valuations of the function field. \emph{(stated; reference)}
\Cref{thm:curves:smooth-model}.
\item \textbf{Function fields} (definition): finitely generated of transcendence degree one with $k$ algebraically closed in it. \Cref{def:curves:function-field}.
\item \textbf{Valuations.} Each point gives $\ord_P$; a function has finitely many zeros and poles; a function without poles is constant. \emph{(proved here, given the finiteness of
proper closed subsets)} \Cref{prop:curves:function-field}.
\item \textbf{Extension of morphisms.} A rational map from a smooth curve to a projective variety is a morphism; every $f\in\bar k(C)$ is a morphism $C\to\PP^1$. \emph{(proved here)}
\Cref{thm:curves:extend-morphism}.
\item \textbf{Curves and function fields.} $C\mapsto k(C)$ is a contravariant equivalence between smooth projective curves with nonconstant morphisms over $k$ and function fields of one
variable over $k$; a morphism of degree one is an isomorphism. \emph{(stated; reference)} \Cref{thm:curves:equivalence}.
\item \textbf{Galois action.} $C(\bar k)^{G_L}=C(L)$; $(\sigma f)(\sigma P)=\sigma(f(P))$ and $\ord_{\sigma P}\sigma f=\ord_Pf$; morphisms over $k$ commute with $G_k$. \emph{(proved here; the fixed
field $k(C)$ stated)} \Cref{prop:curves:galois-action}.
\item \textbf{Divisors} (definition): $\Div(C)$, degree, $\Div_k(C)=\Div(C)^{G_k}$, $\operatorname{div}(f)=\sum\ord_P(f)(P)$. \Cref{def:curves:divisor}.
\item \textbf{Principal divisors.} $\operatorname{div}f=0$ iff $f$ is constant; $\deg\operatorname{div}f=0$, with $\deg\operatorname{div}_0f=\deg\operatorname{div}_\infty f=[\bar k(C):\bar k(f)]$; on $\PP^1$,
$\operatorname{div}(x-a)=(a)-(\infty)$. \emph{(proved here, given \Cref{thm:curves:degree-formula})} \Cref{prop:curves:principal}.
\item \textbf{Picard group} (definition): $\Pic(C)=\Div(C)/\operatorname{Prin}(C)$, $\Pic^0(C)$ the degree-zero part. \Cref{def:curves:picard}.
\item \textbf{The Picard group.} $\Pic^0(\PP^1)=0$; $G_k$ acts; $\Div^0_k(C)/\operatorname{div}(k(C)^\times)\to\Pic^0(C)^{G_k}$ is injective \emph{(proved here)}, and bijective if $C(k)\ne\emptyset$
\emph{(stated; reference)}. \Cref{prop:curves:picard}.
\item \textbf{Differentials} (definition): $\Omega_C$, $\ord_P\omega$, $\operatorname{div}\omega$, regular differentials. \Cref{def:curves:differentials}.
\item \textbf{Differentials and the canonical class.} $\Omega_C$ is one-dimensional, $\ord_P$ is well defined \emph{(stated; reference)}; all $\operatorname{div}\omega$ form one class $K_C$; on $\PP^1$,
$\operatorname{div}(dx)=-2(\infty)$; on $y^2=h(x)$, $\operatorname{div}(dx/y)=0$ \emph{(proved here)}. \Cref{prop:curves:differentials}.
\item \textbf{Riemann-Roch spaces} (definition): $\mathcal L(D)=\{f:\operatorname{div}f+D\ge0\}\cup\{0\}$, $l(D)=\dim\mathcal L(D)$. \Cref{def:curves:riemann-roch-space}.
\item \textbf{Riemann-Roch spaces.} $\mathcal L(0)=\bar k$; $\mathcal L(D)=0$ if $\deg D<0$; $l(D+(P))\le l(D)+1$, so $l(D)\le\deg D+1$ for $D\ge0$; $l$ depends on the class; for $D$ over $k$,
$\mathcal L(D)$ has a $k$-rational basis \emph{(proved here; the rational basis stated)}. \Cref{prop:curves:riemann-roch-space}.
\item \textbf{Riemann-Roch.} $l(D)-l(K_C-D)=\deg D-g+1$. \emph{(stated; reference, with the shape of the proof)} \Cref{thm:curves:riemann-roch}.
\item \textbf{Consequences.} $l(K_C)=g$; $\deg K_C=2g-2$; $l(D)=\deg D-g+1$ if $\deg D>2g-2$; genus $0$ with a rational point means $\PP^1$; for $g=1$, $l(D)=\deg D$ ($\deg D\ge1$), $K_C\sim0$, and
regular differentials have no zeros. \emph{(proved here)} \Cref{cor:curves:rr-consequences}.
\item \textbf{Plane curves.} A smooth plane curve of degree $d$ has genus $(d-1)(d-2)/2$, in any characteristic. \emph{(proof sketched here)} \Cref{thm:curves:plane-genus}.
\item \textbf{Weierstrass equations.} A genus-one curve $C$ with $O\in C(k)$ is isomorphic over $k$, by $[x:y:1]$ with $x\in\mathcal L(2(O))$, $y\in\mathcal L(3(O))$, to a smooth Weierstrass cubic
$y^2+a_1xy+a_3y=x^3+a_2x^2+a_4x+a_6$, $a_i\in k$, with $O\mapsto[0:1:0]$; conversely. \emph{(proved here, given that images of projective varieties are closed;
the converse in characteristic $2$ rests on the sketched \Cref{thm:curves:plane-genus})} \Cref{thm:curves:genus-one-weierstrass}.
\item \textbf{Genus one is $\Pic^0$.} For $g=1$ and $O\in C(k)$, $P\mapsto[(P)-(O)]$ is a $G_k$-equivariant bijection $C(\bar k)\to\Pic^0(C)$. \emph{(proved here)}
\Cref{prop:curves:genus-one-pic}.
\item \textbf{Nonconstant morphisms.} A morphism of curves is constant or surjective \emph{(stated; reference)}; $\varphi^*$ embeds $\bar k(C_2)$ in $\bar k(C_1)$ with finite index \emph{(proved here)}.
\Cref{thm:curves:constant-or-surjective}.
\item \textbf{Degree, ramification, pullback} (definition): $\deg\varphi=[\bar k(C_1):\varphi^*\bar k(C_2)]$; $e_\varphi(P)=\ord_P(\varphi^*t_{\varphi(P)})$; $\varphi^*(Q)=\sum_{\varphi(P)=Q}e_\varphi(P)(P)$; $\varphi_*(P)=(\varphi(P))$.
\Cref{def:curves:degree}.
\item \textbf{Degree formula.} $\sum_{P\mapsto Q}e_\varphi(P)=\deg\varphi$ for every $Q$. \emph{(proof sketched here)} \Cref{thm:curves:degree-formula}.
\item \textbf{The degree in three ways.} Field degree; $\#\varphi^{-1}(Q)=\deg\varphi$ off the branch points for separable $\varphi$; $\deg\varphi^*D=\deg\varphi\deg D$, $\varphi_*\varphi^*=\deg\varphi$,
$\varphi^*\operatorname{div}f=\operatorname{div}\varphi^*f$, so $\varphi_*\varphi^*=[\deg\varphi]$ on $\Pic^0$; degrees multiply; degree one means isomorphism; a function without poles is constant.
\emph{(proved here, given \Cref{thm:curves:degree-formula}; the norm formula and the inseparable case stated)} \Cref{cor:curves:degree-three-ways}.
\item \textbf{Separability and Frobenius.} Characteristic $0$ morphisms are separable \emph{(proved here)}; separable iff $\varphi^*\omega\ne0$, and the factorisation through Frobenius
\emph{(stated; reference)}; $\pi_q$ is a purely inseparable bijection of degree $q$ with $e=q$ everywhere, fixed points $C(\F_q)$, acting on points as $\Frob_q$ \emph{(proved here; the degree
sketched)}. \Cref{prop:curves:separability-frobenius}.
\item \textbf{Ramification} (definition): unramified, branch points, tame and wild. \Cref{def:curves:ramification}.
\item \textbf{Local computation.} For separable $\varphi$, $\ord_P\varphi^*\omega=e\,\ord_Q\omega+\ord_P(d(\varphi^*t_Q)/dt_P)$ with the last term $\ge e-1$, equality iff tame; finitely many ramified points.
\emph{(proved here, given \Cref{prop:curves:separability-frobenius}(b) and \Cref{thm:curves:degree-formula})} \Cref{lem:curves:local-differential}.
\item \textbf{Riemann-Hurwitz.} $2g_1-2\ge\deg\varphi(2g_2-2)+\sum_P(e_\varphi(P)-1)$ for separable $\varphi$, with equality iff all ramification is tame. \emph{(proved here, given \Cref{thm:curves:degree-formula} and Riemann-Roch)}
\Cref{thm:curves:riemann-hurwitz}.
\item \textbf{Consequences.} $g_1\ge g_2$ if $g_2\ge1$; morphisms of genus-one curves are unramified; hyperelliptic and Fermat genera; for $\varphi\colon X\to E$, $\sum(e_P-1)=2g(X)-2$
(for $\varphi_{\curve{121a1}}$: $10$; $\varphi_{\curve{539c1}}$: $96$). \emph{(proved here, given \Cref{thm:curves:riemann-hurwitz})} \Cref{cor:curves:rh-consequences}.
\item \textbf{Riemann surfaces} (definition): charts with holomorphic transitions; meromorphic functions; holomorphic maps. \Cref{def:curves:riemann-surface}.
\item \textbf{Holomorphic maps.} Local normal form $w=z^e$; open, surjective, finite fibres; constant degree $\sum_{P\mapsto Q}\operatorname{mult}_P$; finitely many ramification points;
orientation-preserving. \emph{(proved here)} \Cref{thm:curves:holomorphic-maps}.
\item \textbf{Riemann existence and GAGA.} Compact Riemann surfaces and smooth projective complex curves are equivalent, compatibly with functions, orders, multiplicities, degrees,
differentials, divisors and genus. \emph{(stated; reference)} \Cref{thm:curves:riemann-existence}.
\item \textbf{Topology.} A compact Riemann surface is a sphere with $g$ handles, $\chi=2-2g$ \emph{(stated; reference)}; $\chi(X)=d\chi(Y)-\sum(\operatorname{mult}_P-1)$ \emph{(proof sketched here)}.
\Cref{thm:curves:topology}.
\item \textbf{$1$-forms} (definition): meromorphic and holomorphic $1$-forms, residues, $H^0(X,\Omega^1)$, smooth forms and Stokes. \Cref{def:curves:holomorphic-differential}.
\item \textbf{Residue theorem.} $\sum_P\operatorname{res}_P\omega=0$. \emph{(proved here)} \Cref{thm:curves:residue}.
\item \textbf{Argument principle.} $\deg\operatorname{div}f=\sum\operatorname{res}(df/f)=0$. \emph{(proved here)} \Cref{cor:curves:argument-principle}.
\item \textbf{Genus and forms.} $\dim H^0(X,\Omega^1)=g$; $\deg\operatorname{div}\omega=2g-2$. \emph{(proved here from Riemann-Roch and GAGA, both stated)} \Cref{thm:curves:genus-dimension}.
\item \textbf{Bilinear relations.} $\int\omega\wedge\eta=0$ and $\frac i2\int\omega\wedge\bar\omega>0$ for holomorphic forms; $\int\theta\wedge\eta=\sum_j(A_j(\theta)B_j(\eta)-B_j(\theta)A_j(\eta))$ \emph{(proof sketched
here)}; Riemann's relations; the periods form a lattice $\Lambda_X$ \emph{(proved here, given the sketched identity)}. \Cref{thm:curves:bilinear-relations}.
\item \textbf{First homology} (definition): $H_1(X,\Z)$, $\varphi_*$, $H_1(X,A)=H_1\otimes A$. \Cref{def:curves:homology}.
\item \textbf{Symplectic bases.} $H_1(X,\Z)\simeq\Z^{2g}$ with a symplectic basis; the intersection pairing is perfect and alternating, $a_i\cdot b_j=\delta_{ij}$. \emph{(stated; reference)}
\Cref{thm:curves:symplectic-basis}.
\item \textbf{First cohomology} (definition): $H^1(X,A)=\Hom(H_1(X,\Z),A)$; $\varphi^*$ the transpose of $\varphi_*$; de Rham classes; $H^2$. \Cref{def:curves:cohomology}.
\item \textbf{Cohomology.} $H^1(X,\Zl)=H^1(X,\Z)\otimes\Zl$ free of rank $2g$ \emph{(proved here)}; de Rham and $H^2(X,\Z)=\Z$ \emph{(stated; reference)}; Hodge
$H^1(X,\C)=H^0(\Omega^1)\oplus\overline{H^0(\Omega^1)}$ \emph{(proof sketched here)}; $H^1(X,M)=H^1(\pi_1(X),M)$ \emph{(stated; reference)}. \Cref{prop:curves:cohomology}.
\item \textbf{Cup product.} $\langle\theta,\eta\rangle=\int_X\theta\wedge\eta$ is integral, alternating and perfect on $H^1(X,\Z)$ with Gram matrix $\sm0{I_g}{-I_g}0$, and after any base change, in
particular on $H^1(X,\Zl)$; Poincare duality $\operatorname{PD}(a_j^\vee)=b_j$, $\operatorname{PD}(b_j^\vee)=-a_j$ is an isometry. \emph{(proved here, given \Cref{thm:curves:symplectic-basis} and
\Cref{thm:curves:bilinear-relations}(c))} \Cref{thm:curves:cup-product}.
\item \textbf{Transfer and projection formula.} $\varphi^!\colon H_1(Y)\to H_1(X)$ with $\varphi_*\varphi^!=n$ \emph{(proof sketched here)}; $\varphi_*$ on $H^1$ (transpose of $\varphi^!$, the trace of
forms); $\langle\varphi_*\eta,\xi\rangle=\langle\eta,\varphi^*\xi\rangle$, $\varphi_*\varphi^*=n$, $\langle\varphi^*x,\varphi^*y\rangle=n\langle x,y\rangle$, $\varphi^*=n$ on $H^2$ \emph{(proved here)}. \Cref{prop:curves:transfer}.
\item \textbf{Which is dual to which.} $H_1$, $\Lambda_X$, $J$, $J[n]$, $T_\ell J$ covariant; $H^1=\Hom(T_\ell J,\Zl)$, $H^0(\Omega^1)$ contravariant; $\operatorname{PD}$ turns $\mathcal C^*$ into $(\mathcal C^t)_*$;
$H^1_{\mathrm{et}}\simeq T_\ell J(-1)$. \Cref{conv:curves:duality}.
\item \textbf{Jacobian} (definition): $J(C)(\bar k)=\Pic^0(C)$; $J_{\mathrm{an}}(X)=H^0(X,\Omega^1)^\vee/\Lambda_X$. \Cref{def:curves:jacobian}.
\item \textbf{Abel-Jacobi.} $\operatorname{AJ}\colon\Div^0(X)\to J_{\mathrm{an}}(X)$ is a well-defined homomorphism and $\iota\colon X\to J_{\mathrm{an}}$ is a holomorphic injection inducing the period isomorphism on
$H_1$ \emph{(proved here, injectivity given Abel's theorem)}; Abel's theorem and Jacobi inversion, $\Pic^0(X)\simeq J_{\mathrm{an}}(X)$ \emph{(stated; reference)}. \Cref{thm:curves:abel-jacobi}.
\item \textbf{Jacobian variety.} $J$ is an abelian variety of dimension $g$ over $k$ with $J(\bar k)=\Pic^0(C)$, $J(\C)=J_{\mathrm{an}}$, the Albanese universal property, and Mordell-Weil.
\emph{(stated; reference)} \Cref{thm:curves:jacobian-variety}.
\item \textbf{Functoriality.} $\varphi_*\colon J(X)\to J(Y)$ and $\varphi^*\colon J(Y)\to J(X)$ with $\varphi_*\varphi^*=[n]$; over $\C$, $\varphi_*$ is induced by $(\varphi^*)^\vee$ and $\varphi_*$ on $H_1$
\emph{(proved here)}, $\varphi^*$ by the trace and $\varphi^!$ \emph{(proof sketched here)}; for $Y=E$, $\varphi=\varphi_*\circ\iota_X$. \Cref{prop:curves:jacobian-functoriality}.
\item \textbf{Tori.} For a surjective map of complex tori, $\pi_0(\ker)\simeq M/A(\Lambda)$; for $B\colon W\to V$ $\C$-linear and injective with $B(M)\subseteq\Lambda$, $\ker\simeq\operatorname{sat}(B(M))/B(M)$; maps of one-dimensional tori are $z\mapsto cz$ of
degree $[\Lambda':c\Lambda]$. \emph{(proved here)} \Cref{lem:curves:tori}.
\item \textbf{Degree and cup product.} $\langle\varphi^*x,\varphi^*y\rangle=n\langle x,y\rangle$; $\varphi_*\varphi^*=n$; $\varphi^*H^1(Y,\Z)$ saturated iff $\varphi_*$ onto on $H_1$ iff $\ker(J(X)\to J(Y))$
connected iff $\varphi^*$ injective on $J$, the indices all equal; for $Y=E$, the cup product on $\varphi^*H^1(E)$ has ideal of values $n\Z$, Gram determinant $n^2$, $[H^1:M\oplus M^\perp]=n^2$ if
saturated, and ideal $(n/m_0)\Z$ on the saturation in general, $m_0=\#\pi_0(\ker(J(X)\to E))$; an isogeny of degree $m$ destroys saturation. \emph{(proved here; the identification of $\varphi^*$ on $J$ sketched)}
\Cref{prop:curves:degree-cup-product}.
\item \textbf{Torsion.} $J(X)[n]=\frac1n\Lambda_X/\Lambda_X\simeq H_1(X,\Z/n)\simeq(\Z/n)^{2g}$; $J[\mf m]$ is a $\T/\mf m$-vector space \emph{(proved here)}; $J[n](\bar k)\simeq(\Z/n)^{2g}$ over $k$,
$\operatorname{char}k\nmid n$ \emph{(stated; reference)}. \Cref{prop:curves:torsion}.
\item \textbf{Tate modules} (definition): $T_\ell J=\varprojlim J[\ell^m]$, $V_\ell J=T_\ell J\otimes\Ql$. \Cref{def:curves:tate-module}.
\item \textbf{Tate module and $H^1$.} $T_\ell J\simeq H_1(X,\Zl)$, free of rank $2g$; $H^1(X,\Zl)=\Hom_{\Zl}(T_\ell J,\Zl)$; $\operatorname{PD}$ gives $H^1\simeq T_\ell J$, not Galois-equivariant, turning $\mathcal C^*$
into $(\mathcal C^t)_*$ \emph{(proved here)}; over $k$, $T_\ell J$ free of rank $2g$ with $G_k$-action \emph{(stated; reference)}. \Cref{prop:curves:tate-module}.
\item \textbf{Weil pairing and etale cohomology.} $e_\ell\colon T_\ell J\times T_\ell J\to\Zl(1)$ perfect alternating equivariant; $H^1_{\mathrm{et}}=(T_\ell J)^\vee\simeq T_\ell J(-1)$, comparison with
$H^1(C(\C),\Z)\otimes\Zl$; for $E$, $H^1(E,\Zl)\simeq T_\ell E(-1)$ with determinant $\chi_\ell^{-1}$; Weil's theorem over $\F_q$ and $a_q=\tr(\Frob_q\mid T_\ell E)=\tr(\Frob_q^{-1}\mid H^1)$. \emph{(stated;
reference)} \Cref{thm:curves:weil-pairing-etale}.
\item \textbf{Correspondences} (definition): $\mathcal C=(C,\alpha,\beta)$, transpose $(C,\beta,\alpha)$, graphs, $u\mathcal Cv$. \Cref{def:curves:correspondence}.
\item \textbf{Actions} (definition): $\mathcal C_*=\beta_*\alpha^*$ on divisors and $J$, $\mathcal C_*=\beta_*\alpha^!$ on $H_1$, $\mathcal C^*=\alpha_*\beta^*$ on $H^1$ and on differentials.
\Cref{def:curves:correspondence-actions}.
\item \textbf{Compatibilities.} $\mathcal C^*$ on $H^1$ is the transpose of $\mathcal C_*$ on $H_1$; graphs act by $\varphi^*$, $\varphi_*$, $\varphi^!$; $u\mathcal Cv$ acts by $u_*\mathcal C_*v_*$ and
$v^*\mathcal C^*u^*$ \emph{(proved here)}; the action on $J$ and $T_\ell J$ is $\mathcal C_*$ \emph{(proof sketched here)}. \Cref{prop:curves:correspondence-actions}.
\item \textbf{The adjoint is the transpose.} $\langle\mathcal C^*\xi,\eta\rangle_X=\langle\xi,(\mathcal C^t)^*\eta\rangle_Y$, that is $(\mathcal C^*)^\dagger=(\mathcal C^t)^*$; automorphisms are isometries;
$\operatorname{PD}\circ\mathcal C^*=(\mathcal C^t)_*\circ\operatorname{PD}$. \emph{(proved here, given \Cref{prop:curves:transfer})} \Cref{thm:curves:adjoint-transpose}.
\item \textbf{Twisted pairing.} If $t^t=w^*tw^*$ for every $t$ in a commutative $\Zl$-subalgebra $\T$ of $\End_{\Zl}(H^1(X,\Zl))$ spanned by actions of correspondences, $w$ an involution of $X$, then $\langle x,w^*y\rangle$ is perfect and alternating
on $H^1(X,\Zl)$, every $t\in\T$ is self-adjoint for it, and $H^1(X,\Zl)=\bigoplus_{\mf m}H^1(X,\Zl)_{\mf m}$ is orthogonal with the pairing perfect on each summand. \emph{(proved here)}
\Cref{cor:curves:twisted-pairing}.
\end{enumerate}

\section*{Where the paper uses this chapter}

\begin{itemize}
\item \textbf{Abstract and Section 1} (the modular parametrisation $\varphi_E\colon X_0(N)\to E$, the $X_0(N)$-optimal curve ``for which the modular parametrisation has minimal degree'', the modular
degree $\degphi$): morphisms and their degree in every sense, \Cref{def:curves:degree}, \Cref{thm:curves:degree-formula}, \Cref{cor:curves:degree-three-ways}, \Cref{tab:curves:degree},
\Cref{thm:curves:riemann-existence}; Riemann-Hurwitz checks of the quoted degrees, \Cref{cor:curves:rh-consequences}(e); why the optimal curve has minimal degree,
\Cref{prop:curves:degree-cup-product}(iv), (v) and \Cref{pl:curves:pullback}.
\item \textbf{Remark rem:degree (Section 4, sec:quant)} (the formula of Zagier and Watkins: Petersson norm over the area of the period lattice): positivity and periods, \Cref{thm:curves:bilinear-relations},
\Cref{ex:curves:torus}.
\item \textbf{Section 1.1 and Proposition prop:ns} ($\rho_{E,\ell}$ and the action of $G_\Q$ on $E[3]$ and on the roots of the $3$-division polynomial): the Galois action on points,
\Cref{prop:curves:galois-action}, \Cref{pl:curves:galois}.
\item \textbf{Section 1 and the Appendix} (curves given by Weierstrass equations, \texttt{ellinit}): \Cref{thm:curves:genus-one-weierstrass}, \Cref{prop:curves:genus-one-pic},
\Cref{pl:curves:weierstrass}.
\item \textbf{Section 4.1 (sec:dfg)} (``perfect and alternating on $H^1(X_0(M^\Sigma),\Zl)$, and the Hecke operators, $U_p$ included, are self-adjoint for it, because $w_{M^\Sigma}T_nw_{M^\Sigma}^{-1}$ is the adjoint of
$T_n$ for the cup product''): \Cref{thm:curves:cup-product}, \Cref{def:curves:correspondence-actions}, \Cref{thm:curves:adjoint-transpose}, \Cref{cor:curves:twisted-pairing},
\Cref{pl:curves:cup}, \Cref{pl:curves:hecke}.
\item \textbf{Proposition prop:pullback} ($\varphi^*$ injective with image $H^1(X_0(N),\Zl)[I_f]$; $\langle\varphi^*a,w_N^*\varphi^*b\rangle=\varepsilon_N(f)\degphi\langle a,b\rangle_E$; the congruence ideal is $\degphi\,\Zl$):
optimality as connectedness, saturation and surjectivity on $H_1$, \Cref{lem:curves:tori}, \Cref{prop:curves:degree-cup-product}(iii); the projection formula, \Cref{prop:curves:transfer}(iv);
Hodge, \Cref{prop:curves:cohomology}(iii); unimodularity, \Cref{thm:curves:cup-product}; step-by-step reading in \Cref{pl:curves:pullback}.
\item \textbf{The paragraph after Proposition prop:pullback and Section 4.2} ($M_{f,\lambda}=H^1(E,\Zl)$, which is $\Hom(T_\ell E,\Zl)\simeq T_\ell E(-1)$; ``its adjoint is $\ad T_\ell E$''):
\Cref{prop:curves:tate-module}(ii);
\Cref{thm:curves:weil-pairing-etale}(b), (c); \Cref{warn:curves:not-galois}; \Cref{pl:curves:free}.
\item \textbf{Lemma lem:hecke, Lemma lem:ddt417 as applied, and Step 2 of the proof of Theorem thm:quant} ($L=H^1(X_0(N),\Zl)_{\mf m}$ ``with a perfect $\mathcal O$-bilinear pairing for which $\T$ is
self-adjoint''; ``the idempotent $e_{\mf m}$ is self-adjoint for $\hat\delta_!$, so $H^1(X_0(N),\Zl)$ splits orthogonally into its local components and the pairing is perfect on each''):
\Cref{cor:curves:twisted-pairing}(iii).
\item \textbf{Remark rem:free} ($J_0(N)[\mf m]$ two-dimensional, the Tate module free of rank two, $T_\ell J_0(N)$ versus $H^1(X_0(N),\Zl)$ through Poincare duality composed with $w_N$, not Galois
equivariant): \Cref{prop:curves:torsion}, \Cref{prop:curves:tate-module}, \Cref{thm:curves:adjoint-transpose}(c), \Cref{warn:curves:not-galois}, \Cref{pl:curves:free}.
\end{itemize}
\section{Affine and projective varieties}\label{sec:curves:varieties}

This section fixes the language. A reader who knows complex analysis and linear algebra meets here algebraic geometry and
algebraic topology at once, so the chapter keeps two languages and one dictionary. The \emph{algebraic} side is the classical
language of Silverman \cite[Chapters I and II]{SilvermanAEC}: a variety is the set of its points over an algebraic closure
$\bar k$, cut out by polynomials, with the Galois group $G_k$ acting on coordinates; no schemes, no sheaves and no sheaf
cohomology are used. The \emph{analytic} side (\Cref{sec:curves:riemann-surfaces,sec:curves:differentials-analytic,sec:curves:cohomology})
is the theory of compact Riemann surfaces with charts, as in a first course in complex analysis. The dictionary between
the two (Riemann's existence theorem and GAGA) is stated once, in \Cref{thm:curves:riemann-existence}, and every other
comparison is proved from it. What is proved: every statement that is a computation with polynomials, valuations, local
normal forms of holomorphic maps, change of variables in integrals, or linear algebra of lattices. What is sketched: the
formula $\sum e_P=\deg\varphi$, the genus of plane curves, the degree of the Frobenius morphism, the topological form of Riemann-Hurwitz, Riemann's bilinear
relations, the Hodge decomposition, the transfer, the description of $\varphi^*$ on the Jacobian by the trace and the transfer, and the action of correspondences
on the Jacobian.
What is stated with a reference: the Nullstellensatz, Bezout's theorem, the equivalence between curves and function fields,
Riemann-Roch, Riemann existence and GAGA, the classification of surfaces, Abel's theorem and Jacobi inversion, the
algebraicity of the Jacobian, the Weil pairing and etale cohomology. In the paper the objects of this section are the
background of every sentence that mentions $X_0(N)$, $E$ or $E(\bar\Q)$.

\begin{convention}[conventions of this chapter]\label{conv:curves:conventions}
\leavevmode
\begin{enumerate}[label=(C\arabic*)]
\item \emph{Fields.} $k$ is a perfect field ($\Q$, $\Q_p$, a number field, a finite field $\F_q$ with $q$ a power of a prime
$p$: here $q$ is always the size of a finite field, \Cref{ch:notation}, B-C1); $\bar k$ is an algebraic closure and
$G_k=\Gal(\bar k/k)$ (\Cref{ch:galois}). Over $\C$ we write $X$ for a compact Riemann surface and identify it with the smooth
projective curve it defines (\Cref{thm:curves:riemann-existence}).
\item \emph{Curves.} A \emph{curve over $k$} is a projective variety of dimension one defined over $k$; ``variety'' includes
``irreducible over $\bar k$''. Unless said otherwise curves are smooth and projective; a singular or affine curve
birational to $C$ is called a \emph{model} of $C$. Morphisms are defined over $k$ when $k$ is named, over $\bar k$ otherwise.
The genus is $g$.
\item \emph{Points and Galois.} $C(L)$ is the set of points with coordinates in $L$, for $k\subseteq L\subseteq\bar k$; $\sigma\in G_k$ acts on
$C(\bar k)$ coordinatewise, on the left, and $C(L)=C(\bar k)^{G_L}$. This is the action behind $\rho_{E,\ell}$ (\Cref{ch:ellcurves,ch:galrep}).
\item \emph{Pullback, pushforward, transpose, adjoint.} For a nonconstant morphism $\varphi\colon X\to Y$, $\varphi^*$ is pullback (of functions,
divisors, divisor classes, differentials, classes in $H^1$) and $\varphi_*$ is pushforward (of divisors, divisor classes, classes in
$H_1$, and on $H^1$ the transpose of the transfer, \Cref{prop:curves:transfer}). A correspondence is $\mathcal C=(C,\alpha,\beta)$ and its
transpose $\mathcal C^t=(C,\beta,\alpha)$. The adjoint of an endomorphism $A$ for a pairing is $A^\dagger$, never $A^*$, because a star after
the name of a map is a pullback (\Cref{ch:notation}, B-C4).
\item \emph{Duals.} $M^\vee=\Hom(M,\Z)$ or $\Hom_{\Zl}(M,\Zl)$ for lattices; $V^\vee$ for linear duals and contragredients (\Cref{ch:reps}).
Pontryagin duals do not occur in this chapter.
\item \emph{Orientation and the cup product.} A Riemann surface is oriented by its complex structure: in a chart $z=x+iy$ the form
$dx\wedge dy=\frac i2dz\wedge d\bar z$ is positive. The cup product of classes in $H^1$ is $\langle\theta,\eta\rangle=\int_X\theta\wedge\eta$ (de Rham
representatives). The intersection number $a\cdot b$ of transverse oriented loops counts crossings with sign $+1$ when (tangent of $a$,
tangent of $b$) is a positive basis. Poincare duality $\operatorname{PD}\colon H^1\to H_1$ is fixed by $\int_{\operatorname{PD}(\theta)}\eta=\langle\theta,\eta\rangle$.
Other books use other signs; the paper needs only properties that do not change under a global sign (perfectness, alternation,
adjointness, ideals of values, Gram determinants in even rank), \Cref{ch:notation}, B-C5(d).
\item \emph{Frobenius.} The arithmetic Frobenius $\Frob_q$ of $G_{\F_q}$ raises coordinates of $\bar\F_q$-points to the $q$-th power; the
$q$-power Frobenius \emph{morphism} of a curve over $\F_q$ is $\pi_q$; the geometric Frobenius is $\Frob_q^{-1}$ (\Cref{ch:notation}, B-C2).
\item \emph{Notation.} We write $l(D)$, not $\ell(D)$, for the dimension of a Riemann-Roch space, because $\ell$ is a prime; $K_C$ is a
canonical divisor and $K$ a field; $D$ is always a divisor. ``Riemann-Roch'', ``Riemann-Hurwitz'', ``Abel-Jacobi'' take plain
hyphens.
\end{enumerate}
\end{convention}

\begin{definition}[affine varieties]\label{def:curves:affine-variety}
Let $\A^n=\A^n(\bar k)=\bar k^n$, with $G_k$ acting coordinatewise, so that $\A^n(L)=\A^n(\bar k)^{G_L}$. For a set $S\subseteq\bar k[x_1,\dots,x_n]$ its zero set
$V(S)\subseteq\A^n$ is an \emph{algebraic set}; for $Y\subseteq\A^n$, $I(Y)$ is the ideal of polynomials vanishing on $Y$. The algebraic sets are the closed sets
of a topology, the \emph{Zariski topology} (finite unions and arbitrary intersections of algebraic sets are algebraic: $V(S)\cup V(S')=V(SS')$,
$\bigcap V(S_i)=V(\bigcup S_i)$). An \emph{affine variety} is an algebraic set $V$ with $I(V)$ prime. It is \emph{defined over $k$} if $I(V)$ is generated by
polynomials in $k[x]$. Its \emph{coordinate ring} is the domain $\bar k[V]=\bar k[x]/I(V)$, its \emph{function field} $\bar k(V)=\Frac\bar k[V]$; if $V$ is defined
over $k$, $k[V]=k[x]/(I(V)\cap k[x])$ and $k(V)=\Frac k[V]$. A function $f\in\bar k(V)$ is \emph{regular at} $P\in V$ if $f=g/h$ with $g,h\in\bar k[V]$ and
$h(P)\ne0$; then $f(P)=g(P)/h(P)$ does not depend on the choice. The \emph{dimension} of $V$ is the transcendence degree of $\bar k(V)$ over $\bar k$
(\cite[Chapter VIII, Section 1]{Lang}; \Cref{ch:galois} treats algebraic and transcendental elements).
\end{definition}

\begin{theorem}[Hilbert's Nullstellensatz; stated]\label{thm:curves:nullstellensatz}
For every ideal $\mf a$ of $\bar k[x_1,\dots,x_n]$, $I(V(\mf a))=\sqrt{\mf a}$. Hence $V\mapsto I(V)$ is a bijection from algebraic sets to radical ideals, under
which affine varieties correspond to prime ideals and points to maximal ideals. For $k$ perfect, an affine variety $V$ is defined over $k$ iff
$\sigma(V)=V$ for all $\sigma\in G_k$.
\end{theorem}

We do not prove this; see \cite[Chapter 1]{Fulton} and \cite[Chapter I, Section 1]{Hartshorne} for the Nullstellensatz, and
\cite[Chapter I, Section 1]{SilvermanAEC} with its exercises for the Galois criterion. Examples: $\A^1$, with $\bar k[\A^1]=\bar k[x]$; the
parabola $y=x^2$, with $\bar k[V]=\bar k[x,y]/(y-x^2)\simeq\bar k[x]$ and $\dim V=1$; the affine Weierstrass cubics of \Cref{ex:curves:plane-examples}.

Why points over $\bar k$ and not over $k$? Over a field that is not algebraically closed a polynomial can have too few zeros to be
recovered from them ($x^2+y^2+1$ has no zero in $\R^2$, and neither has $1$), so the dictionary between ideals and zero sets of the Nullstellensatz fails;
over $\bar k$ it holds, and the arithmetic over $k$ is then recorded by the action of $G_k$. Two small examples over $k=\Q$ show how the
definitions interact. The set $V(x^2-2)\subseteq\A^1$ is $\{\sqrt2,-\sqrt2\}$; its ideal $(x^2-2)\bar k[x]=(x-\sqrt2)(x+\sqrt2)\bar k[x]$ is not prime, so it is an
algebraic set but not a variety, although $x^2-2$ is irreducible over $\Q$. Each of its two points is a variety, $V(x\mp\sqrt2)$, but neither is defined over
$\Q$, and indeed $\sigma(\sqrt2)=-\sqrt2$ for some $\sigma\in G_\Q$, so the Galois criterion of the theorem fails for each point while the union is $G_\Q$-stable. In the
same way $V(x^2+1)\subseteq\A^2$ is the union of the two lines $x=\pm i$, exchanged by complex conjugation: the polynomial $x^2+1$ is irreducible over $\Q$ but
not over $\bar\Q$. This is why ``variety'' means ``irreducible over $\bar k$'' in this chapter (\Cref{conv:curves:conventions}(C2)).

\begin{definition}[projective space and projective varieties]\label{def:curves:projective-space}
$\PP^n=\PP^n(\bar k)=(\bar k^{n+1}\setminus\{0\})/\bar k^\times$, with homogeneous coordinates $[x_0:\dots:x_n]$; for $n=1$ this is the set of lines of
\Cref{ch:groups}, on which $\PGL_2$ acts by Moebius transformations. $G_k$ acts coordinatewise, and $\PP^n(L)$ is the set of points with a
representative in $L^{n+1}$. An ideal of $\bar k[x_0,\dots,x_n]$ is \emph{homogeneous} if it is generated by homogeneous polynomials; for a set $S$ of
homogeneous polynomials $V(S)\subseteq\PP^n$ is well defined (whether $F(x)=0$ does not depend on the representative $x$) and is a \emph{projective
algebraic set}; a \emph{projective variety} is a projective algebraic set whose homogeneous ideal $I(V)$ is prime, \emph{defined over $k$} if
$I(V)$ is generated by polynomials with coefficients in $k$. The sets $U_i=\{x_i\ne0\}$ are the standard charts, $U_i\simeq\A^n$ by dividing by
$x_i$. The \emph{projective closure} $\bar V$ of an affine variety $V\subseteq\A^n=U_0$ is the zero set of the homogenisations of the elements of
$I(V)$; conversely $V\cap U_i$ is obtained by dehomogenising. The \emph{function field} $\bar k(V)$ of a projective variety is the field of
quotients $F/G$ of homogeneous polynomials of the same degree with $G\notin I(V)$, two quotients being equal when $FG'-F'G\in I(V)$; such an $f$ is
\emph{regular at} $P\in V$ if $f=F/G$ with $G(P)\ne0$, and then $f(P)=F(P)/G(P)$.
\end{definition}

\begin{proposition}[points over subfields]\label{prop:curves:projective-points}
Let $k\subseteq L\subseteq\bar k$.
\begin{enumerate}[label=(\roman*)]
\item $\PP^n(\bar k)^{G_L}=\PP^n(L)$.
\item $\PP^n=\A^n\sqcup\PP^{n-1}$ (the chart $U_0$ and the hyperplane $x_0=0$); in particular $\PP^1=\A^1\sqcup\{\infty\}$.
\item For a projective variety $V$ meeting $U_i$, the map $F/G\mapsto(F/G)(x_0,\dots,1,\dots,x_n)$ (the $1$ in position $i$) is an isomorphism of $\bar k(V)$
onto the function field of the affine variety $V\cap U_i$; for $P\in V\cap U_i$, a function is regular at $P$ on $V$ iff its image is regular at $P$ on
$V\cap U_i$, with the same value.
\end{enumerate}
\end{proposition}

\begin{proof}
(i) A point with a representative in $L^{n+1}$ is fixed by $G_L$. Conversely let $P=[x_0:\dots:x_n]$ be fixed by $G_L$; choose $i$ with $x_i\ne0$ and
scale so that $x_i=1$. For $\sigma\in G_L$, $\sigma P=P$ means $\sigma(x_j)=cx_j$ for all $j$ with one $c\in\bar k^\times$; the coordinate $j=i$ gives $c=1$. So
every $x_j$ is fixed by $G_L$, hence lies in $L$ ($\bar k^{G_L}=L$ by Galois theory, \Cref{ch:galois}, $k$ being perfect). (ii) A point $[x_0:\dots:x_n]$
either has $x_0\ne0$, and then it has exactly one representative with $x_0=1$, namely $(1,x_1/x_0,\dots,x_n/x_0)$, which identifies $U_0$ with $\A^n$; or it has
$x_0=0$, and then $[x_1:\dots:x_n]$ is a well-defined point of $\PP^{n-1}$ (not all $x_j$ vanish, and scaling acts on both sides in the same way). For $n=1$
the second set is the single point $[0:1]$, which is called $\infty$. (iii) A quotient
$F/G$ of forms of degree $d$ equals $(F/x_i^d)/(G/x_i^d)$, a quotient of polynomials in the affine coordinates $x_j/x_i$; conversely a quotient of
polynomials in the affine coordinates is homogenised by multiplying numerator and denominator by a power of $x_i$. The two constructions are
inverse, and both respect the ideals. If $G(P)\ne0$ then $(G/x_i^d)(P)\ne0$, as $x_i(P)\ne0$; conversely, if $f=g/h$ with $h(P)\ne0$ on $V\cap U_i$,
homogenising $g$ and $h$ to forms $F$, $G$ of the same degree $d$ (multiplying the one of lower degree by a power of $x_i$) gives $f=F/G$ with
$G(P)=x_i(P)^dh(P)\ne0$. The values agree, since both are computed from the same representative of $P$ with $x_i=1$.
\end{proof}

\begin{definition}[curves, rational maps, morphisms]\label{def:curves:curve-morphism}
A \emph{curve over $k$} is a projective variety of dimension one defined over $k$ (\Cref{conv:curves:conventions}). Let $V_1$, $V_2\subseteq\PP^n$
be projective varieties. A \emph{rational map} $\varphi=[f_0:\dots:f_n]\colon V_1\dashrightarrow V_2$ is given by $f_i\in\bar k(V_1)$, not all zero, such that
$\varphi(P)=[f_0(P):\dots:f_n(P)]\in V_2$ wherever all $f_i$ are defined and not all zero; $[f_0:\dots:f_n]$ and $[gf_0:\dots:gf_n]$ define the same map
for $g\in\bar k(V_1)^\times$. It is \emph{regular} at $P$ if there is $g$ with all $gf_i$ regular at $P$ and some $(gf_i)(P)\ne0$; a \emph{morphism} is a
rational map regular at every point; an \emph{isomorphism} is a morphism with an inverse morphism. $\varphi$ is \emph{defined over $k$} if it can be
written with all $f_i\in k(V_1)$; for $k$ perfect this holds iff $\varphi(\sigma P)=\sigma(\varphi(P))$ for all $\sigma\in G_k$ and all $P$
(\cite[Chapter I, Section 3]{SilvermanAEC}; stated).
\end{definition}

\begin{example}[a conic is a line]\label{ex:curves:conic}
The map $\psi=[s^2:st:t^2]\colon\PP^1\to\PP^2$ is a morphism onto the conic $C\colon XZ=Y^2$, defined over every $k$. Its inverse is
$[X:Y]=[Y:Z]$, a rational map on $C$ that is regular everywhere: at points with $X\ne0$ use $[X:Y]$, at points with $Z\ne0$ use $[Y:Z]$; the two
agree on $C$ because $XZ=Y^2$, and $X=Z=0$ forces $Y=0$, which is not a point. So $C\simeq\PP^1$ over $k$.
\end{example}

\section{Plane curves and smoothness}\label{sec:curves:plane}

The curves of the paper are given by equations: the elliptic curves by Weierstrass equations (the Appendix of the paper calls
\texttt{ellinit([1,1,1,-30,-76])}), and $X_0(N)$ by the analytic construction of \Cref{ch:modforms}, whose smooth projective
model is the curve of \Cref{thm:curves:riemann-existence}. This section says when an equation defines a smooth curve and proves the
local fact behind every order of a zero or a pole in the rest of the book: at a smooth point the local ring is a discrete
valuation ring (\Cref{thm:curves:smooth-dvr}). Bezout's theorem is stated for the degree of the divisor of a linear form.

\begin{definition}[plane curves; singular points]\label{def:curves:plane-curve}
A \emph{plane curve of degree $d$} is $C=V(F)\subseteq\PP^2$ for a homogeneous $F\in k[X,Y,Z]$ of degree $d$ that is irreducible in $\bar k[X,Y,Z]$; its
affine part is $f(x,y)=F(x,y,1)=0$. A point $P\in C(\bar k)$ is \emph{singular} if $F_X(P)=F_Y(P)=F_Z(P)=0$ (subscripts are partial derivatives) and
\emph{smooth} otherwise; $C$ is \emph{smooth} if it has no singular point in $C(\bar k)$. For a general variety $V\subseteq\PP^n$ with $I(V)=(F_1,\dots,F_m)$
of dimension one, $P$ is smooth if the Jacobian matrix $(\partial F_i/\partial x_j(P))$ has rank $n-1$ \cite[Chapter I, Section 1]{SilvermanAEC}.
\end{definition}

\begin{proposition}[the affine criterion; tangent lines]\label{prop:curves:euler}
Let $C=V(F)$ be a plane curve of degree $d$.
\begin{enumerate}[label=(\roman*)]
\item (Euler) $XF_X+YF_Y+ZF_Z=dF$. Hence a point $P=(a,b)$ of the affine part is singular iff $f(P)=f_x(P)=f_y(P)=0$; and if $\operatorname{char}k\nmid d$,
every common zero in $\PP^2$ of $F_X$, $F_Y$, $F_Z$ lies on $C$.
\item At a smooth point $P$ the \emph{tangent line} is $F_X(P)X+F_Y(P)Y+F_Z(P)Z=0$; on the affine part it is $f_x(P)(x-a)+f_y(P)(y-b)=0$.
\end{enumerate}
\end{proposition}

\begin{proof}
(i) Both sides of Euler's identity are additive in $F$, and for a monomial $X^iY^jZ^m$ with $i+j+m=d$ the left side is $(i+j+m)X^iY^jZ^m$. On the affine
part $f_x(a,b)=F_X(a,b,1)$ and $f_y(a,b)=F_Y(a,b,1)$, and at a point of $C$, Euler gives $F_Z(a,b,1)=-aF_X(a,b,1)-bF_Y(a,b,1)$; so the three
partials vanish iff $f_x$ and $f_y$ do. If $\operatorname{char}k\nmid d$ and the three partials vanish at $P$, Euler gives $dF(P)=0$, so $F(P)=0$. (ii) The
projective equation is a definition. On the affine part put $Z=1$, $X=x$, $Y=y$ and $P=(a,b,1)$; using $F_Z(P)=-af_x(P)-bf_y(P)$ from (i), the
projective equation becomes
\[
f_x(P)x+f_y(P)y-af_x(P)-bf_y(P)=f_x(P)(x-a)+f_y(P)(y-b)=0,
\]
the affine tangent line of calculus.
\end{proof}

So for an affine equation one never needs the homogenisation away from infinity: smoothness at $(a,b)$ is the familiar condition that the gradient
$(f_x,f_y)$ does not vanish there, and (in the real picture) the tangent line is orthogonal to that gradient. Only the points at infinity require the projective partials, and
\Cref{ex:curves:plane-examples} shows how they are handled.

\begin{example}[the paper's curves are smooth; the node and the cusp are not]\label{ex:curves:plane-examples}
\leavevmode
\begin{enumerate}[label=(\alph*)]
\item The \emph{nodal cubic} $y^2=x^3+x^2$ and the \emph{cuspidal cubic} $y^2=x^3$ have the single singular point $(0,0)$: for $f=y^2-x^3-x^2$,
$f_y=2y$ and $f_x=-3x^2-2x$ vanish together on the curve only at $(0,0)$ when $\operatorname{char}k\ne2$ (at $x=-1$, $f_x=-1$), and at infinity ($[0:1:0]$, where
$F_Z=Y^2-X^2\ne0$) the curve is smooth; likewise for $y^2=x^3$. Over $\R$ the node is a loop crossing itself at the origin with two
tangent lines $y=\pm x$, and the cusp is a curve with a sharp point at the origin with the single tangent $y=0$ counted twice.
\item The \emph{Fermat curve} $X^d+Y^d+Z^d=0$ is smooth iff $\operatorname{char}k\nmid d$. If $\operatorname{char}k\nmid d$, the partials $dX^{d-1}$, $dY^{d-1}$, $dZ^{d-1}$ vanish
together only at $X=Y=Z=0$. If $d=p=\operatorname{char}k$ the polynomial is $(X+Y+Z)^p$, which is not even irreducible; for $p\mid d$ in general all
partials vanish identically.
\item The Weierstrass cubics
\begin{gather*}
\curve{11a1}\colon y^2+y=x^3-x^2-10x-20,\qquad \curve{20a1}\colon y^2=x^3+x^2+4x+4,\\
\curve{121a1}\colon y^2+xy+y=x^3+x^2-30x-76
\end{gather*}
are smooth over $\bar\Q$, including the point $[0:1:0]$ at infinity. At infinity: for $F=Y^2Z+a_1XYZ+a_3YZ^2-X^3-a_2X^2Z-a_4XZ^2-a_6Z^3$ one has
$F_Z(0,1,0)=1$, so $[0:1:0]$ is always smooth. On the affine part of \curve{11a1}, $f_y=2y+1$ and $f_x=-3x^2+2x+10$; a singular point would have
$y=-\frac12$ and then $x$ would be a double root of $4x^3-4x^2-40x-79$, whose discriminant is $-2^4\cdot11^5\ne0$ (computed with PARI). The general
criterion (a Weierstrass cubic is smooth iff its discriminant $\Delta$ is nonzero; $\Delta=-11^5$ for \curve{11a1}, $-2^8\cdot5^2$ for \curve{20a1},
$-11^2$ for \curve{121a1}) is proved in \Cref{ch:ellcurves}.
\end{enumerate}
\end{example}

\begin{theorem}[smooth points and discrete valuation rings]\label{thm:curves:smooth-dvr}
Let $C$ be a curve, $P\in C(\bar k)$, $\bar k[C]_P=\{f\in\bar k(C):f\text{ regular at }P\}$ (\Cref{def:curves:projective-space}) its local ring and $\mf m_P$ the maximal ideal of
functions vanishing at $P$. The following are equivalent: (a) $P$ is smooth; (b) $\dim_{\bar k}\mf m_P/\mf m_P^2=1$; (c) $\bar k[C]_P$ is a discrete
valuation ring. At a smooth point $P=(a,b)$ of an affine plane model $f=0$, the class of a linear form $\lambda=u(x-a)+v(y-b)$ is a uniformiser iff the
line $\lambda=0$ is not the tangent line; for instance $x-a$ is a uniformiser if $f_y(P)\ne0$. The valuation $\ord_P\colon\bar k(C)^\times\to\Z$ of $\bar k[C]_P$
is the \emph{order of vanishing} at $P$, and $g\in\bar k(C)$ is regular at $P$ iff $g=0$ or $\ord_Pg\ge0$.
\end{theorem}

\begin{proof}
We prove the theorem for a point of an affine plane model $f(x,y)=0$, $f$ irreducible; for a point of a curve in $\PP^n$ the same argument with the
Jacobian matrix gives the general case \cite[Chapter II, Proposition 1.1]{SilvermanAEC}. Translate so that $P=(0,0)$. Let $A=\bar k[x,y]$,
$\mf M=(x,y)$, $R=\bar k[C]=A/(f)$ and $\mf m=\mf M/(f)$. The local ring $\bar k[C]_P$ is the localisation $R_{\mf m}$ (\Cref{ch:modules}; regularity may be tested on an affine chart by
\Cref{prop:curves:projective-points}(iii)), with maximal ideal
$\mf m_P=\mf mR_{\mf m}$. Since $\mf m/\mf m^2$ is killed by $\mf m$, it is a vector space over $R/\mf m=\bar k$, and an element $s\in R\setminus\mf m$ acts on it as
the nonzero scalar $s(P)$ (because $s-s(P)\in\mf m$), hence invertibly. So inverting the elements of $R\setminus\mf m$ changes nothing:
$\mf m_P/\mf m_P^2=\mf m/\mf m^2=\mf M/(\mf M^2+(f))$.
Now $\mf M/\mf M^2$ is the two-dimensional space with basis the classes of $x$, $y$, and by Taylor's formula $f\equiv f_x(P)x+f_y(P)y\pmod{\mf M^2}$ (as
$f(P)=0$). So $\mf m_P/\mf m_P^2$ is the quotient of $\bar kx\oplus\bar ky$ by the span of $f_x(P)x+f_y(P)y$: of dimension $1$ if $P$ is smooth and $2$ if not
(\Cref{prop:curves:euler}(i)). This proves (a)$\Leftrightarrow$(b) and shows that a linear form is nonzero in $\mf m_P/\mf m_P^2$ iff it is not proportional to the
tangent form.

The ring $\bar k[C]_P$ is a Noetherian local domain (a localisation of a quotient of the Noetherian ring $A$), and not a field: $C$ has dimension one,
so $x$ and $y$ are not both constant on $C$, and a nonconstant one of them is a nonzero element of $\mf m_P$. (b)$\Rightarrow$(c): if $t\in\mf m_P$ spans
$\mf m_P/\mf m_P^2$, then $\mf m_P=t\bar k[C]_P+\mf m_P^2$, and Nakayama's lemma (\Cref{ch:modules}) gives $\mf m_P=(t)$; a Noetherian local domain that is not a
field and has principal maximal ideal is a discrete valuation ring with uniformiser $t$ (the criterion of \Cref{ch:modules}). (c)$\Rightarrow$(b): if
$\mf m_P=(t)$ then $\mf m_P/\mf m_P^2$ is spanned by the class of $t$, and it is nonzero because $\mf m_P\ne\mf m_P^2$ by Nakayama. The last sentence is the
description of a discrete valuation ring by its valuation (\Cref{ch:modules}).
\end{proof}

Two computations show the theorem at work. On the smooth cubic $y^2=h(x)$, $h$ squarefree of degree $3$, $\operatorname{char}k\ne2$, take a root $e$ of $h$ and
$P=(e,0)$. There $f_y=2y$ vanishes and $f_x=-h'(e)\ne0$, so the tangent line is the vertical line $x=e$, and $x-e$ is \emph{not} a uniformiser; $y$ is one.
Indeed $x-e=y^2/q(x)$ with $q(x)=h(x)/(x-e)$, and $q(e)=h'(e)\ne0$, so $q(x)$ is a unit at $P$ and $\ord_P(x-e)=2$: the vertical tangent meets the curve
doubly at $P$. At the node of $y^2=x^3+x^2$ the computation of the proof gives $\mf m_P/\mf m_P^2=\bar kx\oplus\bar ky$ (the gradient vanishes), of dimension two,
so the local ring at the node is not a discrete valuation ring and no single order of vanishing can be attached to a function there; the smooth model of
\Cref{thm:curves:smooth-model} replaces the node by two points, each with its own valuation.

\begin{theorem}[Bezout; stated]\label{thm:curves:bezout}
Two plane curves of degrees $m$ and $n$ without a common component meet in exactly $mn$ points of $\PP^2(\bar k)$ counted with intersection multiplicity.
In particular a line not contained in a plane curve of degree $d$ meets it in $d$ points counted with multiplicity, and a line meets a cubic in three
points (the chord-tangent construction of \Cref{ch:ellcurves}).
\end{theorem}

We do not prove this; see \cite[Chapter 5, Section 3]{Fulton} and \cite[Chapter I, Section 7]{Hartshorne}. On a smooth plane curve the
intersection multiplicity of the curve with a line $L=0$ at $P$ is $\ord_P(L/M)$ for any linear form $M$ with $M(P)\ne0$ \cite[Chapter 3]{Fulton}, so
the divisor of zeros of $L/M$ (\Cref{def:curves:divisor}) has degree $d$. For example, on a Weierstrass cubic the line at infinity $Z=0$ meets the curve
only at $O=[0:1:0]$ (put $Z=0$ in the equation: $X^3=0$), and in the chart $Y=1$ the function $Z/Y=v$ has $\ord_Ov=3$ (computed for $y^2=h(x)$ in the proof of
\Cref{prop:curves:differentials}(iv)): the line at infinity is a flex tangent, meeting the cubic at $O$ with multiplicity $3=d$, as Bezout requires.

\begin{theorem}[smooth projective models; stated]\label{thm:curves:smooth-model}
Every curve, possibly singular or affine, is birational over $k$ to a smooth projective curve, which is unique up to isomorphism over $k$. Its
points over $\bar k$ correspond bijectively to the discrete valuations of $\bar k(C)$ trivial on $\bar k$.
\end{theorem}

We do not prove this; it follows from \Cref{thm:curves:equivalence}, see \cite[Chapter I, Section 6]{Hartshorne} and \cite[Chapter 7]{Fulton}.
Example (proved): the nodal cubic $y^2=x^3+x^2$ is birational to $\PP^1$ through $t=y/x$, with inverse $x=t^2-1$, $y=t(t^2-1)$: substituting $y=tx$
in the equation gives $x^2(t^2-x-1)=0$. The two branches through the node correspond to $t=\pm1$: the smooth model separates them. In the paper
``the curve given by an equation'' always means the smooth projective model.

\section{Function fields, points over extensions, the Galois action}\label{sec:curves:function-fields}

This section turns a curve into its field of rational functions, and back. Its uses in the paper are two. The modular parametrisation
$\varphi_E\colon X_0(N)\to E$ is the same thing as the field embedding $\varphi_E^*\colon\Q(E)\to\Q(X_0(N))$ (\Cref{thm:curves:equivalence}), and its degree is
the degree of that field extension (\Cref{def:curves:degree}). And the Galois representation $\rho_{E,\ell}$ of Section 1.1 of the paper, and the action
of $G_\Q$ on $E[3]$ in Proposition prop:ns, are the action of $G_k$ on points of \Cref{prop:curves:galois-action}.

\begin{definition}[function fields of one variable]\label{def:curves:function-field}
A \emph{function field of one variable over $k$} is a finitely generated field extension $F/k$ of transcendence degree one in which $k$ is
algebraically closed. For a curve $C$ over $k$, $k(C)$ is one, and $\bar k(C)=\bar k\cdot k(C)$ is one over $\bar k$ \cite[Chapter II, Section 1]{SilvermanAEC}.
\end{definition}

The examples to keep in mind are $k(\PP^1)=k(x)$, and, for a Weierstrass cubic $E$, $k(E)=k(x)[y]/(y^2+a_1xy+a_3y-x^3-a_2x^2-a_4x-a_6)$, a quadratic
extension of $k(x)$. ``Transcendence degree one'' says that the curve has dimension one. The condition that $k$ be algebraically closed in $F$ excludes, for
$k=\Q$, a field such as $\Q(i)(y)$: it is finitely generated of transcendence degree one over $\Q$, and it is the function field over $\Q$ of the plane set
$x^2+1=0$, which over $\bar\Q$ falls apart into the two conjugate lines $x=\pm i$ (the second example after \Cref{thm:curves:nullstellensatz}). A function field
of one variable over $k$ is thus the function field of a curve that stays irreducible over $\bar k$.

\begin{proposition}[valuations of a function field]\label{prop:curves:function-field}
Let $C$ be a smooth curve.
\begin{enumerate}[label=(\roman*)]
\item Each $P\in C(\bar k)$ gives a discrete valuation $\ord_P$ of $\bar k(C)$ with $\ord_P(c)=0$ for $c\in\bar k^\times$.
\item For $f\in\bar k(C)^\times$, $\ord_Pf=0$ for all but finitely many $P$.
\item A rational function without poles is constant (proved in \Cref{cor:curves:degree-three-ways}(vi)).
\end{enumerate}
\end{proposition}

\begin{proof}
(i) is \Cref{thm:curves:smooth-dvr}; a nonzero constant is a unit of every local ring. (ii) Cover $C$ by finitely many affine pieces (the charts
$U_i$). On one piece write $f=g/h$ with $g$, $h$ regular functions, not zero. Off the zero sets $V(g)$ and $V(h)$ of the piece, $f$ is a unit of the
local ring, so $\ord_Pf=0$. The zero set of a nonzero regular function on an irreducible curve is a proper closed subset, hence finite (a proper
closed subset of an irreducible variety of dimension one has dimension zero; \cite[Chapter I, Section 1]{Hartshorne}, stated). (iii) is proved
after \Cref{thm:curves:degree-formula}.
\end{proof}

\begin{theorem}[rational maps from smooth curves are morphisms]\label{thm:curves:extend-morphism}
Let $C$ be a smooth curve, $V\subseteq\PP^n$ a projective variety and $\varphi=[f_0:\dots:f_n]\colon C\dashrightarrow V$ a rational map. Then $\varphi$ is a morphism.
In particular a rational map between smooth projective curves is a morphism, and every $f\in\bar k(C)$ defines a morphism $[f:1]\colon C\to\PP^1$
(constant if $f$ is), with $\infty=[1:0]$ the image of the poles.
\end{theorem}

\begin{proof}
Let $P\in C$, $t$ a uniformiser at $P$ (\Cref{thm:curves:smooth-dvr}) and $m=\min\{\ord_Pf_i:f_i\ne0\}$. Then each $t^{-m}f_i$ has $\ord_P\ge0$, so is
regular at $P$, and for an index with $\ord_Pf_i=m$ the function $t^{-m}f_i$ is a unit at $P$, nonzero at $P$. So $g=t^{-m}$ shows that $\varphi$ is regular
at $P$. For $f\in\bar k(C)$ nonconstant this applies to $[f:1]$; at a pole $m=\ord_Pf<0$ and $[t^{-m}f:t^{-m}](P)=[1:0]$.
\end{proof}

Smoothness of the source is essential. On the nodal cubic $y^2=x^3+x^2$ the rational map $t=y/x$ to $\PP^1$ is not regular at the node. The parametrisation
$\psi(t)=(t^2-1,t(t^2-1))$ of the example after \Cref{thm:curves:smooth-model} is a morphism $\PP^1\to C$ (extended to $\infty$ by this theorem, since $\PP^1$ is smooth), and
$(y/x)\circ\psi$ is the identity of $\PP^1$. If $y/x$ were regular at the node $S$ with some value $c$, then evaluating $(y/x)\circ\psi$ at $t=1$ and at $t=-1$,
both of which map to $S$, would give $c=1$ and $c=-1$. The two branches through the node ask for two different values, which a single point cannot
provide; the smooth model has two points there.

\begin{theorem}[curves and function fields; stated]\label{thm:curves:equivalence}
The functor $C\mapsto k(C)$, $(\varphi\colon C_1\to C_2)\mapsto(\varphi^*\colon k(C_2)\to k(C_1),\ f\mapsto f\circ\varphi)$ is a contravariant equivalence between the
category of smooth projective curves over $k$ with nonconstant morphisms defined over $k$, and the category of function fields of one variable over $k$
with $k$-embeddings. In particular two smooth projective curves with $k$-isomorphic function fields are isomorphic over $k$, and a nonconstant
morphism of degree one (\Cref{def:curves:degree}) is an isomorphism.
\end{theorem}

We do not prove this; see \cite[Chapter II, Theorem 2.4]{SilvermanAEC} and \cite[Chapter I, Corollary 6.12]{Hartshorne}. The idea: the points of
the smooth projective model are the discrete valuations of the function field trivial on the constants (\Cref{thm:curves:smooth-model}), and a field
embedding pulls back valuations.

\begin{proposition}[the Galois action on points and functions]\label{prop:curves:galois-action}
Let $C$ be a curve over $k$.
\begin{enumerate}[label=(\roman*)]
\item $G_k$ acts on $C(\bar k)$, and $C(\bar k)^{G_L}=C(L)$ for $k\subseteq L\subseteq\bar k$.
\item $G_k$ acts on $\bar k(C)$ through the coefficients: for $f=F/G$, $\sigma f=F^\sigma/G^\sigma$. For $f$ regular at $P$, $\sigma f$ is regular at $\sigma P$ and
$(\sigma f)(\sigma P)=\sigma(f(P))$; for $C$ smooth, $\ord_{\sigma P}(\sigma f)=\ord_Pf$. The fixed field of $G_k$ on $\bar k(C)$ is $k(C)$ (stated).
\item A morphism $\varphi$ defined over $k$ satisfies $\varphi(\sigma P)=\sigma(\varphi(P))$; hence $\varphi(C_1(L))\subseteq C_2(L)$, and the fibre $\varphi^{-1}(Q)$ over $Q\in C_2(k)$ is
$G_k$-stable.
\end{enumerate}
\end{proposition}

\begin{proof}
(i) $C(\bar k)$ is $G_k$-stable: if $F(P)=0$ for every $F$ in a set of generators of $I(C)$ with coefficients in $k$, then $F(\sigma P)=\sigma(F(P))=0$. The
fixed points are given by \Cref{prop:curves:projective-points}(i). (ii) $F^\sigma$ (apply $\sigma$ to the coefficients) is well defined on $\bar k(C)$ because
$I(C)$ is generated by polynomials with coefficients in $k$, so $\sigma$ maps $I(C)\otimes\bar k$ to itself. For $f=F/G$ with $G(P)\ne0$,
$F^\sigma(\sigma P)=\sigma(F(P))$ and $G^\sigma(\sigma P)=\sigma(G(P))\ne0$. So $f\mapsto\sigma f$ is a field automorphism of $\bar k(C)$ (semilinear over $\sigma$) that maps the
local ring at $P$ onto the local ring at $\sigma P$ and $\mf m_P$ onto $\mf m_{\sigma P}$; it therefore maps a uniformiser at $P$ to one at $\sigma P$ and preserves
valuations. That the fixed field is $k(C)$ is Hilbert's Theorem 90 for the vector space $\bar k(C)$ over $\bar k$ \cite[Chapter II, Section 5]{SilvermanAEC}.
(iii) If $\varphi=[f_0:\dots:f_n]$ with $f_i\in k(C_1)$, then $\sigma f_i=f_i$, and by (ii) $\sigma(\varphi(P))=[\sigma f_0(\sigma P):\dots]=\varphi(\sigma P)$. If some $g\in\bar k(C_1)^\times$ is needed to evaluate at $P$, that is $\varphi(P)=[(gf_0)(P):\dots:(gf_n)(P)]$, then
$\sigma(\varphi(P))=[(\sigma g\,f_0)(\sigma P):\dots:(\sigma g\,f_n)(\sigma P)]$, and $\sigma g$ does the same job at $\sigma P$, so this is $\varphi(\sigma P)$. The rest follows
from (i).
\end{proof}

\begin{warning}[a curve need not have a rational point]\label{warn:curves:rational-points}
The conic $X^2+Y^2+Z^2=0$ over $\Q$ has no point over $\R$, so $C(\Q)=\emptyset$; but $[1:i:0]\in C(\Q(i))$. It is a smooth curve of genus $0$
(\Cref{ex:curves:genus-examples}) that is isomorphic to $\PP^1$ over $\Q(i)$ (project from $[1:i:0]$, as in \Cref{cor:curves:rr-consequences}(d)) but not over
$\Q$, since an isomorphism over $\Q$ would map $\PP^1(\Q)\ne\emptyset$ into $C(\Q)$ (\Cref{prop:curves:galois-action}(iii)). Selmer's cubic $3X^3+4Y^3+5Z^3=0$
is smooth of genus one with no point over $\Q$, although it has points over $\R$ and over every $\Q_p$ (stated; see \cite[Chapter X]{SilvermanAEC}, section not verified).
So \Cref{thm:curves:genus-one-weierstrass,prop:curves:genus-one-pic} assume a rational point, and the paper's curves have one: $O$ on $E$, and the cusp
$\infty$ on $X_0(N)$.
\end{warning}

\begin{paperlink}\label{pl:curves:galois}
Section 1.1 of the paper: $\rho_{E,\ell}$ is the action of $G_\Q$ on $E[\ell^n]\subseteq E(\Qbar)$ of \Cref{prop:curves:galois-action}(i), extended to
$T_\ell E=\varprojlim E[\ell^n]$ (\Cref{ch:ellcurves,ch:galrep}). This action respects the group law because the group law is given by rational functions with
coefficients in $\Q$ (\Cref{prop:curves:genus-one-pic}). Proposition prop:ns: $G_\Q$ acts on the eight points of $E[3]\setminus\{O\}$ and hence on the four
roots of the $3$-division polynomial, the $x$-coordinates of those points, by \Cref{prop:curves:galois-action}(ii), (iii) applied to the function $x$,
which is defined over $\Q$: $x(\sigma P)=\sigma(x(P))$.
\end{paperlink}

\section{Divisors, differentials, the Picard group}\label{sec:curves:divisors}

Divisors record zeros and poles; the Picard group is the group of divisors modulo divisors of functions, and its degree-zero part is the
Jacobian of \Cref{sec:curves:jacobian}. Differentials give the canonical class and, through Riemann-Roch, the genus. In the paper these
objects stand behind the group law of $E$ (\Cref{prop:curves:genus-one-pic}), behind the Jacobian $J_0(N)$ of Remark rem:free, and behind the
differential $\omega_E$ whose pullback $\varphi_E^*\omega_E=c_E\cdot2\pi i\,f(z)\,dz$ appears in the proof of Proposition prop:pullback (\Cref{ch:moddeg}).
Throughout, $C$ is a smooth projective curve over $k$.

\begin{definition}[divisors]\label{def:curves:divisor}
$\Div(C)$ is the free abelian group on $C(\bar k)$; its elements are finite sums $D=\sum n_P(P)$. The \emph{degree} is $\deg D=\sum n_P$; $D\ge0$
(\emph{effective}) means all $n_P\ge0$; $\operatorname{supp}D=\{P:n_P\ne0\}$. $G_k$ acts by $\sigma D=\sum n_P(\sigma P)$, and $\Div_k(C)=\Div(C)^{G_k}$ is the group
of divisors \emph{defined over $k$}; for instance the sum of the points of a $G_k$-orbit. The \emph{divisor} of $f\in\bar k(C)^\times$ is
$\operatorname{div}(f)=\sum_P\ord_P(f)(P)$, a finite sum by \Cref{prop:curves:function-field}(ii); its positive and negative parts are the divisors of zeros
$\operatorname{div}_0(f)$ and of poles $\operatorname{div}_\infty(f)$. Divisors of functions are \emph{principal}; they form a subgroup $\operatorname{Prin}(C)$, since
$\operatorname{div}(fg)=\operatorname{div}(f)+\operatorname{div}(g)$.
\end{definition}

\begin{proposition}[principal divisors]\label{prop:curves:principal}
Let $f\in\bar k(C)^\times$.
\begin{enumerate}[label=(\roman*)]
\item $\operatorname{div}(f)=0$ iff $f\in\bar k^\times$.
\item $\deg\operatorname{div}(f)=0$; more precisely, for $f$ nonconstant, $\deg\operatorname{div}_0(f)=\deg\operatorname{div}_\infty(f)=[\bar k(C):\bar k(f)]$.
\item $\operatorname{div}(\sigma f)=\sigma\operatorname{div}(f)$ for $\sigma\in G_k$.
\item On $\PP^1$ with coordinate $x$: $\operatorname{div}(x-a)=(a)-(\infty)$, and $\operatorname{div}\bigl(c\prod(x-a_i)^{n_i}\bigr)=\sum n_i(a_i)-(\sum n_i)(\infty)$.
\end{enumerate}
\end{proposition}

\begin{proof}
(i) If $\operatorname{div}(f)=0$, $f$ has no pole, so it is constant (\Cref{prop:curves:function-field}(iii)); constants have divisor $0$. (ii) is
\Cref{cor:curves:degree-three-ways}(vi), which uses the degree formula of \Cref{sec:curves:morphisms}; over $\C$ a second proof is
\Cref{cor:curves:argument-principle}. (iii) is \Cref{prop:curves:galois-action}(ii). (iv) At $a\in\bar k$ the function $x-a$ is a uniformiser; at $\infty$ the
function $t=1/x$ is one (in the chart $[1:t]$), and $x-a=t^{-1}(1-at)$ with $1-at$ a unit at $\infty$, so $\ord_\infty(x-a)=-1$. The second formula follows
by additivity.
\end{proof}

\emph{Example: divisors on a cubic.}
Let $C$ be the smooth cubic $y^2=h(x)=(x-e_1)(x-e_2)(x-e_3)$, $e_i$ distinct, $\operatorname{char}k\ne2$, $O=[0:1:0]$ and $T_i=(e_i,0)$. The orders
needed are computed in the remarks after \Cref{thm:curves:smooth-dvr} and in the proof of \Cref{prop:curves:differentials}(iv) below: $y$ is a uniformiser at
each $T_i$, $\ord_O y=-3$, $\ord_O x=-2$, and at an affine point with $y\ne0$ the function $x-a$ is a uniformiser. Hence
\[
\operatorname{div}(y)=(T_1)+(T_2)+(T_3)-3(O),\qquad \operatorname{div}(x-e_1)=2(T_1)-2(O),\qquad \operatorname{div}(x-a)=(P)+(P')-2(O),
\]
the last for $a$ not a root of $h$, with $P,P'=(a,\pm b)$ and $b^2=h(a)\ne0$. All three have degree zero, as \Cref{prop:curves:principal}(ii) requires, and
the degrees of the zero divisors, $3$ and $2$, are the field degrees $[\bar k(C):\bar k(y)]$ and $[\bar k(C):\bar k(x)]$ (the second because $y$ satisfies a
quadratic equation over $\bar k(x)$ and is not in $\bar k(x)$, $h$ not being a square). Note also $\operatorname{div}(y^2)=2\operatorname{div}(y)=\sum_i\operatorname{div}(x-e_i)$,
as it must be, since $y^2=\prod(x-e_i)$.

\begin{definition}[linear equivalence; the Picard group]\label{def:curves:picard}
$D\sim D'$ (\emph{linearly equivalent}) iff $D-D'$ is principal. $\Pic(C)=\Div(C)/\operatorname{Prin}(C)$. By \Cref{prop:curves:principal}(ii) the degree
induces $\deg\colon\Pic(C)\to\Z$, with kernel $\Pic^0(C)=\Div^0(C)/\operatorname{Prin}(C)$, $\Div^0$ the divisors of degree zero.
\end{definition}

\begin{proposition}[the Picard group]\label{prop:curves:picard}
\leavevmode
\begin{enumerate}[label=(\roman*)]
\item $\Pic^0(\PP^1)=0$, and $\deg\colon\Pic(\PP^1)\to\Z$ is an isomorphism.
\item $G_k$ acts on $\Pic(C)$ and $\Pic^0(C)$.
\item The natural map $\Div^0_k(C)/\operatorname{div}(k(C)^\times)\to\Pic^0(C)^{G_k}$ is injective; it is bijective when $C(k)\ne\emptyset$ (stated).
\end{enumerate}
\end{proposition}

\begin{proof}
(i) Every $f\in\bar k(\PP^1)^\times=\bar k(x)^\times$ is $c\prod(x-a_i)^{n_i}$ (factor numerator and denominator over $\bar k$), so by
\Cref{prop:curves:principal}(iv) a divisor $\sum n_i(a_i)+n_\infty(\infty)$ of degree zero, that is with $n_\infty=-\sum n_i$, is $\operatorname{div}\prod(x-a_i)^{n_i}$.
(ii) $\sigma$ preserves degrees and, by \Cref{prop:curves:principal}(iii), principal divisors.
(iii) Let $D\in\Div^0_k(C)$ with $D=\operatorname{div}(f)$ for some $f\in\bar k(C)^\times$. Choose a finite Galois extension $L/k$ containing the coefficients of
$f$. For $\sigma\in G_k$, $\operatorname{div}(\sigma f)=\sigma D=D$, so $\sigma f=c_\sigma f$ with $c_\sigma\in\bar k^\times$ by \Cref{prop:curves:principal}(i); in fact $c_\sigma\in L^\times$ and
$c_\sigma$ depends only on $\sigma|_L$. From $\sigma\tau f=\sigma(c_\tau)c_\sigma f$ we get $c_{\sigma\tau}=\sigma(c_\tau)c_\sigma$: a $1$-cocycle of $\Gal(L/k)$ with values in $L^\times$. By
Hilbert's Theorem 90 (\Cref{ch:galois}) there is $b\in L^\times$ with $c_\sigma=\sigma(b)/b$; then $f/b$ is fixed by $G_k$, so lies in $k(C)$
(\Cref{prop:curves:galois-action}(ii)), and $D=\operatorname{div}(f/b)$. This is injectivity. Surjectivity when $C(k)\ne\emptyset$ uses the vanishing of an
obstruction in the Brauer group of $k$; we do not prove it, see \cite[Section 1]{MilneJacobian86} (section number not verified) and
\Cref{ch:galcoh} for Hilbert 90 in cohomological form.
\end{proof}

\begin{definition}[differentials]\label{def:curves:differentials}
The space $\Omega_C$ of (meromorphic) \emph{differentials} is the $\bar k(C)$-vector space generated by symbols $df$, $f\in\bar k(C)$, subject to
$d(f+g)=df+dg$, $d(fg)=f\,dg+g\,df$ and $da=0$ for $a\in\bar k$. For a uniformiser $t$ at $P$, every $\omega\in\Omega_C$ is $\omega=g\,dt$ with a unique
$g=\omega/dt\in\bar k(C)$, and $\ord_P\omega:=\ord_Pg$. Then $\operatorname{div}(\omega)=\sum_P\ord_P(\omega)(P)$; $\omega$ is \emph{regular} (or holomorphic) if
$\operatorname{div}(\omega)\ge0$, and the regular differentials form a $\bar k$-vector space.
\end{definition}

The definition is the algebraic shadow of the expressions $g(z)\,dz$ of complex analysis: the rules are those of the differential calculus, and $\omega/dt$
is ``the coefficient of $dt$'' in a local coordinate. Differentials are introduced because, unlike functions, they can be regular everywhere on a projective
curve without being constant multiples of each other, and the number of independent regular ones turns out to be the genus (\Cref{cor:curves:rr-consequences}(a)).
For instance on $\PP^1$, $d(x^2)=2x\,dx$, and at $a\ne0$ with uniformiser $x-a$, $d(x^2)/d(x-a)=2x$ has order $0$ (if $\operatorname{char}k\ne2$), while at $0$ it has
order $1$.

\begin{proposition}[differentials; the canonical class]\label{prop:curves:differentials}
\leavevmode
\begin{enumerate}[label=(\roman*)]
\item $\Omega_C$ is one-dimensional over $\bar k(C)$, and $df$ is a basis iff $\bar k(C)/\bar k(f)$ is a finite separable extension (in characteristic $p$,
$d(f^p)=pf^{p-1}df=0$). In particular $dt$ is a basis for a uniformiser $t$ (stated).
\item $\ord_P\omega$ does not depend on the uniformiser $t$; if $f$ is regular at $P$, so is $df/dt$; and $\ord_P\omega=0$ for all but finitely many $P$ (stated).
\item $\operatorname{div}(f\omega)=\operatorname{div}(f)+\operatorname{div}(\omega)$. Hence the divisors of nonzero differentials form one linear equivalence class, the
\emph{canonical class}; any $\operatorname{div}(\omega)$, $\omega\ne0$, is a \emph{canonical divisor} $K_C$.
\item Examples. On $\PP^1$, $\operatorname{div}(dx)=-2(\infty)$, so $\PP^1$ has no nonzero regular differential. On the smooth cubic $y^2=h(x)$ with $h$
squarefree of degree $3$ and $\operatorname{char}k\ne2$, $\operatorname{div}(dx/y)=0$.
\end{enumerate}
\end{proposition}

\begin{proof}
(i) and (ii) are \cite[Chapter II, Propositions 4.2 and 4.3]{SilvermanAEC}; we do not prove them. (iii) $f\omega=(fg)\,dt$, and $\ord_P(fg)=\ord_Pf+\ord_Pg$.
If $\omega,\omega'\ne0$ then $\omega'=f\omega$ for some $f\in\bar k(C)^\times$ by (i), so $\operatorname{div}\omega'-\operatorname{div}\omega=\operatorname{div}f$; conversely every divisor
linearly equivalent to $\operatorname{div}\omega$ is $\operatorname{div}(f\omega)$.

(iv) On $\PP^1$, at $a\in\bar k$, $dx=d(x-a)$ has order $0$; at $\infty$, with $t=1/x$, $dx=d(t^{-1})=-t^{-2}dt$ has order $-2$. A regular differential
$f\,dx$ would need $\operatorname{div}f\ge2(\infty)$, of positive degree, impossible by \Cref{prop:curves:principal}(ii) unless $f=0$.

On $y^2=h(x)$ let $P=(a,b)$. If $b\ne0$, then $f_y=2b\ne0$, so $x-a$ is a uniformiser (\Cref{thm:curves:smooth-dvr}), $\ord_Pdx=0$ and $y$ is a unit:
$\ord_P(dx/y)=0$. If $b=0$, then $h(a)=0$ and $f_x(P)=-h'(a)\ne0$ as $h$ is squarefree, so $y$ is a uniformiser; differentiating $y^2=h(x)$ gives
$2y\,dy=h'(x)\,dx$, so $dx=\frac{2y}{h'(x)}dy$ with $h'(x)$ a unit at $P$: $\ord_Pdx=1$ and $\ord_P(dx/y)=0$. At the point $O=[0:1:0]$ at infinity
(write $h=x^3+c_2x^2+c_1x+c_0$ after scaling) use the chart $Y=1$, with coordinates $u=x/y$, $v=1/y$; the equation $Y^2Z=X^3+c_2X^2Z+c_1XZ^2+c_0Z^3$
becomes $v(1-c_2u^2-c_1uv-c_0v^2)=u^3$. Its partial derivative in $v$ at $(0,0)$ is $1$, so $O$ is smooth and $u$ is a uniformiser (the tangent line
is $v=0$). The bracket is a unit at $O$, so $\ord_Ov=3$; hence $\ord_Oy=-3$ and $\ord_Ox=\ord_O(u/v)=-2$. Write $x=u^{-2}w$ with $w=u^3/v$ a unit,
$w(O)=1$; then $dx=(-2u^{-3}w+u^{-2}w')du$ with $w'=dw/du$ regular (by (ii)), and $-2w+uw'$ takes the value $-2\ne0$ at $O$. So $\ord_Odx=-3$ and
$\ord_O(dx/y)=0$.
\end{proof}

The translation invariance of $dx/y$, the \emph{invariant differential}, is proved in \Cref{ch:ellcurves}; for a general Weierstrass equation the
differential is $dx/(2y+a_1x+a_3)$, with the same divisor (\Cref{ex:curves:genus-examples}(b)).

\section{Riemann-Roch and its consequences}\label{sec:curves:riemann-roch}

Riemann-Roch is the one deep theorem of the algebraic side; it is stated with a reference and the shape of its proof, and everything after
it in this section is proved. Its consequences give the genus as the number of independent regular differentials, show that a genus-one curve
with a rational point is a Weierstrass cubic (\Cref{thm:curves:genus-one-weierstrass}), and identify such a curve with its own $\Pic^0$
(\Cref{prop:curves:genus-one-pic}). In the paper these are the facts that ``an elliptic curve over $\Q$ is a Weierstrass equation'' (Section 1 and
the Appendix) and that $E=\Pic^0(E)$ is its own Jacobian, which underlies the factorisation $X_0(N)\to J_0(N)\to E$ used in Proposition prop:pullback.

\begin{definition}[Riemann-Roch spaces]\label{def:curves:riemann-roch-space}
For $D\in\Div(C)$, $\mathcal L(D)=\{f\in\bar k(C)^\times:\operatorname{div}(f)+D\ge0\}\cup\{0\}$, the functions whose poles are bounded by $D$ and which vanish where $D$
requires; $l(D)=\dim_{\bar k}\mathcal L(D)$.
\end{definition}

\begin{proposition}[Riemann-Roch spaces]\label{prop:curves:riemann-roch-space}
\leavevmode
\begin{enumerate}[label=(\roman*)]
\item $\mathcal L(D)$ is a $\bar k$-vector space; $\mathcal L(0)=\bar k$; $\mathcal L(D)=0$ if $\deg D<0$; if $D'=D+\operatorname{div}h$ then $f\mapsto fh^{-1}$ is an isomorphism
$\mathcal L(D)\to\mathcal L(D')$.
\item $l(D+(P))\le l(D)+1$. Hence $l(D)\le\deg D+1$ for $D\ge0$, and $l(D)$ is finite for every $D$.
\item $l(D)$ depends only on the class of $D$ in $\Pic(C)$; so does $l(K_C-D)$.
\item If $D$ is defined over $k$, $\mathcal L(D)$ has a basis in $k(C)$, and $\mathcal L(D)\cap k(C)$ is a $k$-space of dimension $l(D)$ (stated, except the
second clause).
\end{enumerate}
\end{proposition}

\begin{proof}
(i) If $\operatorname{div}f\ge-D$ and $\operatorname{div}g\ge-D$, then $\ord_P(f+g)\ge\min(\ord_Pf,\ord_Pg)\ge-n_P$ for every $P$; scalars do not change divisors.
$\mathcal L(0)$ is the space of functions without poles, the constants (\Cref{prop:curves:function-field}(iii)). If $f\in\mathcal L(D)$, $f\ne0$, then
$\operatorname{div}f+D\ge0$ has degree $\deg D$ (\Cref{prop:curves:principal}(ii)), so $\deg D\ge0$. Finally $\operatorname{div}(fh^{-1})+D'=\operatorname{div}f+D$.

(ii) Let $n$ be the coefficient of $P$ in $D$ and $t$ a uniformiser at $P$. For $f\in\mathcal L(D+(P))$, $\ord_Pf\ge-n-1$, so $t^{n+1}f$ is regular at $P$, and
$\lambda(f)=(t^{n+1}f)(P)$ is a linear map $\mathcal L(D+(P))\to\bar k$. Its kernel is $\{f:\ord_Pf\ge-n\}=\mathcal L(D)$. So $l(D+(P))\le l(D)+1$. For $D\ge0$ add the
points one at a time starting from $l(0)=1$. For general $D$, $\mathcal L(D)\subseteq\mathcal L(D^+)$, where $D^+$ keeps the positive coefficients.

(iii) is (i), applied to $D$ and to $K_C-D$. (iv) The first clause is \cite[Chapter II, Proposition 5.8]{SilvermanAEC} (Hilbert 90 for vector spaces).
Given a basis $f_1,\dots,f_m$ in $k(C)$, if $f=\sum c_if_i\in k(C)$ then $\sigma f=f$ gives $\sum(\sigma c_i-c_i)f_i=0$, so $c_i\in k$: the $k$-span of the $f_i$ is
$\mathcal L(D)\cap k(C)$.
\end{proof}

\begin{theorem}[Riemann-Roch; stated]\label{thm:curves:riemann-roch}
Let $C$ be a smooth projective curve over $k$ and $K_C$ a canonical divisor. There is an integer $g\ge0$, the \emph{genus} of $C$, such that for every
$D\in\Div(C)$
\[
l(D)-l(K_C-D)=\deg D-g+1.
\]
\end{theorem}

We do not prove this; see \cite[Chapter II, Theorem 5.4]{SilvermanAEC}, which refers the proof to Hartshorne and Lang,
\cite[Chapter IV, Theorem 1.3]{Hartshorne}, \cite[Chapter 8]{Fulton} and \cite[Chapter VI]{Miranda} (chapter numbers of Miranda not verified).
\emph{Shape of the proof.} By \Cref{prop:curves:riemann-roch-space}(ii), $\deg D+1-l(D)$ is bounded above as $D$ grows (one compares $\mathcal L(D)$ with
$\mathcal L(D+n(P))$ and uses a function with a single pole); the maximum of this defect is the genus, which gives Riemann's inequality
$l(D)\ge\deg D+1-g$. The correction term $l(K_C-D)$ is the dimension of the space of obstructions to finding a function with prescribed principal parts
bounded by $D$ (the algebraic Mittag-Leffler problem, in the language of repartitions), and residues of differentials make that space dual to
$\mathcal L(K_C-D)$; this duality is Serre duality. Over $\C$ the obstruction space is the sheaf cohomology group $H^1(X,\mathcal O(D))$, the only appearance
of sheaf cohomology in the book; it is never used in a proof.

\begin{corollary}[consequences of Riemann-Roch]\label{cor:curves:rr-consequences}
Let $C$ have genus $g$.
\begin{enumerate}[label=(\alph*)]
\item $l(K_C)=g$: the genus is the dimension of the space of regular differentials.
\item $\deg K_C=2g-2$.
\item If $\deg D>2g-2$ then $l(D)=\deg D-g+1$.
\item $g(\PP^1)=0$. Conversely, a curve $C$ of genus $0$ with a point $P\in C(k)$ is isomorphic over $k$ to $\PP^1$. A curve with two distinct points $P$, $Q$
such that $(P)\sim(Q)$ has genus $0$.
\item If $g=1$: $l(D)=\deg D$ whenever $\deg D\ge1$; $K_C\sim0$; and a nonzero regular differential has no zeros.
\end{enumerate}
\end{corollary}

\begin{proof}
(a) Take $D=0$: $1-l(K_C)=1-g$. The map $f\mapsto f\omega$ identifies $\mathcal L(\operatorname{div}\omega)$ with the regular differentials, because
$\operatorname{div}(f\omega)=\operatorname{div}f+\operatorname{div}\omega$ (\Cref{prop:curves:differentials}(iii)) and every differential is $f\omega$. (b) Take $D=K_C$:
$g-1=\deg K_C-g+1$. (c) $\deg(K_C-D)<0$, so $l(K_C-D)=0$ by \Cref{prop:curves:riemann-roch-space}(i).

(d) $\deg\operatorname{div}(dx)=-2$ on $\PP^1$ (\Cref{prop:curves:differentials}(iv)), so $2g-2=-2$ by (b). Let $g(C)=0$ and $P\in C(k)$. By (c),
$l((P))=1-0+1=2$, and by \Cref{prop:curves:riemann-roch-space}(iv) there is a nonconstant $f\in\mathcal L((P))\cap k(C)$. Its divisor of poles is at most $(P)$ and is
nonzero (\Cref{prop:curves:function-field}(iii)), so it is exactly $(P)$. By \Cref{prop:curves:principal}(ii) the morphism $f\colon C\to\PP^1$, defined over $k$,
has degree $[\bar k(C):\bar k(f)]=\deg\operatorname{div}_\infty f=1$, so it is an isomorphism (\Cref{thm:curves:equivalence}). If $(P)-(Q)=\operatorname{div}f$ with $P\ne Q$,
the same argument applies to $f$, whose divisor of poles is $(Q)$: $C\simeq\PP^1$ over $\bar k$, and the genus, defined through $l$ and degrees, is invariant
under isomorphism, so $g=0$.

(e) By (b), $\deg K_C=0$, so (c) gives $l(D)=\deg D$ for $\deg D\ge1$. By (a), $l(K_C)=1$: there is $f\ne0$ with $\operatorname{div}f+K_C\ge0$; this effective divisor
has degree $0$, so it is $0$, and $K_C\sim0$. If $\omega$ is regular and nonzero, $\operatorname{div}\omega$ is effective of degree $0$, hence $0$.
\end{proof}

\begin{theorem}[genus of a smooth plane curve]\label{thm:curves:plane-genus}
A smooth plane curve of degree $d$ over $k$, in any characteristic, has genus
\[
g=\frac{(d-1)(d-2)}2.
\]
\end{theorem}

\begin{proof}[Sketch of proof]
Choose coordinates in general position: so that $f_y$ is not identically zero, the line at infinity $Z=0$ is tangent to $C$ nowhere (so it meets $C$ in $d$
distinct points, by \Cref{thm:curves:bezout}), and $[0:1:0]\notin C$. On the affine part $f(x,y)=0$ put
$\omega_0=dx/f_y=-dy/f_x$ (the two agree because $f_x\,dx+f_y\,dy=d(f)=0$ on $C$). At an affine point with $f_y\ne0$, $x-a$ is a uniformiser and
$\ord\omega_0=0$; at an affine point with $f_y=0$, $f_x\ne0$ by smoothness, $y-b$ is a uniformiser and $\ord\omega_0=\ord(-dy/f_x)=0$. So $\omega_0$ is regular
and nowhere zero on the affine part. At infinity we compute in a chart, as in \Cref{prop:curves:differentials}(iv). Since $[0:1:0]\notin C$, every point
$P$ of $C$ on the line $Z=0$ has $X\ne0$, so it lies in the chart $X=1$, with affine coordinates $u=Y/X$, $v=Z/X$ and equation $g(u,v)=F(1,u,v)=0$. In
these coordinates $x=X/Z=1/v$ and $y=Y/Z=u/v$. The line $Z=0$ is $v=0$; it is not the tangent at $P$, so $g_u(P)\ne0$ (the tangent is
$g_u(P)(u-u(P))+g_v(P)v=0$, \Cref{prop:curves:euler}(ii)), and $v$ is a uniformiser at $P$ (\Cref{thm:curves:smooth-dvr}). Now $F_Y$ is homogeneous of degree
$d-1$, so
\[
f_y=F_Y(x,y,1)=F_Y(v^{-1},uv^{-1},1)=v^{1-d}F_Y(1,u,v)=v^{1-d}g_u(u,v),
\]
and $g_u$ is a unit at $P$: $\ord_Pf_y=1-d$. Also $dx=d(v^{-1})=-v^{-2}dv$ has order $-2$. Hence $\ord_P\omega_0=-2-(1-d)=d-3$ at each of the $d$ points
at infinity. As $Z/M$ (for a linear form $M$ not vanishing at these points) has a simple zero at each of them and no other zero, this reads
$\operatorname{div}(\omega_0)=(d-3)\cdot\operatorname{div}_0(Z/M)$, a divisor of degree $d(d-3)$ (\Cref{thm:curves:bezout}). By \Cref{cor:curves:rr-consequences}(b),
$2g-2=d(d-3)$, so $g=\frac{d^2-3d+2}2=\frac{(d-1)(d-2)}2$. For $d=3$ the form $\omega_0$ has no zeros or poles at all, as for $dx/(2y+a_1x+a_3)$ on a
Weierstrass cubic. One also finds that the regular
differentials are exactly the $h(x,y)\,\omega_0$ with $\deg h\le d-3$, a space of dimension $\binom{d-1}2$, in agreement with
\Cref{cor:curves:rr-consequences}(a). See \cite[Chapter 8]{Fulton} and \cite[Chapter IV, Section 1, and Chapter I, Section 7]{Hartshorne} (sections not
verified). The argument needs no hypothesis on the characteristic, because such coordinates always exist: the tangent lines of $C$ form at most a one-dimensional
family in the two-dimensional family of lines, so a general line is tangent nowhere; and $F_X$, $F_Y$, $F_Z$ are not all identically zero ($C$ being smooth), so
after a general linear change of coordinates $F_Y$, hence $f_y$, is not identically zero (in characteristic $p$ this can fail for special coordinates, when
$f\in k[x,y^p]$).
\end{proof}

\begin{example}[genera]\label{ex:curves:genus-examples}
\leavevmode
\begin{enumerate}[label=(\alph*)]
\item Lines and smooth conics have genus $0$ (\Cref{thm:curves:plane-genus} with $d=1,2$; for the conic, also \Cref{ex:curves:conic}).
\item Smooth plane cubics have genus $1$; in particular every smooth Weierstrass cubic, such as \curve{11a1}, \curve{20a1}, \curve{121a1}. If $\operatorname{char}k\ne2$ the
substitution $y\mapsto y-\frac12(a_1x+a_3)$ turns $y^2+a_1xy+a_3y=x^3+\dots$ into $y^2=h(x)$ and $dx/(2y+a_1x+a_3)$ into $\frac12dx/y$; so the differential
$\omega=dx/(2y+a_1x+a_3)$ has divisor $0$ by \Cref{prop:curves:differentials}(iv) (smoothness makes $h$ squarefree), and is regular and nowhere zero.
This gives $g=1$ again, from \Cref{cor:curves:rr-consequences}(b).
\item The Fermat quartic $X^4+Y^4+Z^4=0$ has genus $3$ ($\operatorname{char}k\ne2$): it is smooth by \Cref{ex:curves:plane-examples}(b), and
$\frac{(4-1)(4-2)}2=3$.
\item Hyperelliptic curves: the smooth projective model of $y^2=h(x)$, $h$ squarefree of degree $2g+1$ or $2g+2$, $\operatorname{char}k\ne2$, has genus $g$
(\Cref{cor:curves:rh-consequences}(c)); for instance $y^2=x^5-x$ has genus $2$, since $x^5-x=x(x-1)(x+1)(x^2+1)$ has degree $5=2\cdot2+1$ and
discriminant $-2^8$, nonzero for $\operatorname{char}k\ne2$. For $g\ge2$ the projective closure in $\PP^2$ is singular at infinity, so
\Cref{thm:curves:plane-genus} does not apply directly. Indeed, for $\deg h=n\ge4$ the homogenised equation $Y^2Z^{n-2}=Z^nh(X/Z)$ meets $Z=0$ only at
$[0:1:0]$, and in the chart $Y=1$ it reads $Z^{n-2}=(\text{leading coefficient})X^n+(\text{terms divisible by }Z)$, whose lowest-degree part at the origin
has degree $\min(n-2,n)=n-2\ge2$: all partial derivatives vanish there. Were the plane formula applicable it would give $\frac{(n-1)(n-2)}2$, which
for $n=5$ is $6$, not $2$.
\item Singular plane models have smaller genus: the nodal and cuspidal cubics are rational (\Cref{thm:curves:smooth-model} and
\Cref{ex:curves:plane-examples}). In general $g=\frac{(d-1)(d-2)}2-\sum_P\delta_P$, with a positive local invariant $\delta_P$ at each singular point (stated;
\cite[Chapter 7]{Fulton}, \cite[Chapter IV, Section 1]{Hartshorne}, sections not verified).
\item The modular curves are not given as smooth plane curves; their genera come from Riemann-Hurwitz (\Cref{cor:curves:rh-consequences}(f) and
\Cref{ch:modforms}): $g(X_0(11))=g(X_0(20))=1$, $g(X_0(121))=6$, $g(X_0(539))=49$. We computed these from the genus formula of \Cref{ch:modforms} and checked
them with PARI's \texttt{mfdim([N,2],1)}, which is $\dim S_2(\Gamma_0(N))=g(X_0(N))$ (\Cref{thm:curves:genus-dimension}).
\end{enumerate}
\end{example}

\begin{theorem}[a genus-one curve with a rational point is a Weierstrass cubic]\label{thm:curves:genus-one-weierstrass}
Let $C$ be a smooth projective curve of genus one over $k$ and $O\in C(k)$.
\begin{enumerate}[label=(\roman*)]
\item There are $x\in\mathcal L(2(O))$ and $y\in\mathcal L(3(O))$, in $k(C)$, with $\mathcal L(2(O))=\langle1,x\rangle$, $\mathcal L(3(O))=\langle1,x,y\rangle$, $\ord_Ox=-2$, $\ord_Oy=-3$.
\item After rescaling $x$ and $y$ by elements of $k^\times$ they satisfy a Weierstrass equation
\[
y^2+a_1xy+a_3y=x^3+a_2x^2+a_4x+a_6,\qquad a_i\in k.
\]
\item The morphism $\varphi=[x:y:1]\colon C\to\PP^2$ is an isomorphism onto the projective cubic $E$ defined by this equation, $E$ is smooth, and $\varphi(O)=[0:1:0]$.
\item Conversely, a smooth Weierstrass cubic has genus one, and $[0:1:0]$ is a point of it defined over $k$.
\end{enumerate}
\end{theorem}

\begin{proof}
(i) By \Cref{cor:curves:rr-consequences}(e), $l(n(O))=n$ for $n\ge1$; so $\mathcal L((O))=\bar k$ and each step $\mathcal L((n-1)(O))\subset\mathcal L(n(O))$, $n\ge2$, adds exactly
one dimension, that is, a function with a pole of order exactly $n$ at $O$ and no other pole. By \Cref{prop:curves:riemann-roch-space}(iv) the $k$-rational
parts $\mathcal L(n(O))\cap k(C)$ have $k$-dimension $n$, so we can choose $x\in\mathcal L(2(O))\cap k(C)$ not in $\mathcal L((O))$ and $y\in\mathcal L(3(O))\cap k(C)$ not in $\mathcal L(2(O))$.

(ii) The seven functions $1,x,y,x^2,xy,y^2,x^3$ have poles at $O$ only, of orders $0,2,3,4,5,6,6$, so they lie in $\mathcal L(6(O))\cap k(C)$, a $k$-space of
dimension $6$. They are linearly dependent over $k$. In a nontrivial relation the coefficients $A$ of $y^2$ and $B$ of $x^3$ are both nonzero: if one of them were
zero, the other term would be the only one with a pole of order $6$, so its coefficient would vanish too, and then $1,x,y,x^2,xy$, which have
distinct pole orders, would be linearly independent, forcing all coefficients to be zero. Write the relation as
$Ay^2+b_1xy+b_3y=Bx^3+b_2x^2+b_4x+b_6$ and substitute $x=ABx'$, $y=AB^2y'$. Then $Ay^2=A^3B^4y'^2$ and $Bx^3=A^3B^4x'^3$, so after dividing by
$A^3B^4$ both leading terms have coefficient $1$, and the relation becomes
\[
y'^2+\frac{b_1}{AB}x'y'+\frac{b_3}{A^2B^2}y'=x'^3+\frac{b_2}{AB^2}x'^2+\frac{b_4}{A^2B^3}x'+\frac{b_6}{A^3B^4},
\]
a Weierstrass equation with coefficients in $k$. The new functions $x'$, $y'$ still have poles of orders exactly $2$ and $3$ at $O$ and no other pole, so (i)
holds for them; we rename them $x$, $y$. (A scaling $x=\alpha x'$, $y=\beta y'$ makes both leading coefficients equal exactly when $A\beta^2=B\alpha^3$;
the choice $\alpha=AB$, $\beta=AB^2$ is one solution with $\alpha,\beta\in k^\times$.)

(iii) $\varphi$ is a morphism (\Cref{thm:curves:extend-morphism}) defined over $k$, with image in $E$. Writing $\varphi=[x/y:1:1/y]$ and using $\ord_O(x/y)=1$,
$\ord_O(1/y)=3$ gives $\varphi(O)=[0:1:0]$. The Weierstrass polynomial is irreducible over $\bar k$: it is monic of degree two in $y$, so no nonconstant polynomial in $x$ alone divides it, and a
factorisation would be $(y-p(x))(y-r(x))$ with $p+r=-a_1x-a_3$ and $pr=-(x^3+a_2x^2+a_4x+a_6)$. Then $\deg p+\deg r=3$ is odd, so $\deg p\ne\deg r$, and
$\deg(p+r)=\max(\deg p,\deg r)\ge2$, contradicting $\deg(p+r)\le1$. The image of $\varphi$ is an irreducible closed
subset of $E$ (stated: images of projective varieties are closed) and not a point (as $x$ is not constant), so it is all of $E$. By
\Cref{prop:curves:principal}(ii), $[\bar k(C):\bar k(x)]=\deg\operatorname{div}_\infty x=2$ and $[\bar k(C):\bar k(y)]=3$; hence $[\bar k(C):\bar k(x,y)]$ divides $2$ and $3$, so
$\bar k(C)=\bar k(x,y)=\varphi^*\bar k(E)$ and $\varphi$ is birational. If $E$ were singular, it would be birational to $\PP^1$: projecting from the singular point $S$,
a line through $S$ meets $E$ in $S$ with multiplicity at least two and in one further point (\Cref{thm:curves:bezout}), which gives a birational map
$E\dashrightarrow\PP^1$, as for the node in \Cref{thm:curves:smooth-model}. Then $C$ would be birational to $\PP^1$, hence isomorphic to it
(\Cref{thm:curves:equivalence}), and of genus $0$, a contradiction. So $E$ is smooth, $\varphi$ is a morphism of degree one between smooth projective curves,
and it is an isomorphism by \Cref{thm:curves:equivalence}.

(iv) Genus one is \Cref{ex:curves:genus-examples}(b); $[0:1:0]$ has coordinates in $k$.
\end{proof}

\begin{proposition}[a genus-one curve is its own $\Pic^0$]\label{prop:curves:genus-one-pic}
Let $C$ have genus one and $O\in C(k)$. The map
\[
\kappa\colon C(\bar k)\to\Pic^0(C),\qquad P\mapsto\text{the class of }(P)-(O),
\]
is a bijection, and $\kappa(\sigma P)=\sigma\kappa(P)$ for $\sigma\in G_k$. Hence $C(\bar k)$ carries a unique group structure, with identity $O$, for which $\kappa$ is an
isomorphism, and $G_k$ acts by group automorphisms.
\end{proposition}

\begin{proof}
Surjective: for $D\in\Div^0(C)$, $l(D+(O))=1$ by \Cref{cor:curves:rr-consequences}(e), so there is $f\ne0$ with $\operatorname{div}f+D+(O)\ge0$; this
effective divisor has degree one, so it is $(P)$ for a point $P$, and $D\sim(P)-(O)$. Injective: if $(P)-(O)\sim(Q)-(O)$ then $(P)\sim(Q)$, and $P\ne Q$ would
force genus $0$ (\Cref{cor:curves:rr-consequences}(d)). Equivariance: $\sigma((P)-(O))=(\sigma P)-(O)$ because $\sigma O=O$, and $\sigma$ preserves principal divisors
(\Cref{prop:curves:principal}(iii)).
\end{proof}

\Cref{ch:ellcurves} identifies this group law with the chord-tangent law of a Weierstrass cubic: if $P$, $Q$, $R$ are the three intersections of a line
$L=0$ with $E$, then $(P)+(Q)+(R)-3(O)=\operatorname{div}(L/Z)$, so $\kappa(P)+\kappa(Q)+\kappa(R)=0$. Associativity of the chord-tangent law, which is awkward to check
directly, is then automatic.

\begin{paperlink}\label{pl:curves:weierstrass}
Section 1 and the Appendix of the paper: every elliptic curve is given by a Weierstrass equation, as \texttt{ellinit([1,1,1,-30,-76])} for
\curve{121a1}; \Cref{thm:curves:genus-one-weierstrass} is why this loses nothing. \Cref{prop:curves:genus-one-pic} says $E=\Pic^0(E)=J(E)$: an elliptic curve
is its own Jacobian (\Cref{def:curves:jacobian}), which is what makes $\varphi_E$ factor through $J_0(N)$ (\Cref{prop:curves:jacobian-functoriality}(iv)) and
lets the proof of Proposition prop:pullback speak of ``the map $J_0(N)\to E$ induced by $\varphi$''.
\end{paperlink}

\section[Morphisms and their degree; the Frobenius morphism]{Morphisms and their degree; separability; the Frobenius morphism}\label{sec:curves:morphisms}

The modular degree of the paper is the degree of the morphism $\varphi_E\colon X_0(N)\to E$. This section defines the degree of a nonconstant
morphism of curves and proves that its descriptions agree: the degree of the extension of function fields, the number of preimages of a general
point, the factor by which pullback multiplies degrees of divisors, and (over $\C$, in \Cref{sec:curves:riemann-surfaces,sec:curves:cohomology}) the
topological degree on $H^2$. It also proves $\varphi_*\varphi^*=[\deg\varphi]$ on divisor classes, the identity behind ``$E\to J_0(N)\to E$ is multiplication by
$\deg\varphi_E$'' (\Cref{ch:moddeg}), and introduces the Frobenius morphism used in \Cref{ch:ellcurves,ch:moddeg}.

\begin{theorem}[nonconstant morphisms]\label{thm:curves:constant-or-surjective}
A morphism $\varphi\colon C_1\to C_2$ of curves is constant or surjective. If it is nonconstant, $\varphi^*\colon\bar k(C_2)\to\bar k(C_1)$, $f\mapsto f\circ\varphi$, is a
well-defined injective homomorphism of fields fixing $\bar k$, and $\bar k(C_1)$ is a finite extension of $\varphi^*\bar k(C_2)$; if $\varphi$ is defined over $k$, the same
holds for $k(C_2)\to k(C_1)$.
\end{theorem}

\begin{proof}
Surjectivity is \cite[Chapter II, Theorem 2.3]{SilvermanAEC} (it uses that the image of a projective curve is closed); we do not prove it. Let $\varphi$ be
nonconstant, hence surjective. For $f\in\bar k(C_2)$, $f$ is regular off a finite set $S$ (\Cref{prop:curves:function-field}(ii)), and $\varphi^{-1}(S)$ is a proper
closed subset of $C_1$ (it is not all of $C_1$, as $\varphi$ is surjective and $C_2\ne S$), so $f\circ\varphi$ is a rational function on $C_1$. The map $f\mapsto f\circ\varphi$
respects sums and products, so it is a field homomorphism, injective as every field homomorphism is. Both fields are finitely generated of
transcendence degree one over $\bar k$, and $\varphi^*\bar k(C_2)$ contains an element transcendental over $\bar k$ (the pullback of a nonconstant function);
so $\bar k(C_1)$ is algebraic and finitely generated over $\varphi^*\bar k(C_2)$, hence finite. Over $k$ the argument is the same.
\end{proof}

\begin{definition}[degree, ramification index, pullback and pushforward]\label{def:curves:degree}
Let $\varphi\colon C_1\to C_2$ be a nonconstant morphism of smooth curves.
\begin{enumerate}[label=(\roman*)]
\item $\deg\varphi=[\bar k(C_1):\varphi^*\bar k(C_2)]$; $\deg_s\varphi$ and $\deg_i\varphi$ are the separable and inseparable degrees of this extension (\Cref{ch:galois}), so
$\deg\varphi=\deg_s\varphi\cdot\deg_i\varphi$; $\varphi$ is \emph{separable} (\emph{purely inseparable}) if the extension is. A constant map has degree $0$.
\item The \emph{ramification index} of $\varphi$ at $P$ is $e_\varphi(P)=\ord_P(\varphi^*t_{\varphi(P)})\ge1$, for a uniformiser $t_{\varphi(P)}$ at $\varphi(P)$; it does not depend on the
choice, since two uniformisers differ by a unit $u$ and $\varphi^*u$ is a unit at $P$.
\item $\varphi^*\colon\Div(C_2)\to\Div(C_1)$, $(Q)\mapsto\sum_{\varphi(P)=Q}e_\varphi(P)(P)$; $\varphi_*\colon\Div(C_1)\to\Div(C_2)$, $(P)\mapsto(\varphi(P))$; both extended linearly.
\end{enumerate}
\end{definition}

\begin{theorem}[the degree formula]\label{thm:curves:degree-formula}
For a nonconstant morphism $\varphi\colon C_1\to C_2$ of smooth projective curves and every $Q\in C_2(\bar k)$,
\[
\sum_{P\in\varphi^{-1}(Q)}e_\varphi(P)=\deg\varphi.
\]
\end{theorem}

\begin{proof}[Sketch of proof]
Let $A=\bar k[C_2]_Q$, a discrete valuation ring with uniformiser $t=t_Q$, and $B$ the integral closure of $A$ in $\bar k(C_1)$ (identifying $\bar k(C_2)$ with its
image under $\varphi^*$). Two facts are taken from the literature: $B$ is a finitely generated $A$-module (finiteness of integral closure for finitely generated
algebras over a field; \cite[Chapter 5]{AtiyahMacdonald}), and the maximal ideals of $B$ are the $\mf m_P\cap B$ for the points $P\in\varphi^{-1}(Q)$, with $B_{\mf m_P\cap B}=\bar k[C_1]_P$
(\Cref{thm:curves:smooth-model}). Then $B$ is a semilocal Dedekind domain, free over the principal ideal domain $A$ (torsion-free and finitely generated,
\Cref{ch:modules}) of rank $n=\deg\varphi$, and $tB=\prod_P(\mf m_P\cap B)^{e_\varphi(P)}$ by the definition of $e_\varphi(P)$. The argument of \Cref{ch:ant} for $\sum e_if_i=n$
now applies: $\dim_{\bar k}B/tB=n$ by freeness, and by the Chinese remainder theorem $B/tB\simeq\prod_PB/(\mf m_P\cap B)^{e_\varphi(P)}$, a product of spaces of
dimensions $e_\varphi(P)$, because every residue field is $\bar k$ (all residue degrees are $1$). See \cite[Chapter II, Proposition 2.6]{SilvermanAEC} and
\cite[Chapter II, Proposition 6.9]{Hartshorne}. In the language of physics: the number of sheets of a branched cover, counted with multiplicity, does not jump.
\end{proof}

The simplest example shows every feature. Let $\varphi\colon\PP^1\to\PP^1$, $z\mapsto z^n$, with $n\ge1$. Then $\varphi^*\bar k(w)=\bar k(z^n)$, and $z$ is a root of
$T^n-z^n$, which is irreducible over $\bar k(z^n)$ (Eisenstein at the prime $z^n$ of $\bar k[z^n]$), so $\deg\varphi=n$. Suppose $\operatorname{char}k\nmid n$. Over
$a\ne0,\infty$ the fibre is the $n$ distinct roots of $z^n=a$, and at each of them $\varphi^*(w-a)=z^n-a$ has a simple zero: $n$ points with $e=1$. Over $0$ the
fibre is the single point $0$, with $e_\varphi(0)=\ord_0(z^n)=n$; over $\infty$, with uniformiser $1/w$, $e_\varphi(\infty)=\ord_\infty(z^{-n})=n$. In each case
$\sum e_\varphi(P)=n$. If instead $n=p=\operatorname{char}k$, then $z^p-a=(z-a^{1/p})^p$, so every fibre is a single point with $e=p$: the map is a bijection on
points, of degree $p$, and $\deg_s\varphi=1$; this is the Frobenius morphism of \Cref{prop:curves:separability-frobenius}.

\begin{corollary}[the degree in three ways]\label{cor:curves:degree-three-ways}
Let $\varphi\colon C_1\to C_2$ be a nonconstant morphism of smooth projective curves.
\begin{enumerate}[label=(\roman*)]
\item $\deg\varphi=[\bar k(C_1):\varphi^*\bar k(C_2)]$ (the definition).
\item $\#\varphi^{-1}(Q)=\deg_s\varphi$ for all but finitely many $Q\in C_2(\bar k)$; if $\varphi$ is separable (always in characteristic $0$), $\#\varphi^{-1}(Q)=\deg\varphi$ for every $Q$
outside the finite set of branch points (\Cref{def:curves:ramification}).
\item $\deg\varphi^*D=\deg\varphi\cdot\deg D$; $\varphi_*\varphi^*D=(\deg\varphi)D$; $\varphi^*\operatorname{div}(f)=\operatorname{div}(\varphi^*f)$; and $\varphi_*\operatorname{div}(h)=\operatorname{div}(\Nm h)$, with $\Nm$ the norm
from $\bar k(C_1)$ to $\varphi^*\bar k(C_2)$ (stated). Hence $\varphi^*$ and $\varphi_*$ induce homomorphisms $\varphi^*\colon\Pic(C_2)\to\Pic(C_1)$ and $\varphi_*\colon\Pic(C_1)\to\Pic(C_2)$,
preserving $\Pic^0$, and $\varphi_*\varphi^*=[\deg\varphi]$ on $\Pic^0(C_2)$.
\item Over $\C$: $\deg\varphi$ is the analytic degree of the holomorphic map $\varphi$ (\Cref{thm:curves:holomorphic-maps}(c), \Cref{thm:curves:riemann-existence}), and
$\varphi^*$ on $H^2(\cdot,\Z)=\Z$ is multiplication by $\deg\varphi$ (\Cref{prop:curves:transfer}(iv)).
\item $\deg(\psi\circ\varphi)=\deg\psi\cdot\deg\varphi$; and $\deg\varphi=1$ iff $\varphi$ is an isomorphism.
\item For nonconstant $f\in\bar k(C)$, the morphism $f\colon C\to\PP^1$ has degree $[\bar k(C):\bar k(f)]=\deg\operatorname{div}_0(f)=\deg\operatorname{div}_\infty(f)$. Hence
$\deg\operatorname{div}(f)=0$, and a function without poles is constant.
\end{enumerate}
\end{corollary}

\begin{proof}
(iii) The first two identities follow from \Cref{thm:curves:degree-formula}: $\deg\varphi^*(Q)=\sum_{P\mapsto Q}e_\varphi(P)=\deg\varphi$, and
$\varphi_*\varphi^*(Q)=\sum_{P\mapsto Q}e_\varphi(P)(Q)=(\deg\varphi)(Q)$. For the third, let $Q=\varphi(P)$, $e=e_\varphi(P)$, $m=\ord_Qf$, and write $f=ut_Q^m$ with $u$ a unit at $Q$. Then
$\varphi^*f=(\varphi^*u)(\varphi^*t_Q)^m$, with $\varphi^*u$ a unit at $P$ and $\ord_P\varphi^*t_Q=e$, so $\ord_P\varphi^*f=em$: the coefficient of $P$ in $\varphi^*\operatorname{div}(f)$.
The norm formula is \cite[Chapter II, Section 3]{SilvermanAEC} (stated). So $\varphi^*$ and $\varphi_*$ map principal divisors to principal divisors and multiply,
respectively preserve, degrees; they descend to $\Pic$ and $\Pic^0$, and the divisor identity $\varphi_*\varphi^*=\deg\varphi$ descends.

(vi) For $f$ nonconstant, $f^*\bar k(\PP^1)=f^*\bar k(t)=\bar k(f)$, so the degree is $[\bar k(C):\bar k(f)]$. The coordinate $t$ is a uniformiser at $0\in\PP^1$, so
$e_f(P)=\ord_P(f^*t)=\ord_Pf$ at a zero $P$ of $f$; and $1/t$ is a uniformiser at $\infty$, so $e_f(P)=-\ord_Pf$ at a pole. Hence $f^*((0))=\operatorname{div}_0(f)$ and
$f^*((\infty))=\operatorname{div}_\infty(f)$, both of degree $\deg f$ by \Cref{thm:curves:degree-formula}. So $\deg\operatorname{div}f=0$, and a nonconstant function has a pole
divisor of degree $\deg f\ge1$. This proves \Cref{prop:curves:function-field}(iii) and \Cref{prop:curves:principal}(ii); neither was used in the proof of
\Cref{thm:curves:degree-formula}, which rests on the sketched facts named there.

(ii) For separable $\varphi$ only finitely many points are ramified (\Cref{lem:curves:local-differential}), and over every other $Q$ all $e_\varphi(P)=1$, so
$\#\varphi^{-1}(Q)=\deg\varphi$ by \Cref{thm:curves:degree-formula}. In general $\varphi$ factors as a purely inseparable morphism, which is bijective on points, followed
by a separable one of degree $\deg_s\varphi$ (\Cref{prop:curves:separability-frobenius}(b), stated).

(v) The first statement is the tower law for degrees of field extensions. If $\deg\varphi=1$, $\varphi^*$ is an isomorphism of function fields and $\varphi$ is an
isomorphism (\Cref{thm:curves:equivalence}); an isomorphism induces an isomorphism of function fields. (iv) is proved where cited.
\end{proof}

\begin{table}[ht]
\centering
\small
\begin{tabular}{@{}p{0.27\textwidth}p{0.36\textwidth}p{0.27\textwidth}@{}}
\toprule
description of $\deg\varphi$ & statement & where \\
\midrule
field degree & $[\bar k(C_1):\varphi^*\bar k(C_2)]$ & \Cref{def:curves:degree} \\
counting preimages & $\sum_{P\mapsto Q}e_\varphi(P)$ for every $Q$; $\#\varphi^{-1}(Q)$ for general $Q$ ($\varphi$ separable) & \Cref{thm:curves:degree-formula}, \Cref{cor:curves:degree-three-ways}(ii) \\
divisors & $\deg\varphi^*D=\deg\varphi\deg D$; $\varphi_*\varphi^*=\deg\varphi$ & \Cref{cor:curves:degree-three-ways}(iii) \\
analytic (over $\C$) & local multiplicities summed over a fibre & \Cref{thm:curves:holomorphic-maps}(c) \\
cohomology (over $\C$) & $\varphi^*$ on $H^2=\Z$; $\langle\varphi^*x,\varphi^*y\rangle=\deg\varphi\,\langle x,y\rangle$ & \Cref{prop:curves:transfer}, \Cref{prop:curves:degree-cup-product} \\
\bottomrule
\end{tabular}
\caption{The degree of a nonconstant morphism of smooth projective curves.}\label{tab:curves:degree}
\end{table}

\begin{proposition}[separability and the Frobenius morphism]\label{prop:curves:separability-frobenius}
\leavevmode
\begin{enumerate}[label=(\alph*)]
\item In characteristic $0$ every nonconstant morphism of curves is separable.
\item In characteristic $p$, $\varphi$ is separable iff $\varphi^*\omega\ne0$ for one (equivalently every) nonzero differential $\omega$ on $C_2$, where
$\varphi^*(f\,dg)=(\varphi^*f)\,d(\varphi^*g)$. For a power $r$ of $p$ and a curve $C\subseteq\PP^n$, let $C^{[r]}$ (not to be confused with the twist notation $E^{(q)}$ of Appendix B) be the curve whose ideal is generated by the polynomials obtained from generators of $I(C)$ by raising
their coefficients to the $r$-th power, and $\pi_r\colon C\to C^{[r]}$, $[x_0:\dots:x_n]\mapsto[x_0^r:\dots:x_n^r]$, the \emph{$r$-power Frobenius morphism} (it maps $C$ into $C^{[r]}$
because $F^{[r]}(x_0^r,\dots,x_n^r)=F(x_0,\dots,x_n)^r$, where $F^{[r]}$ has the coefficients of $F$ raised to the $r$-th power). Every nonconstant $\varphi$ factors as $C_1\xrightarrow{\pi_r}C_1^{[r]}\xrightarrow{\lambda}C_2$ with $r=\deg_i\varphi$ and $\lambda$
separable (stated). For $C$ defined over $\F_q$, $C^{[q]}=C$ and $\pi_q$ is the morphism of the next item.
\item Let $k=\F_q$ and $C\subseteq\PP^n$ a curve defined over $\F_q$. The \emph{$q$-power Frobenius morphism} $\pi_q\colon[x_0:\dots:x_n]\mapsto[x_0^q:\dots:x_n^q]$ maps $C$ to
itself and is a bijection on $C(\bar\F_q)$; $\pi_q^*\bar\F_q(C)=\bar\F_q(C)^q$; $d(\pi_q^*f)=0$ for every $f$; $\pi_q$ is purely inseparable of degree $q$ and
$e_{\pi_q}(P)=q$ at every point; its fixed points are $C(\F_q)$; and on $C(\bar\F_q)$ it coincides with the action of the arithmetic Frobenius $\Frob_q\in G_{\F_q}$.
\end{enumerate}
\end{proposition}

\begin{proof}
(a) Fields of characteristic $0$ are perfect, and every algebraic extension of a perfect field is separable (\Cref{ch:galois}); here
$\varphi^*\bar k(C_2)$ has characteristic $0$. (b) is \cite[Chapter II, Proposition 4.2(c) and Corollary 2.12]{SilvermanAEC}; we do not prove it.

(c) If $F$ is a homogeneous polynomial with coefficients in $\F_q$, then $F(x_0^q,\dots,x_n^q)$ equals $F(x_0,\dots,x_n)^q$ (the freshman's dream, \Cref{ch:rings},
and $a^q=a$ for $a\in\F_q$); so $\pi_q$ maps $C$ into $C$. On $\bar\F_q$ the map $a\mapsto a^q$ is a bijection, so $\pi_q$ is a bijection of $\PP^n(\bar\F_q)$ preserving
$C(\bar\F_q)$; in particular it is nonconstant. For $f=F/G\in\bar\F_q(C)$ the pullback $\pi_q^*f=F(x^q)/G(x^q)$ equals $(F'/G')^q$, where the coefficients of $F'$, $G'$ are the $q$-th
roots of those of $F$, $G$ ($\bar\F_q$ is perfect); conversely every $(F'/G')^q$ is of this form. So $\pi_q^*\bar\F_q(C)=\bar\F_q(C)^q$, and $d(h^q)=qh^{q-1}dh=0$. The degree
$[\bar\F_q(C):\bar\F_q(C)^q]$ is $q$: for a separating transcendental element $t$ (one with $\bar\F_q(C)/\bar\F_q(t)$ separable, which exists, \cite[Chapter II,
Section 1]{SilvermanAEC}), $\bar\F_q(C)=\bar\F_q(C)^q(t)$ and $t$ has minimal polynomial $X^q-t^q$ over $\bar\F_q(C)^q$ (sketched; \cite[Chapter II, Proposition
2.11]{SilvermanAEC}). Every point has exactly one preimage, so $e_{\pi_q}(P)=q$ by \Cref{thm:curves:degree-formula}. A point $[x_0:\dots:x_n]$, scaled so that some
$x_i=1$, is fixed iff $x_j^q=x_j$ for all $j$, that is, iff all $x_j\in\F_q$. Finally $\Frob_q$ acts on coordinates by $a\mapsto a^q$ (\Cref{conv:curves:conventions}(C7)).
\end{proof}

On $\PP^1$ over $\F_q$ the sketched degree computation is immediate: $\pi_q^*\bar\F_q(t)=\bar\F_q(t^q)$, and $t$ is a root of $X^q-t^q=(X-t)^q$, irreducible over
$\bar\F_q(t^q)$ by Eisenstein, so $\deg\pi_q=q$, while $d(t^q)=0$ shows that the extension is purely inseparable. Two features of (c) matter later. First, the
morphism $\pi_q$ and the Galois element $\Frob_q$ have the same effect on points but act differently on functions: $\pi_q^*f=f\circ\pi_q$ is
$\bar\F_q$-linear in $f$, whereas $\Frob_q$ acts on $\bar\F_q(C)$ through the coefficients (\Cref{prop:curves:galois-action}(ii)), semilinearly. Second, because
they agree on points, the arithmetic Frobenius is what $\pi_q$ induces on torsion points and on $T_\ell J$, and its inverse, the geometric Frobenius, is what
appears on $H^1$ (\Cref{thm:curves:weil-pairing-etale}(d), \Cref{conv:curves:duality}(e)). Since $H^1$ is the dual of $T_\ell J$, at a prime $p$ of good reduction
$\det(1-\Frob_p^{-1}X\mid H^1)=\det(1-\Frob_pX\mid T_\ell J)$; using $\Frob_p$ itself on $H^1$ would replace each eigenvalue $\alpha$ by $\alpha^{-1}$. This is
the convention behind the book's form $\det(1-\Frob_p^{-1}p^{-s}\mid H^1)$ of the $L$-factors.

\begin{paperlink}\label{pl:curves:modular-degree}
Section 1 of the paper: ``the $X_0(N)$-optimal curve, for which the modular parametrisation $\varphi_E\colon X_0(N)\to E$ has minimal degree'' and
``$\degphi$, its modular degree''. The degree is that of \Cref{def:curves:degree}, and by \Cref{cor:curves:degree-three-ways} and
\Cref{tab:curves:degree} it is at once $[\Q(X_0(N)):\varphi_E^*\Q(E)]$ (degrees of field extensions do not change under $\otimes_\Q\Qbar$ for function fields
of curves defined over $\Q$), the number of points of $X_0(N)(\C)$ over a general point of $E(\C)$, and the factor $\deg\varphi_E$ in
$\varphi_{E*}\varphi_E^*=[\deg\varphi_E]$ on $E=\Pic^0(E)$. The proof of Proposition prop:pullback uses the cohomological form, ``$\varphi^*$ multiplies the fundamental
class of $H^2$ by $\degphi$'', which is \Cref{prop:curves:transfer}(iv). PARI's \texttt{ellmoddegree} computes this number (\Cref{ch:computation}).
\end{paperlink}

\section{Ramification and the Riemann-Hurwitz formula}\label{sec:curves:hurwitz}

The Riemann-Hurwitz formula relates the genera of two curves joined by a morphism. In the book it gives the genus of $X_0(N)$ (\Cref{ch:modforms}), shows that
isogenies are unramified (\Cref{ch:ellcurves}), and checks the modular degrees quoted in the paper against the genera of the modular curves
(\Cref{cor:curves:rh-consequences}(e)).

\begin{definition}[ramification]\label{def:curves:ramification}
Let $\varphi\colon C_1\to C_2$ be nonconstant. $\varphi$ is \emph{unramified} at $P$ if $e_\varphi(P)=1$ and \emph{ramified} otherwise, and $\varphi(P)$ is then a \emph{branch
point}. The ramification at $P$ is \emph{tame} if $\operatorname{char}k\nmid e_\varphi(P)$ and \emph{wild} otherwise. $\varphi$ is \emph{unramified} if it is unramified at every point.
\end{definition}

\begin{lemma}[the local computation]\label{lem:curves:local-differential}
Let $\varphi\colon C_1\to C_2$ be nonconstant and separable, $\omega\ne0$ a differential on $C_2$, $P\in C_1$, $Q=\varphi(P)$, $e=e_\varphi(P)$, and $t_P$, $t_Q$ uniformisers. Then
$\varphi^*\omega\ne0$ and
\[
\ord_P(\varphi^*\omega)=e\,\ord_Q(\omega)+\ord_P\Bigl(\frac{d(\varphi^*t_Q)}{dt_P}\Bigr),\qquad \ord_P\Bigl(\frac{d(\varphi^*t_Q)}{dt_P}\Bigr)\ge e-1,
\]
with equality iff $\operatorname{char}k\nmid e$. In particular only finitely many points of $C_1$ are ramified.
\end{lemma}

\begin{proof}
$\varphi^*\omega\ne0$ by \Cref{prop:curves:separability-frobenius}(b) (in characteristic $0$ by \Cref{prop:curves:differentials}(i): $d$ of a separating element is
nonzero). Write $\omega=g\,dt_Q$ and $\varphi^*t_Q=ut_P^e$ with $u$ a unit at $P$. Then $\varphi^*\omega=(\varphi^*g)\,d(ut_P^e)$ and
\[
\frac{d(ut_P^e)}{dt_P}=eut_P^{e-1}+t_P^e\frac{du}{dt_P},
\]
where $du/dt_P$ is regular at $P$ (\Cref{prop:curves:differentials}(ii)). If $\operatorname{char}k\nmid e$, the first term has order exactly $e-1$ and the second order at
least $e$, so the sum has order $e-1$. If $\operatorname{char}k\mid e$, the first term vanishes and the order is at least $e$. Since
$\ord_P\varphi^*g=e\,\ord_Qg$ (\Cref{cor:curves:degree-three-ways}(iii)), the formula follows. At a ramified $P$ the second term is $\ge e-1\ge1$; but it equals
$\ord_P\varphi^*\omega-e\,\ord_Q\omega$, which is nonzero only at the finitely many zeros and poles of $\varphi^*\omega$ and points over zeros and poles of $\omega$
(\Cref{prop:curves:differentials}(ii), \Cref{thm:curves:constant-or-surjective}: fibres are finite by \Cref{thm:curves:degree-formula}).
\end{proof}

\begin{theorem}[Riemann-Hurwitz]\label{thm:curves:riemann-hurwitz}
Let $\varphi\colon C_1\to C_2$ be a nonconstant separable morphism of smooth projective curves of genera $g_1$, $g_2$. Then
\[
2g_1-2\ \ge\ \deg\varphi\,(2g_2-2)+\sum_{P\in C_1}\bigl(e_\varphi(P)-1\bigr),
\]
with equality iff every ramification is tame; in particular with equality in characteristic $0$.
\end{theorem}

\begin{proof}
Take $\omega\ne0$ on $C_2$ and sum the formula of \Cref{lem:curves:local-differential} over $P\in C_1$. The first terms give
$\sum_Q\ord_Q(\omega)\sum_{P\mapsto Q}e_\varphi(P)=\deg\varphi\cdot\deg\operatorname{div}(\omega)$ by \Cref{thm:curves:degree-formula}, that is $\deg\varphi\,(2g_2-2)$ by
\Cref{cor:curves:rr-consequences}(b); the left side is $\deg\operatorname{div}(\varphi^*\omega)=2g_1-2$. The second terms are $\ge\sum(e_\varphi(P)-1)$, with equality iff equality
holds at every $P$, that is, iff all ramification is tame.
\end{proof}

A topological proof over $\C$ is sketched in \Cref{thm:curves:topology}(c): triangulate $C_2$ with the branch points among the vertices, lift the triangulation and
count $V-E+F$.

For $z\mapsto z^n$ on $\PP^1$ with $\operatorname{char}k\nmid n$ (the example after \Cref{thm:curves:degree-formula}) the formula reads $-2=n(-2)+(n-1)+(n-1)$,
with the two ramification points $0$ and $\infty$. Wild ramification is genuinely different. In characteristic $p$ let $\varphi\colon\PP^1\to\PP^1$ be
$y\mapsto x=y^p-y$ (the Artin-Schreier map). It has degree $p$ ($y$ is a root of $T^p-T-x$, irreducible over $\bar k(x)$ because the map is totally ramified at
$\infty$, as we now see) and is separable, because $d(y^p-y)=-dy\ne0$ (\Cref{prop:curves:separability-frobenius}(b)). At finite points $dx/dy=-1$ is a unit, so
$\varphi$ is unramified there. At $\infty$, with uniformisers $t_P=1/y$ and $t_Q=1/x$,
\[
\varphi^*t_Q=\frac1{y^p-y}=\frac{t_P^p}{1-t_P^{p-1}},\qquad\frac{d(\varphi^*t_Q)}{dt_P}=\frac{pt_P^{p-1}}{1-t_P^{p-1}}+\frac{(p-1)t_P^{2p-2}}{(1-t_P^{p-1})^2}=\frac{-t_P^{2p-2}}{(1-t_P^{p-1})^2},
\]
so $e_\varphi(\infty)=p$, wild, and the order in \Cref{lem:curves:local-differential} is $2p-2$, strictly larger than $e-1=p-1$. With $2p-2$ the exact formula
of the proof gives $-2=p(-2)+(2p-2)$, which is correct, while the tame count $\sum(e_\varphi(P)-1)=p-1$ gives the strict inequality $-2>-2p+(p-1)$ of the theorem.
(The degree claim: the fibre over $\infty$ is one point with $e=p$, so $\deg\varphi=p$ by \Cref{thm:curves:degree-formula}, and $T^p-T-x$, having $y$ as a root and
degree $p$, is the minimal polynomial.)

\begin{corollary}[consequences and examples]\label{cor:curves:rh-consequences}
Let $\varphi\colon C_1\to C_2$ be nonconstant and separable.
\begin{enumerate}[label=(\alph*)]
\item If $g_2\ge1$ then $g_1\ge g_2$. If $g_1=0$ then $g_2=0$ (Lueroth's theorem, for separable maps).
\item If $g_1=g_2=1$ then $\varphi$ is unramified, and $\#\varphi^{-1}(Q)=\deg\varphi$ for every $Q$. In characteristic $0$ this applies to every nonconstant morphism
between genus-one curves, in particular to every isogeny (\Cref{ch:ellcurves}).
\item The double cover $x\colon C\to\PP^1$ of the hyperelliptic curve $C$ given by $y^2=h(x)$, $h$ squarefree of degree $2g+1$ or $2g+2$,
$\operatorname{char}k\ne2$, is branched exactly at the $2g+2$ points consisting of the roots of $h$ and, if $\deg h$ is odd, $\infty$; so $g(C)=g$.
\item The projection $[X:Y:Z]\mapsto[X:Y]$ of the Fermat curve $X^d+Y^d+Z^d=0$, $\operatorname{char}k\nmid d$, has degree $d$ and is totally ramified ($e=d$) over the $d$
roots of $X^d+Y^d=0$ and unramified elsewhere; so $2g-2=-2d+d(d-1)$ and $g=\frac{(d-1)(d-2)}2$, as in \Cref{thm:curves:plane-genus}.
\item Over $\C$, for a nonconstant $\varphi\colon X\to E$ with $g(E)=1$: $\sum_P(e_\varphi(P)-1)=2g(X)-2$, whatever $\deg\varphi$. For the paper's curves (genera from
\Cref{ex:curves:genus-examples}(f), degrees from Cremona's tables, checked with PARI's \texttt{ellmoddegree}):
\begin{center}
\small
\begin{tabular}{@{}lccc@{}}
\toprule
$\varphi_E$ & $g(X_0(N))$ & $\deg\varphi_E$ & $\sum_P(e_\varphi(P)-1)$ \\
\midrule
$X_0(11)\to\curve{11a1}$ & $1$ & $1$ & $0$ \\
$X_0(20)\to\curve{20a1}$ & $1$ & $1$ & $0$ \\
$X_0(121)\to\curve{121a1}$ & $6$ & $6$ & $10$ \\
$X_0(539)\to\curve{539c1}$ & $49$ & $288=2^5\cdot3^2$ & $96$ \\
\bottomrule
\end{tabular}
\end{center}
A morphism of degree one is an isomorphism, so $X_0(11)\simeq\curve{11a1}$ and $X_0(20)\simeq\curve{20a1}$ over $\Q$: these modular curves are elliptic curves.
\item (Forward pointer.) The genus formula for $X_0(N)$ in \Cref{ch:modforms} is \Cref{thm:curves:riemann-hurwitz} for the morphism $X_0(N)\to X_0(1)\simeq\PP^1$, of degree
$[\SL_2(\Z):\Gamma_0(N)]$, ramified only over $j=0$, $1728$, $\infty$.
\end{enumerate}
\end{corollary}

\begin{proof}
(a) If $g_2\ge1$ the theorem gives $2g_1-2\ge\deg\varphi(2g_2-2)\ge2g_2-2$. If $g_1=0$ and $g_2\ge1$, the right side is $\ge0>-2$, a contradiction.
(b) $0\ge0+\sum(e_\varphi(P)-1)$, so every $e_\varphi(P)=1$, and \Cref{thm:curves:degree-formula} gives $\#\varphi^{-1}(Q)=\deg\varphi$ for every $Q$.

(c) $y\notin\bar k(x)$ because $h$ is not a square, so $[\bar k(C):\bar k(x)]=2$, and the map is separable as $\operatorname{char}k\ne2$. Over $a\in\bar k$ the points of $C$ are
$(a,\pm\sqrt{h(a)})$: two points, hence $e=1$ at both, if $h(a)\ne0$, and one point, hence $e=2$, if $h(a)=0$. Over $\infty$: the smooth projective model is
covered by the affine curve $y^2=h(x)$ and the affine curve $w^2=H(s)$, $H(s)=s^{2g+2}h(1/s)$, glued by $s=1/x$, $w=ys^{g+1}$ (\cite[Chapter II,
Exercise 2.14]{SilvermanAEC}, exercise number not verified); $H$ is squarefree, $H(0)$ is the leading coefficient of $h$ if $\deg h=2g+2$, and $H$ has a simple
zero at $0$ if $\deg h=2g+1$. So $\infty$ is a branch point iff $\deg h$ is odd. All ramification is tame, and the theorem gives $2g(C)-2=2(-2)+(2g+2)$.

(d) The map is defined everywhere ($X=Y=0$ forces $Z=0$). The fibre over $[a:b]$ is $\{[a:b:z]:z^d=-(a^d+b^d)\}$: $d$ distinct points if $a^d+b^d\ne0$, and the
single point $[a:b:0]$ otherwise. The extension $\bar k(C)/\bar k(X/Y)$ is generated by $Z/Y$, a root of $T^d+(X/Y)^d+1$, which is irreducible (Eisenstein at a
simple root of $s^d+1$) and separable ($\operatorname{char}k\nmid d$); so the degree is $d$, and by \Cref{thm:curves:degree-formula} the single point over a root of $X^d+Y^d$
has $e=d$. All ramification is tame. (e) is the theorem with $g_2=1$; the table uses $2g-2=0,0,10,96$.
\end{proof}

\section[Compact Riemann surfaces and complex curves]{Compact Riemann surfaces and smooth projective complex curves}\label{sec:curves:riemann-surfaces}

$X_0(N)$ is built analytically in \Cref{ch:modforms}, as a compact Riemann surface, and used algebraically in \Cref{ch:moddeg} and in the paper, as a curve
over $\Q$; the modular parametrisation $\varphi_E$ is at once a holomorphic map and a morphism. This section introduces compact Riemann surfaces, proves the local
normal form of holomorphic maps (which gives the analytic degree), and states the dictionary between the two languages (\Cref{thm:curves:riemann-existence}),
summarised in \Cref{tab:curves:dictionary}. The prerequisites are those of \Cref{ch:prereq}: holomorphic functions, zeros of finite order, the inverse function
theorem, compactness and connectedness.

\begin{definition}[Riemann surfaces]\label{def:curves:riemann-surface}
A \emph{Riemann surface} is a connected Hausdorff topological space $X$ with a countable basis, together with a maximal atlas of \emph{charts}
$z\colon U\to z(U)\subseteq\C$ (homeomorphisms from open sets of $X$ onto open subsets of $\C$, the $U$ covering $X$) whose transition maps $z_2\circ z_1^{-1}$ are holomorphic
where defined. A function $f\colon U\to\C$ on an open subset is \emph{holomorphic} (\emph{meromorphic}) if $f\circ z^{-1}$ is, for every chart $z$; $\mathcal M(X)$ is the field of
meromorphic functions on $X$ (for $X$ connected, \Cref{ch:prereq}). A continuous map $F\colon X\to Y$ of Riemann surfaces is \emph{holomorphic} if $w\circ F\circ z^{-1}$ is
holomorphic for all charts $z$ of $X$ and $w$ of $Y$.
\end{definition}

\begin{example}[Riemann surfaces]\label{ex:curves:riemann-surfaces}
\leavevmode
\begin{enumerate}[label=(\alph*)]
\item $\C$, and every connected open subset of $\C$, with the identity chart.
\item The Riemann sphere $\PP^1(\C)=\C\cup\{\infty\}$, with the charts $z$ on $\C$ and $w=1/z$ on $\PP^1\setminus\{0\}$; the transition map $w=1/z$ is holomorphic on
$\C^\times$. It is compact.
\item $\C/\Lambda$ for a lattice $\Lambda=\Z\omega_1+\Z\omega_2$ ($\omega_1$, $\omega_2$ linearly independent over $\R$): the charts are local inverses of the projection
$\C\to\C/\Lambda$ on discs of radius less than half the shortest nonzero vector of $\Lambda$; the transition maps are translations $z\mapsto z+\lambda$. It is compact
(the image of a closed parallelogram) and, topologically, a torus. \Cref{ch:ellcurves} proves that every elliptic curve over $\C$ is of this form.
\item The smooth points of an affine plane curve $f(x,y)=0$ over $\C$: near a point with $f_y\ne0$, the holomorphic implicit function theorem writes the curve
as a graph $y=y(x)$, and $x$ is a chart; near a point with $f_x\ne0$, $y$ is a chart; the transition maps are holomorphic.
\item Hence $C(\C)$ is a compact Riemann surface for every smooth projective curve $C$ over $\C$: it is covered by the affine charts, it is closed in the compact
space $\PP^n(\C)$, and it is connected because $C$ is irreducible (stated; \cite[Appendix B]{Hartshorne}, \cite[Chapter VII]{Miranda}, chapter number not verified).
\item (Forward pointer, \Cref{ch:modforms}.) $\Gamma_0(N)\backslash\HH$, compactified by adding the finitely many cusps, with special charts at the cusps and at the
elliptic points, is the compact Riemann surface $X_0(N)(\C)$.
\end{enumerate}
\end{example}

\begin{theorem}[holomorphic maps of compact Riemann surfaces]\label{thm:curves:holomorphic-maps}
Let $F\colon X\to Y$ be a nonconstant holomorphic map of compact Riemann surfaces.
\begin{enumerate}[label=(\alph*)]
\item (Local normal form.) For each $P\in X$ there are charts $z$ at $P$ and $w$ at $F(P)$, centred at $0$, in which $F$ is $w=z^e$. The integer
$e=\operatorname{mult}_P(F)\ge1$ does not depend on the charts.
\item $F$ is open and surjective, and its fibres are finite.
\item The number $d=\sum_{P\in F^{-1}(Q)}\operatorname{mult}_P(F)$ does not depend on $Q\in Y$; it is the \emph{degree} of $F$.
\item $\operatorname{mult}_P(F)=1$ except at finitely many points of $X$; so $\#F^{-1}(Q)=d$ except for finitely many $Q$.
\item $F$ preserves orientation: in local coordinates its real Jacobian determinant is $\abs{h'(z)}^2\ge0$.
\end{enumerate}
\end{theorem}

\begin{proof}
(a) Take any charts $z$ at $P$ and $w$ at $F(P)$, centred at $0$, and put $h=w\circ F\circ z^{-1}$, holomorphic near $0$ with $h(0)=0$. $h$ is not identically zero near $0$:
otherwise, by the identity theorem and connectedness of $X$, $F$ would be constant. So $h(z)=z^eg(z)$ with $e\ge1$ and $g(0)\ne0$. On a small disc $g$ has a
holomorphic $e$-th root $g_1$ (\Cref{ch:prereq}: a nonvanishing function on a disc has a holomorphic logarithm), and $\zeta=zg_1(z)$ has $\zeta'(0)=g_1(0)\ne0$, so it is a
chart near $P$ by the inverse function theorem. In the charts $\zeta$, $w$ the map is $w=\zeta^e$. The exponent $e$ is the order of vanishing at $0$ of $w\circ F$ in any
chart at $P$; changing the chart at $P$ or at $F(P)$ composes with a biholomorphism fixing $0$, whose derivative at $0$ is nonzero, and this does not change the order
of a zero.

(b) The map $\zeta\mapsto\zeta^e$ is open, so by (a) $F$ is open. Then $F(X)$ is open, and compact, hence closed; $Y$ is connected, so $F(X)=Y$. By (a), $P$ is the only
preimage of $F(P)$ near $P$; so fibres are discrete, and they are closed in the compact $X$, hence finite.

(c) Let $F^{-1}(Q)=\{P_1,\dots,P_r\}$, with multiplicities $e_i$ and normal-form charts on pairwise disjoint neighbourhoods $U_i$ of $P_i$. There is a
neighbourhood $V$ of $Q$ with $F^{-1}(V)\subseteq\bigcup U_i$: otherwise there are $x_n\notin\bigcup U_i$ with $F(x_n)\to Q$; a subsequence converges (compactness) to
some $x\notin\bigcup U_i$ (the complement is closed) with $F(x)=Q$, which is absurd. Shrinking $V$, every $Q'\in V\setminus\{Q\}$ has exactly $e_i$ preimages in $U_i$
(the $e_i$-th roots of $w(Q')\ne0$), each of multiplicity one since $\frac d{d\zeta}\zeta^{e_i}\ne0$ for $\zeta\ne0$. So $Q'\mapsto\sum_{F(P)=Q'}\operatorname{mult}_P(F)$ is constant on
$V$; it is locally constant, hence constant on the connected $Y$.

(d) In the normal form at $P$, every point near $P$ other than $P$ has multiplicity one. So the set of points of multiplicity $\ge2$ is discrete; it is also closed
(its complement is open by the same remark), hence finite since $X$ is compact. Over $Q$ outside its finite image, $\#F^{-1}(Q)=d$ by (c).

(e) For $h=u+iv$ holomorphic, $\det\begin{psmallmatrix}u_x&u_y\\v_x&v_y\end{psmallmatrix}=u_x^2+v_x^2=\abs{h'}^2$ by the Cauchy-Riemann equations.
\end{proof}

The proofs of (a) and the openness in (b) use only that the target is a Riemann surface, not that it is compact. Hence \emph{a holomorphic function
$f\colon X\to\C$ on a compact Riemann surface is constant}: otherwise $f(X)$ would be open in $\C$, and also compact, hence closed and bounded, which is impossible for a
nonempty open subset of $\C$. (The same follows from the maximum modulus principle, \Cref{ch:prereq}.)

\begin{theorem}[Riemann existence and GAGA for curves; stated]\label{thm:curves:riemann-existence}
\leavevmode
\begin{enumerate}[label=(\alph*)]
\item (Riemann.) On a compact Riemann surface $X$ the meromorphic functions separate points and give local coordinates at every point, and $\mathcal M(X)$ is a function field of
one variable over $\C$.
\item (Equivalence.) $X\mapsto$ the smooth projective curve with function field $\mathcal M(X)$, and $C\mapsto C(\C)$, are inverse equivalences between compact Riemann
surfaces with nonconstant holomorphic maps and smooth projective curves over $\C$ with nonconstant morphisms. Under them: meromorphic functions are rational
functions; the order of a zero or pole is $\ord_P$; $\operatorname{mult}_P(F)=e_F(P)$; the analytic degree of \Cref{thm:curves:holomorphic-maps}(c) is $\deg F$ of
\Cref{def:curves:degree}; meromorphic and holomorphic $1$-forms (\Cref{def:curves:holomorphic-differential}) are the differentials and the regular differentials of
\Cref{def:curves:differentials}; divisors, linear equivalence and $\Pic$ agree.
\item The algebraic genus $g=l(K_C)$ equals the topological genus of \Cref{thm:curves:topology}.
\end{enumerate}
\end{theorem}

We do not prove this; see \cite[Chapters VI and VII]{Miranda} (Riemann existence and the algebraic curve of a Riemann surface), \cite[Appendix B]{Hartshorne}
(GAGA), and \cite[Sections 14 and 16]{Forster81} (all section and chapter numbers not verified). Once (b) is granted, $\operatorname{mult}_P(F)=e_F(P)$ is a tautology:
the multiplicity is the order of vanishing of $w\circ F$ at $P$ for a local coordinate $w$ at $F(P)$, which is $\ord_P(F^*t_{F(P)})$. In particular
$\deg\varphi_E$ means the same in the analytic and the algebraic sense.

\begin{table}[ht]
\centering
\small
\begin{tabular}{@{}p{0.42\textwidth}p{0.48\textwidth}@{}}
\toprule
smooth projective curve $C$ over $\C$ & compact Riemann surface $X=C(\C)$ \\
\midrule
function field $\C(C)$ & meromorphic functions $\mathcal M(X)$ \\
$\ord_Pf$ (valuation of $\bar k[C]_P$) & order of zero or pole in a chart \\
ramification index $e_\varphi(P)$ & multiplicity $\operatorname{mult}_P(\varphi)$, normal form $w=z^e$ \\
$\deg\varphi=[\C(C_1):\varphi^*\C(C_2)]$ & number of sheets; degree on $H^2$ \\
regular differentials & holomorphic $1$-forms $H^0(X,\Omega^1)$ \\
genus $g=l(K_C)$ & $\dim H^0(X,\Omega^1)$; number of handles; $\frac12\rank H_1(X,\Z)$ \\
$\Pic^0(C)$ & $H^0(X,\Omega^1)^\vee/H_1(X,\Z)$ (Abel-Jacobi) \\
\bottomrule
\end{tabular}
\caption{The dictionary of \Cref{thm:curves:riemann-existence,thm:curves:topology,thm:curves:genus-dimension,thm:curves:abel-jacobi}.}\label{tab:curves:dictionary}
\end{table}

\begin{theorem}[topology of compact Riemann surfaces]\label{thm:curves:topology}
\leavevmode
\begin{enumerate}[label=(\alph*)]
\item A compact Riemann surface is an orientable compact surface, homeomorphic to a sphere with $g$ handles for a unique $g\ge0$, its \emph{topological genus};
its Euler characteristic is $\chi(X)=2-2g$ (stated).
\item This $g$ is the genus of \Cref{thm:curves:riemann-roch} (stated; \Cref{thm:curves:riemann-existence}(c)).
\item (Topological Riemann-Hurwitz.) For $F\colon X\to Y$ holomorphic of degree $d$, $\chi(X)=d\,\chi(Y)-\sum_P(\operatorname{mult}_P(F)-1)$.
\end{enumerate}
\end{theorem}

\begin{proof}[Sketch of proof of (c)]
Triangulate $Y$ so that the finitely many branch points are vertices. Away from the branch points $F$ is a $d$-sheeted covering
(\Cref{thm:curves:holomorphic-maps}(c), (d)), so each open face and each open edge has exactly $d$ lifts, and these lifts triangulate $X$. A vertex $Q$ has
$\#F^{-1}(Q)=d-\sum_{F(P)=Q}(\operatorname{mult}_P(F)-1)$ lifts. Counting, $V_X-E_X+F_X=d(V_Y-E_Y+F_Y)-\sum_P(\operatorname{mult}_P(F)-1)$. For (a) and (b) see
\cite[Chapters I and II]{Miranda} and \cite[Section 17]{Forster81} (not verified). This is Euler's $V-E+F$ for a branched cover; with $\chi=2-2g$ it is
\Cref{thm:curves:riemann-hurwitz} over $\C$.
\end{proof}

For example, triangulate the sphere $Y=\PP^1(\C)$ as a double pyramid over a triangle, with apexes $0$ and $\infty$ and the equator $\abs w=1$ divided into three
edges: $V=5$, $E=9$, $F=6$, $\chi=2$. The map $z\mapsto z^2$ is branched only over $0$ and $\infty$, each with one preimage of multiplicity $2$. The three
equatorial vertices have two lifts each, every edge and face has two lifts, and the apexes have one lift each: $V=8$, $E=18$, $F=12$, so $\chi(X)=2=2\cdot2-(1+1)$,
in agreement with (c) and with $X\simeq\PP^1$.

\section{Holomorphic differentials, residues, periods}\label{sec:curves:differentials-analytic}

Holomorphic $1$-forms are the analytic face of regular differentials. In the book they are the weight-two cusp forms ($f\mapsto f(z)\,dz$ identifies
$S_2(\Gamma_0(N))$ with $H^0(X_0(N),\Omega^1)$, \Cref{ch:modforms}), and their periods give the period lattice $\Lambda_E$ and the Jacobian. In the paper, positivity
(\Cref{thm:curves:bilinear-relations}(b)) is the positivity of the Petersson norm $\langle f,f\rangle$ and of the area $\vol(\Lambda_E)$ in the formula of Zagier and
Watkins (Remark rem:degree; \Cref{ch:moddeg}), and $\varphi^*H^1(E,\C)$ in the proof of Proposition prop:pullback is spanned by $\varphi^*\omega_E$ and its conjugate
(\Cref{prop:curves:cohomology}(iii)).

\begin{definition}[$1$-forms on a Riemann surface]\label{def:curves:holomorphic-differential}
A \emph{meromorphic} (\emph{holomorphic}) $1$-form $\omega$ on $X$ is given in each chart $z$ by $f(z)\,dz$ with $f$ meromorphic (holomorphic), the expressions in two charts
being related by $f(z)\,dz=f(z(w))z'(w)\,dw$. $\ord_P\omega=\ord_Pf$ in a chart at $P$; $\operatorname{div}\omega=\sum\ord_P(\omega)(P)$; the \emph{residue} $\operatorname{res}_P\omega$ is the coefficient
of $z^{-1}$ in the Laurent expansion of $f$ in a chart centred at $P$. $H^0(X,\Omega^1)$ is the complex vector space of holomorphic $1$-forms. Smooth real or complex
$1$-forms $p\,dx+q\,dy$ and $2$-forms $r\,dx\wedge dy$, the exterior derivative $d$, the wedge product and $\int_X$ of $2$-forms (with the orientation of
\Cref{conv:curves:conventions}(C6)) are defined chartwise as on open subsets of $\R^2$, using a partition of unity for $\int_X$; Stokes' theorem $\int_M d\theta=\int_{\partial M}\theta$
holds for a compact surface $M$ with smooth boundary (\Cref{ch:prereq}; \cite[Chapter IV]{Miranda}). A $1$-form $\theta$ is \emph{closed} if $d\theta=0$, \emph{exact} if
$\theta=dg$ for a smooth function $g$; line integrals $\int_\gamma\theta$ along piecewise smooth paths are defined chartwise.
\end{definition}

The residue does not depend on the chart: $\operatorname{res}_P\omega=\frac1{2\pi i}\oint_\gamma\omega$ for a small positively oriented loop $\gamma$ around $P$ in a chart, by the
residue theorem of complex analysis, and the line integral is chart-independent. A holomorphic $1$-form is closed: $d(f\,dz)=\frac{\partial f}{\partial\bar z}\,d\bar z\wedge dz=0$.
Examples: $dz$ on $\C/\Lambda$ is a holomorphic $1$-form without zeros (the transition maps are translations, with derivative $1$); on $\PP^1$, $dz=-w^{-2}dw$ has a double
pole at $\infty$, and there is no nonzero holomorphic $1$-form, as in \Cref{prop:curves:differentials}(iv).

\begin{theorem}[the residue theorem on a compact Riemann surface]\label{thm:curves:residue}
For a meromorphic $1$-form $\omega$ on a compact Riemann surface $X$, $\sum_{P\in X}\operatorname{res}_P\omega=0$.
\end{theorem}

\begin{proof}
$\omega$ has finitely many poles $P_1,\dots,P_r$ (they are isolated in the compact $X$). Choose pairwise disjoint closed coordinate discs $\Delta_j$ around $P_j$ and put
$M=X\setminus\bigcup_j\operatorname{int}\Delta_j$, a compact surface with boundary $\partial M=\bigcup\partial\Delta_j$, each circle oriented clockwise as seen from its disc (the
boundary orientation of $M$). On $M$, $\omega$ is holomorphic, hence closed, so Stokes gives $0=\int_Md\omega=\int_{\partial M}\omega=-\sum_j\oint_{\partial\Delta_j}\omega=-2\pi i\sum_j\operatorname{res}_{P_j}\omega$,
the circles $\partial\Delta_j$ now oriented counterclockwise.
\end{proof}

On $\PP^1$ the form $dz/z$ has residue $1$ at $0$, and in the chart $w=1/z$ at $\infty$ it is $-w^{-2}dw/w^{-1}=-dw/w$, with residue $-1$: the sum is $0$. Compactness is
what makes the theorem work: on $\C$ alone, $dz/z$ has the single residue $1$. The point at infinity carries the missing residue, just as a closed surface
admits no net flux out of it.

\begin{corollary}[the argument principle]\label{cor:curves:argument-principle}
For $f\in\mathcal M(X)^\times$, $\deg\operatorname{div}(f)=\sum_P\operatorname{res}_P(df/f)=0$.
\end{corollary}

\begin{proof}
In a chart at $P$ write $f=z^mu$ with $m=\ord_Pf$ and $u$ holomorphic and nonzero near $0$; then $df/f=m\,dz/z+du/u$, and $du/u$ is holomorphic, so
$\operatorname{res}_P(df/f)=m$. Apply \Cref{thm:curves:residue}.
\end{proof}

This is a second proof of \Cref{prop:curves:principal}(ii) over $\C$, independent of \Cref{thm:curves:degree-formula}. In the language of physics: the total charge on a
closed surface is zero.

\begin{theorem}[the genus counts holomorphic forms]\label{thm:curves:genus-dimension}
For a compact Riemann surface $X$ of genus $g$, $\dim_\C H^0(X,\Omega^1)=g$, and $\deg\operatorname{div}\omega=2g-2$ for every meromorphic $1$-form $\omega\ne0$.
\end{theorem}

\begin{proof}
By \Cref{thm:curves:riemann-existence}(b), holomorphic $1$-forms are regular differentials and meromorphic ones are differentials, so this is
\Cref{cor:curves:rr-consequences}(a), (b) (proved there from the stated Riemann-Roch theorem).
\end{proof}

Examples: $H^0(\C/\Lambda,\Omega^1)=\C\,dz$ (a holomorphic form is $f\,dz$ with $f$ holomorphic on the compact $\C/\Lambda$, hence constant); $H^0(\PP^1,\Omega^1)=0$.
Forward pointer: $\dim S_2(\Gamma_0(N))=g(X_0(N))$ (\Cref{ch:modforms}), for instance $\dim S_2(\Gamma_0(121))=6$.

\begin{theorem}[Riemann's bilinear relations]\label{thm:curves:bilinear-relations}
Let $X$ be a compact Riemann surface of genus $g\ge1$ and $a_1,\dots,a_g,b_1,\dots,b_g$ a symplectic basis of $H_1(X,\Z)$ (\Cref{thm:curves:symplectic-basis}). For a closed
$1$-form $\theta$ put $A_j(\theta)=\int_{a_j}\theta$ and $B_j(\theta)=\int_{b_j}\theta$.
\begin{enumerate}[label=(\alph*)]
\item For $\omega,\eta\in H^0(X,\Omega^1)$, $\int_X\omega\wedge\eta=0$.
\item For $\omega\in H^0(X,\Omega^1)$, $\omega\ne0$: $\frac i2\int_X\omega\wedge\bar\omega>0$.
\item For closed $1$-forms $\theta$, $\eta$: $\displaystyle\int_X\theta\wedge\eta=\sum_{j=1}^g\bigl(A_j(\theta)B_j(\eta)-B_j(\theta)A_j(\eta)\bigr)$.
\item For $\omega,\eta\in H^0(X,\Omega^1)$, and for $\omega\ne0$ respectively,
\[
\sum_j\bigl(A_j(\omega)B_j(\eta)-B_j(\omega)A_j(\eta)\bigr)=0,\qquad\sum_j\operatorname{Im}\bigl(\overline{A_j(\omega)}B_j(\omega)\bigr)>0.
\]
\item A holomorphic $1$-form with all $A$-periods zero is zero. The period map $H_1(X,\Z)\to H^0(X,\Omega^1)^\vee$, $\gamma\mapsto(\omega\mapsto\int_\gamma\omega)$, is injective, and its image
$\Lambda_X$ is a lattice in the $g$-dimensional complex vector space $H^0(X,\Omega^1)^\vee$.
\end{enumerate}
\end{theorem}

\begin{proof}
(a) Locally $\omega\wedge\eta=fg\,dz\wedge dz=0$. (b) Locally $\omega=f\,dz$ and $\frac i2\omega\wedge\bar\omega=\frac i2\abs f^2dz\wedge d\bar z=\abs f^2dx\wedge dy$, since $dz\wedge d\bar z=-2i\,dx\wedge dy$;
the integral of this nonnegative density is positive unless $f\equiv0$. (In physical terms $\frac i2\int\omega\wedge\bar\omega=\int\abs f^2\,dx\,dy$ is a squared norm; the Petersson
product of \Cref{ch:modforms} is $\langle f,g\rangle=\frac i2\int_{X_0(N)}f\,dz\wedge\overline{g\,dz}$, \Cref{ch:notation}, B-C8.)
(d) The first identity is (c) with (a). For the second apply (c) to $\theta=\omega$, $\eta=\bar\omega$, using $A_j(\bar\omega)=\overline{A_j(\omega)}$:
$\frac i2\int\omega\wedge\bar\omega=\frac i2\sum_j\bigl(A_j\bar B_j-\overline{A_j\bar B_j}\bigr)=\frac i2\cdot2i\operatorname{Im}\sum_jA_j\bar B_j=\sum_j\operatorname{Im}(\bar A_jB_j)$, positive by (b).
(e) If all $A_j(\omega)=0$, the sum in (d) is zero, so $\omega=0$. If $\int_\gamma\omega=0$ for all $\omega$, then also $\int_\gamma\bar\omega=0$, so $\gamma$ pairs to zero with
$H^0(X,\Omega^1)\oplus\overline{H^0(X,\Omega^1)}$, which is all of $H^1(X,\C)$ (\Cref{prop:curves:cohomology}(iii)); as $H_1(X,\Z)$ is free and $H^1(X,\C)=\Hom(H_1(X,\Z),\C)$,
$\gamma=0$. The same argument with real coefficients shows that the $\R$-linear extension $H_1(X,\Z)\otimes\R\to H^0(X,\Omega^1)^\vee$ is injective; both sides have real dimension
$2g$ (\Cref{thm:curves:symplectic-basis}, \Cref{thm:curves:genus-dimension}), so it is an isomorphism, and the image of $H_1(X,\Z)$ is a lattice.
\end{proof}

\begin{proof}[Sketch of proof of (c)]
Cut $X$ along loops representing the symplectic basis, all through one point; the result is a $4g$-gon $\Pi$ whose sides are identified in the pattern
$a_1b_1a_1^{-1}b_1^{-1}\cdots a_gb_ga_g^{-1}b_g^{-1}$. On the simply connected $\Pi$, $\theta=dF$ for a smooth function $F$, and Stokes gives
$\int_X\theta\wedge\eta=\int_\Pi d(F\eta)=\int_{\partial\Pi}F\eta$, because $d(F\eta)=dF\wedge\eta+F\,d\eta=\theta\wedge\eta$ ($\eta$ is closed); here $\partial\Pi$ is traversed
counterclockwise, which for the torus of \Cref{ex:curves:torus} is the word $aba^{-1}b^{-1}$ with $a\cdot b=+1$. Fix $j$ and label the vertices of the block
$a_jb_ja_j^{-1}b_j^{-1}$ as $v_0\xrightarrow{a_j}v_1\xrightarrow{b_j}v_2\xrightarrow{a_j^{-1}}v_3\xrightarrow{b_j^{-1}}v_4$. A point $p(s)$ on the first side and the point $p'(s)$ on the
third side that is glued to it lie over the same point of $a_j$, the third side running from $v_3$ ($s=0$) to $v_2$ ($s=1$). The difference $F(p'(s))-F(p(s))$ is
constant in $s$: its derivative is $\theta$ evaluated on the same tangent vector of $a_j$, minus itself. At $s=1$ it is $F(v_2)-F(v_1)=\int_{v_1}^{v_2}dF=B_j(\theta)$. As
the third side is traversed against $a_j$, the two sides contribute
\[
\int_{a_j}\bigl(F(p)-F(p')\bigr)\eta=-B_j(\theta)A_j(\eta).
\]
In the same way a point $q(s)$ of the second side and the glued point $q'(s)$ of the fourth (running from $v_4$ at $s=0$ to $v_3$ at $s=1$) have
$F(q')-F(q)$ constant, equal at $s=1$ to $F(v_3)-F(v_2)=-A_j(\theta)$, and the two sides contribute $\int_{b_j}(F(q)-F(q'))\eta=A_j(\theta)B_j(\eta)$. Summing over
$j$ gives (c). On the torus, with $\theta=ds$, $\eta=dt$ and $F=s$, the two contributions are $0$ and $1$. What is not proved here is that cutting along the
loops of a symplectic basis gives such a polygon. See \cite[Chapter VIII]{Miranda}
and \cite[Chapter 2, Section 2]{GriffithsHarris78} (sections not verified). On the torus (\Cref{ex:curves:torus}) the formula is checked directly.
\end{proof}

In the language of physics: for $X=\C/\Lambda$ and $\omega=dz$, (b) is the area of a period parallelogram, $\frac i2\int dz\wedge d\bar z=\operatorname{Im}(\bar\omega_1\omega_2)$, and
the lattice $\Lambda_X$ of (e) is the crystal lattice of periods (\Cref{ex:curves:torus}).

\section{\texorpdfstring{$H_1$, $H^1$}{H1, H1} and the cup product}\label{sec:curves:cohomology}

This section is the main input of the chapter to Section 4.1 of the paper. The paper works with $H^1(X_0(M^\Sigma),\Zl)$ and states that the cup product is ``perfect
and alternating'' on it and that the Hecke operators are self-adjoint for the cup product composed with $w_{M^\Sigma}$; Proposition prop:pullback computes the cup
product on $\varphi_E^*H^1(E,\Zl)$. We define $H_1$ and $H^1$ of a compact Riemann surface without sheaves ($H^1:=\Hom(H_1,\Z)$, with de Rham forms to compute), prove
that the cup product is perfect and alternating with Gram matrix $\sm0{I_g}{-I_g}0$ in a symplectic basis, and set up pullback, pushforward and the transfer with the
projection formula, which is the cohomological meaning of the degree. The convention box \Cref{conv:curves:duality} says which object is dual to which. Throughout,
$X$ and $Y$ are compact Riemann surfaces of genera $g_X$, $g_Y$ (or $g$).

\begin{definition}[first homology]\label{def:curves:homology}
$H_1(X,\Z)$ is the abelianisation of the fundamental group $\pi_1(X)$; equivalently, singular homology in degree one. For the reader: an element is represented by a
finite formal sum of closed oriented loops, two sums being identified when their difference bounds a (sum of oriented) surface pieces; a loop's class depends only on
its free homotopy class. A holomorphic map $\varphi\colon X\to Y$ induces $\varphi_*\colon H_1(X,\Z)\to H_1(Y,\Z)$, $[\gamma]\mapsto[\varphi\circ\gamma]$. For an abelian group $A$,
$H_1(X,A)=H_1(X,\Z)\otimes A$; in particular $H_1(X,\Zl)=H_1(X,\Z)\otimes\Zl$ and $H_1(X,\Z/n)=H_1(X,\Z)/n$.
\end{definition}

\begin{theorem}[symplectic bases; stated]\label{thm:curves:symplectic-basis}
Let $X$ have genus $g$.
\begin{enumerate}[label=(\roman*)]
\item $H_1(X,\Z)$ is free of rank $2g$, with a \emph{symplectic basis} $a_1,\dots,a_g,b_1,\dots,b_g$: on the model ``sphere with $g$ handles'', $a_j$ goes around the $j$-th handle
and $b_j$ through it.
\item The intersection pairing $H_1(X,\Z)\times H_1(X,\Z)\to\Z$, which on transverse oriented loops is the signed count of crossings of \Cref{conv:curves:conventions}(C6), is
well defined, bilinear, alternating and perfect, with $a_i\cdot b_j=\delta_{ij}$ and $a_i\cdot a_j=b_i\cdot b_j=0$.
\item Every class in $H_1(X,\Z)$ is represented by a sum of loops avoiding any given finite set, and the kernel of $H_1(X\setminus S,\Z)\to H_1(X,\Z)$, for a finite set
$S$, is generated by small loops around the points of $S$.
\end{enumerate}
\end{theorem}

We do not prove this; see \cite[Chapter VIII, Section 1]{Miranda}, \cite[Section 21]{Forster81}, \cite[Chapter III, Section 1]{FarkasKra92} and, for (iii), the
Mayer-Vietoris sequence in \cite[Chapter 2]{Hatcher02} (all section numbers not verified). In the language of physics, the intersection number is a signed
crossing count, like a winding or linking number.

\begin{definition}[first cohomology]\label{def:curves:cohomology}
$H^1(X,\Z)=\Hom(H_1(X,\Z),\Z)$, and for an abelian group (or ring) $A$, $H^1(X,A)=\Hom(H_1(X,\Z),A)$. A holomorphic map $\varphi\colon X\to Y$ induces
$\varphi^*\colon H^1(Y,A)\to H^1(X,A)$, the transpose of $\varphi_*$: $(\varphi^*\xi)(\gamma)=\xi(\varphi_*\gamma)$. A closed $1$-form $\theta$ defines the class $[\theta]\in H^1(X,\C)$,
$\gamma\mapsto\int_\gamma\theta$ (well defined by Stokes: the integral over a boundary is zero). $H^2(X,\R)$ is the space of smooth $2$-forms modulo exact ones, and $H^2(X,\Z)$
the subgroup of classes with integral integral.
\end{definition}

\begin{proposition}[cohomology of a compact Riemann surface]\label{prop:curves:cohomology}
\leavevmode
\begin{enumerate}[label=(\roman*)]
\item $H^1(X,\Z)$ is free of rank $2g$, and $H^1(X,A)=H^1(X,\Z)\otimes A$ for every abelian group $A$; in particular $H^1(X,\Zl)=H^1(X,\Z)\otimes\Zl$ is free of rank $2g$ over
$\Zl$. With the dual basis of a symplectic basis, written $a_1^\vee,\dots,a_g^\vee,b_1^\vee,\dots,b_g^\vee$, $H^1(X,\Z)=\bigoplus\Z a_j^\vee\oplus\bigoplus\Z b_j^\vee$.
\item (de Rham; stated.) Every class in $H^1(X,\R)$ is $[\theta]$ for a closed real $1$-form $\theta$, unique up to an exact form; $H^1(X,\Z)$ is the lattice of classes with
integral periods. $\int_X\colon H^2(X,\R)\to\R$ is an isomorphism, so $H^2(X,\Z)\simeq\Z$, generated by the class of integral $1$; this fixes the isomorphism $H^2(X,\Z)=\Z$.
\item (Hodge.) $H^1(X,\C)=H^0(X,\Omega^1)\oplus\overline{H^0(X,\Omega^1)}$: every class is $[\omega+\bar\eta]$ for unique holomorphic $\omega$, $\eta$.
\item (Stated.) For $g\ge1$ the universal cover of $X$ is $\C$ or the unit disc, so $X$ is a $K(\pi_1(X),1)$, and $H^1(X,M)=H^1(\pi_1(X),M)$ in the sense of \Cref{ch:cohomology}
for every abelian group $M$ with trivial action, compatibly with cup products. This is the bridge to the group cohomology of $\Gamma_0(N)$ used by Diamond, Flach and Guo
(\Cref{ch:moddeg}).
\end{enumerate}
\end{proposition}

\begin{proof}
(i) $H_1(X,\Z)\simeq\Z^{2g}$ (\Cref{thm:curves:symplectic-basis}), and $\Hom(\Z^{2g},A)=A^{2g}=\Hom(\Z^{2g},\Z)\otimes A$. The universal coefficient theorem, with $H_0$ free,
identifies this with singular cohomology (stated, \cite[Section 3.1]{Hatcher02}). (ii) is \cite[Chapter VIII]{Miranda} and \cite[Sections 10 and 19]{Forster81}
(not verified); we do not prove it.

(iii) (Sketch, given (ii) and \Cref{thm:curves:genus-dimension}.) The map $(\omega,\eta)\mapsto[\omega+\bar\eta]$ from $H^0(X,\Omega^1)\oplus H^0(X,\Omega^1)$ to $H^1(X,\C)$ is
injective: if $\omega+\bar\eta=dF$, then $\int_X(\omega+\bar\eta)\wedge\bar\omega=\int_Xd(F\bar\omega)=0$ by Stokes ($\bar\omega$ is closed), while
$\int_X\bar\eta\wedge\bar\omega=\overline{\int_X\eta\wedge\omega}=0$ by \Cref{thm:curves:bilinear-relations}(a); so $\int_X\omega\wedge\bar\omega=0$ and $\omega=0$ by
\Cref{thm:curves:bilinear-relations}(b); then $\bar\eta=dF$ and the same argument gives $\eta=0$. Both sides have complex dimension $2g$, by
\Cref{thm:curves:genus-dimension} and (i) with (ii). (The proof is complete given the stated inputs; it is labelled a sketch only because de Rham's theorem is.)
(iv) We do not prove this; see \cite[Section 1.B]{Hatcher02} (not verified).
\end{proof}

\begin{theorem}[the cup product]\label{thm:curves:cup-product}
The \emph{cup product}
\[
\langle\ ,\ \rangle\colon H^1(X,\Z)\times H^1(X,\Z)\to H^2(X,\Z)=\Z,\qquad\langle\theta,\eta\rangle=\int_X\theta\wedge\eta
\]
(de Rham representatives) is well defined, integral, bilinear, alternating and perfect. In the basis dual to a symplectic basis,
\[
\langle a_i^\vee,b_j^\vee\rangle=\delta_{ij},\qquad\langle a_i^\vee,a_j^\vee\rangle=\langle b_i^\vee,b_j^\vee\rangle=0,
\]
so the Gram matrix is $\sm0{I_g}{-I_g}0$, of determinant $1$. The map $\operatorname{PD}\colon H^1(X,\Z)\to H_1(X,\Z)$ defined by $\xi(\operatorname{PD}(\theta))=\langle\theta,\xi\rangle$ for all
$\xi$ is an isomorphism (\emph{Poincare duality}) with $\operatorname{PD}(a_j^\vee)=b_j$ and $\operatorname{PD}(b_j^\vee)=-a_j$, and an isometry:
$\langle\theta,\eta\rangle=\operatorname{PD}(\theta)\cdot\operatorname{PD}(\eta)$. For every commutative ring $R$ (in particular $\Zl$, $\Fl$, $\Ql$, $\C$) the $R$-bilinear extension to
$H^1(X,R)$ is perfect and alternating.
\end{theorem}

\begin{proof}
Changing $\theta$ by $dF$ changes $\int\theta\wedge\eta$ by $\int d(F\eta)=0$ (Stokes, $\eta$ closed), so the pairing is well defined; it is bilinear, and alternating because
$\theta\wedge\theta=0$. By \Cref{thm:curves:bilinear-relations}(c),
$\langle\theta,\eta\rangle=\sum_j(\theta(a_j)\eta(b_j)-\theta(b_j)\eta(a_j))$, an integer for integral classes; on the dual basis this gives the displayed values and the Gram matrix,
whose determinant is $1$. So the pairing is perfect (\Cref{ch:modules}: perfect iff the Gram determinant is a unit). For $\operatorname{PD}$: $\xi\mapsto\langle\theta,\xi\rangle$ is an
element of $\Hom(H^1(X,\Z),\Z)$, which is $H_1(X,\Z)$ (a free module of finite rank is its double dual); this defines $\operatorname{PD}$, an isomorphism by perfectness. The class
$b_j$ has $a_i^\vee(b_j)=0$ and $b_i^\vee(b_j)=\delta_{ij}$, which are the values $\langle a_j^\vee,a_i^\vee\rangle=0$ and $\langle a_j^\vee,b_i^\vee\rangle=\delta_{ij}$; so $\operatorname{PD}(a_j^\vee)=b_j$, and
likewise $\operatorname{PD}(b_j^\vee)=-a_j$. Then $\operatorname{PD}(a_i^\vee)\cdot\operatorname{PD}(b_j^\vee)=b_i\cdot(-a_j)=a_j\cdot b_i=\delta_{ij}=\langle a_i^\vee,b_j^\vee\rangle$, and the other basis pairs give $0$ on both
sides; by bilinearity $\operatorname{PD}$ is an isometry. Base change: the Gram matrix over $R$ is the image of the integral one, of determinant $1$ (\Cref{ch:modules}), and
$\langle x,x\rangle=0$ holds for $x\in H^1(X,R)$ because it holds on a basis and $\langle x,y\rangle=-\langle y,x\rangle$.
\end{proof}

The input that is only stated is the existence of a symplectic basis and the bilinear identity of \Cref{thm:curves:bilinear-relations}(c); perfectness is Poincare duality for
surfaces, which in general is \cite[Section 3.3]{Hatcher02}. Over $\R$ nondegeneracy also follows from positivity: by \Cref{prop:curves:cohomology}(iii), $\langle[\omega],[\bar\omega]\rangle\ne0$
for $\omega\ne0$.

\begin{example}[the torus]\label{ex:curves:torus}
Let $X=\C/\Lambda$, $\Lambda=\Z\omega_1+\Z\omega_2$ with $\operatorname{Im}(\omega_2/\omega_1)>0$, and real coordinates $z=s\omega_1+t\omega_2$, $s,t\in\R/\Z$. From
$dz=\omega_1ds+\omega_2dt$ one computes
\[
dz\wedge d\bar z=(\omega_1ds+\omega_2dt)\wedge(\bar\omega_1ds+\bar\omega_2dt)=(\omega_1\bar\omega_2-\omega_2\bar\omega_1)\,ds\wedge dt=-2i\operatorname{Im}(\bar\omega_1\omega_2)\,ds\wedge dt,
\]
so $dx\wedge dy=\frac i2dz\wedge d\bar z=\operatorname{Im}(\bar\omega_1\omega_2)\,ds\wedge dt$, and $\operatorname{Im}(\bar\omega_1\omega_2)=\abs{\omega_1}^2\operatorname{Im}(\omega_2/\omega_1)>0$ (as $\bar\omega_1\omega_2=\abs{\omega_1}^2\,\omega_2/\omega_1$);
so $ds\wedge dt$ is positive and $\int_Xds\wedge dt=1$ (the unit square in $(s,t)$ is a fundamental domain).
\begin{enumerate}[label=(\alph*)]
\item $H_1(X,\Z)=\Lambda$ (a path from $0$ to $\lambda$ in $\C$ projects to a loop), with $a$ the image of $[0,\omega_1]$ and $b$ that of $[0,\omega_2]$. They meet only at $0$, with
tangents $\partial_s$, $\partial_t$, a positive basis: $a\cdot b=+1$.
\item $H^1(X,\Z)$ has basis $[ds]$, $[dt]$, dual to $a$, $b$; $\langle ds,dt\rangle=\int ds\wedge dt=1$; $\operatorname{PD}(ds)=b$, $\operatorname{PD}(dt)=-a$; and
\Cref{thm:curves:bilinear-relations}(c) reads $1=A(ds)B(dt)-B(ds)A(dt)=1\cdot1-0\cdot0$.
\item The holomorphic form $dz=\omega_1ds+\omega_2dt$ has periods $A=\omega_1$, $B=\omega_2$, and $\frac i2\int dz\wedge d\bar z=\operatorname{Im}(\bar\omega_1\omega_2)$ is the area of the period
parallelogram: \Cref{thm:curves:bilinear-relations}(b), (d). For $\Lambda=\Z+\tau\Z$ the inequality of (d) is $\operatorname{Im}\tau>0$.
\item $[-1]\colon z\mapsto-z$ acts by $-1$ on $H_1$ and $H^1$, and is an isometry: $\langle-x,-y\rangle=\langle x,y\rangle$. Multiplication by an integer $m\ne0$ acts by $m$ on $H_1$ and on
$H^1$ (it pulls $ds$ back to $m\,ds$), and by $m^2$ on $H^2$ (it pulls $ds\wedge dt$ back to $m^2ds\wedge dt$), so it has degree $m^2$
(\Cref{prop:curves:transfer}(iv)); directly, the fibre over $0$ is $\frac1m\Lambda/\Lambda$, with $m^2$ points, all unramified.
\end{enumerate}
In the language of physics, $\Lambda$ is a crystal lattice, $H^1(X,\Z)=\Hom(\Lambda,\Z)$ its reciprocal lattice, and $ds\wedge dt$ the normalised volume form of the periodic
box.
\end{example}

\begin{proposition}[pullback, pushforward, transfer; the projection formula]\label{prop:curves:transfer}
Let $\varphi\colon X\to Y$ be a nonconstant holomorphic map of compact Riemann surfaces, $n=\deg\varphi$, $B\subseteq Y$ its finite set of branch points, $Y'=Y\setminus B$,
$X'=\varphi^{-1}(Y')$.
\begin{enumerate}[label=(\roman*)]
\item On de Rham classes $\varphi^*[\theta]=[\varphi^*\theta]$: $\int_\gamma\varphi^*\theta=\int_{\varphi_*\gamma}\theta$ (change of variables). $\varphi\colon X'\to Y'$ is an $n$-sheeted covering map,
locally the disjoint union of $n$ biholomorphisms, each orientation-preserving.
\item (Transfer.) There is a unique homomorphism $\varphi^!\colon H_1(Y,\Z)\to H_1(X,\Z)$ sending the class of a loop $\gamma$ in $Y'$ to the class of $\varphi^{-1}(\gamma)$, the sum of the
closed loops formed by the $n$ lifts of $\gamma$. It satisfies $\varphi_*\varphi^!=n$.
\item (Pushforward on $H^1$.) Define $\varphi_*\colon H^1(X,A)\to H^1(Y,A)$ as the transpose of $\varphi^!$: $(\varphi_*\eta)(\gamma)=\eta(\varphi^!\gamma)$. Every class in $H^1(X,\R)$ is
represented by a closed form $\theta$ vanishing near $\varphi^{-1}(B)$, and then $\varphi_*[\theta]$ is represented by the trace $\Tr\theta$: on $Y'$, $\Tr\theta=\sum_is_i^*\theta$ over the
local sections $s_1,\dots,s_n$ of $\varphi$, extended by $0$ near $B$.
\item (Projection formula.) For $\eta\in H^1(X,\Z)$ and $\xi,\xi'\in H^1(Y,\Z)$:
\[
\langle\varphi_*\eta,\xi\rangle_Y=\langle\eta,\varphi^*\xi\rangle_X,\qquad\varphi_*\varphi^*=n\text{ on }H^1(Y,\Z),\qquad\langle\varphi^*\xi,\varphi^*\xi'\rangle_X=n\,\langle\xi,\xi'\rangle_Y,
\]
and $\int_X\varphi^*\rho=n\int_Y\rho$ for every smooth $2$-form $\rho$ on $Y$: $\varphi^*$ is multiplication by $n$ on $H^2(\cdot,\Z)=\Z$. In particular $\varphi^*$ is injective on
$H^1(Y,\Z)$ and $\varphi_*\colon H_1(X,\Z)\to H_1(Y,\Z)$ has finite cokernel, killed by $n$.
\item $(\psi\circ\varphi)_*=\psi_*\varphi_*$, $(\psi\circ\varphi)^*=\varphi^*\psi^*$, $(\psi\circ\varphi)^!=\varphi^!\psi^!$; for an automorphism $w$, $w^!=(w^{-1})_*$ and $w_*=(w^{-1})^*$ on $H^1$.
\end{enumerate}
\end{proposition}

\begin{proof}
(i) The integral identity is the change of variables formula for line integrals (on each chart piece of $\gamma$). By \Cref{thm:curves:holomorphic-maps}(c), (d), every
point of $Y'$ has exactly $n$ preimages, each of multiplicity one, at which $\varphi$ is a local biholomorphism; the proof of \Cref{thm:curves:holomorphic-maps}(c) gives a
neighbourhood $V$ with $\varphi^{-1}(V)$ a disjoint union of $n$ open sets mapped biholomorphically onto $V$. Orientation is preserved by \Cref{thm:curves:holomorphic-maps}(e).

(ii) (Sketch.) By \Cref{thm:curves:symplectic-basis}(iii), classes in $H_1(Y,\Z)$ are represented by loops in $Y'$. For a covering, the preimage of a loop is a union of closed
loops covering it $n$ times in total, and the covering transfer $H_1(Y',\Z)\to H_1(X',\Z)$ is well defined (\cite[Section 3.G]{Hatcher02}, not verified). By
\Cref{thm:curves:symplectic-basis}(iii), the kernel of $H_1(Y',\Z)\to H_1(Y,\Z)$ is generated by small loops around the points of $B$; the preimage of such a loop is a sum of small loops
around the points of $\varphi^{-1}(B)$ (by the normal form $w=z^e$, a loop around $Q$ lifts to loops around the points over $Q$), which are null-homologous in $X$. So the transfer
descends to $\varphi^!$. Pushing forward the $n$ lifts of $\gamma$ gives $n\gamma$, so $\varphi_*\varphi^!=n$.

(iii) Let $\theta$ be a closed form on $X$ and $P\in\varphi^{-1}(B)$. On a coordinate disc $U$ around $P$, $\theta=dF$ for a smooth function $F$ (Poincare lemma on a disc).
With $\rho$ a smooth function equal to $1$ near $P$ and supported in $U$, $\theta-d(\rho F)$ is in the same class and vanishes near $P$. Doing this at the finitely many points of
$\varphi^{-1}(B)$ gives a representative vanishing near $\varphi^{-1}(B)$. For such $\theta$, $\Tr\theta$ is smooth and closed on $Y'$ (locally a finite sum of pullbacks of closed forms
by biholomorphisms) and vanishes near $B$, so it extends by $0$ to a closed form on $Y$. For a loop $\gamma$ in $Y'$, $\int_\gamma\Tr\theta=\sum_i\int_{s_i\circ\gamma}\theta=\int_{\varphi^{-1}(\gamma)}\theta$,
which is $[\theta](\varphi^!\gamma)$; as loops in $Y'$ represent all of $H_1(Y,\Z)$, $[\Tr\theta]=\varphi_*[\theta]$.

(iv) For a $2$-form $\psi$ on $X$, let $\Tr\psi=\sum_is_i^*\psi$ on $Y'$. On $V$ as in (i), $\int_V\Tr\psi=\sum_i\int_{s_i(V)}\psi=\int_{\varphi^{-1}(V)}\psi$ (each $s_i$ is an orientation-preserving
diffeomorphism onto its image); covering $Y'$ by countably many such $V$ and using a partition of unity, $\int_{Y'}\Tr\psi=\int_{X'}\psi$. Finite sets have measure zero, so
$\int_{X'}=\int_X$ and $\int_{Y'}=\int_Y$ for smooth forms. Now take $\theta$ as in (iii) representing $\eta$, and $\xi$ represented by a closed form, also written $\xi$. On $Y'$,
$\Tr(\theta)\wedge\xi=\sum_is_i^*\theta\wedge s_i^*\varphi^*\xi=\Tr(\theta\wedge\varphi^*\xi)$, since $\varphi\circ s_i=\mathrm{id}$. Hence
$\langle\varphi_*\eta,\xi\rangle=\int_Y\Tr(\theta)\wedge\xi=\int_X\theta\wedge\varphi^*\xi=\langle\eta,\varphi^*\xi\rangle$. For a $2$-form $\rho$ on $Y$, $\Tr(\varphi^*\rho)=n\rho$ on $Y'$, so
$\int_X\varphi^*\rho=n\int_Y\rho$. Choosing the representative $\xi$ of a class in $H^1(Y)$ to vanish near $B$ (as in (iii), with $Y$ in place of $X$), $\varphi^*\xi$ vanishes near
$\varphi^{-1}(B)$ and $\Tr(\varphi^*\xi)=n\xi$; so $\varphi_*\varphi^*=n$ by (iii). Then $\langle\varphi^*\xi,\varphi^*\xi'\rangle=\langle\varphi_*\varphi^*\xi,\xi'\rangle=n\langle\xi,\xi'\rangle$ (alternatively,
$\int_X\varphi^*(\xi\wedge\xi')=n\int_Y\xi\wedge\xi'$). Since $H^1(Y,\Z)$ is torsion-free and $\varphi_*\varphi^*=n$, $\varphi^*$ is injective; dually $\varphi_*\varphi^!=n$ on $H_1$ shows that
$nH_1(Y,\Z)\subseteq\varphi_*H_1(X,\Z)$.

(v) Functoriality of $\varphi_*$ on loops is clear, $\varphi^*$ is its transpose, and $\varphi^!$ is compatible with composition on loops avoiding the branch points of both maps.
An automorphism $w$ is unramified of degree one, so $w^{-1}(\gamma)=(w^{-1})\circ\gamma$; transposing gives the last statement.
\end{proof}

The transfer is easiest to see on the torus $X=Y=\C/\Lambda$ with $\varphi=[m]$, $z\mapsto mz$, of degree $n=m^2$ (\Cref{ex:curves:torus}(d)), which is unramified, so no
loop has to be moved. The loop $a$, the image of $s\mapsto s\omega_1$, $s\in[0,1]$, has the $m^2$ lifts $s\mapsto(s\omega_1+\lambda)/m$, one for each class
$\lambda\in\Lambda/m\Lambda$. The lift indexed by $\lambda$ ends where the lift indexed by $\lambda+\omega_1$ begins, so the lifts link up in chains of $m$ (indexed by
$\lambda,\lambda+\omega_1,\dots,\lambda+(m-1)\omega_1$), each chain forming a closed loop from $\lambda/m$ to $\lambda/m+\omega_1$, whose class in $H_1=\Lambda$ is $\omega_1$, that is $a$.
There are $m^2/m=m$ such loops, so $\varphi^!a=m\,a$, and likewise $\varphi^!b=m\,b$. As $\varphi_*=m$ on $H_1$, $\varphi_*\varphi^!=m^2=n$, as (ii) says; on $H^1$,
$\varphi^*=m$ and $\varphi_*$ (the transpose of $\varphi^!$) is also $m$, and $\varphi_*\varphi^*=m^2$ as in (iv). For a general $\varphi$ the transfer is not a multiple of
the identity, and it is the transfer, not $\varphi_*$, that is injective on $H_1$.

\begin{convention}[which is dual to which]\label{conv:curves:duality}
\leavevmode
\begin{enumerate}[label=(\alph*)]
\item \emph{Covariant objects}: $H_1(X,\Z)$, the period lattice $\Lambda_X$, the Jacobian $J(X)=H^0(X,\Omega^1)^\vee/\Lambda_X$, its torsion $J[n]=H_1(X,\Z/n)$ and its Tate module
$T_\ell J=H_1(X,\Zl)$ (\Cref{sec:curves:jacobian}). A morphism $\varphi$ acts by $\varphi_*$, a correspondence $\mathcal C=(C,\alpha,\beta)$ by $\mathcal C_*=\beta_*\alpha^!$
(\Cref{prop:curves:correspondence-actions}).
\item \emph{Contravariant objects}: $H^1(X,\Z)=\Hom(H_1,\Z)$, $H^1(X,\Zl)=\Hom_{\Zl}(T_\ell J,\Zl)$, $H^0(X,\Omega^1)$. A morphism acts by $\varphi^*$, a correspondence by
$\mathcal C^*=\alpha_*\beta^*$, and $\mathcal C^*$ is the transpose of $\mathcal C_*$ for the evaluation pairing.
\item The two sides are identified by Poincare duality (\Cref{thm:curves:cup-product}), an isometry for the cup product and the intersection pairing, which turns $\mathcal C^*$ into
$(\mathcal C^t)_*$ (\Cref{thm:curves:adjoint-transpose}(c)); it is canonical up to the global orientation sign.
\item With Galois action (\Cref{thm:curves:weil-pairing-etale}): $H^1_{\mathrm{et}}=(T_\ell J)^\vee\simeq T_\ell J(-1)$. For an elliptic curve, $V=V_\ell E$ (the paper's $V$) has determinant
$\chi_\ell$, and $H^1(E,\Ql)=V^\vee\simeq V(-1)$ has determinant $\chi_\ell^{-1}$; $\ad V$ is the same for $V$, $V^\vee$ and $V(n)$ (\Cref{ch:reps}).
\item Frobenius (\Cref{conv:curves:conventions}(C7)): the arithmetic Frobenius on $T_\ell J$ and the geometric Frobenius on $H^1$ have the same characteristic polynomial $P(T)$,
with $\#C(\F_q)=q+1-\tr$.
\end{enumerate}
\end{convention}

\begin{paperlink}\label{pl:curves:cup}
Section 4.1 of the paper (sec:dfg): the pairing $\hat\delta_!$ ``is perfect and alternating on $H^1(X_0(M^\Sigma),\Zl)$'', being $\frac12$ times the cup product composed with
$w_{M^\Sigma}$ on classes from $X_0(M^\Sigma)$. The cup product is perfect and alternating on $H^1(X_0(M^\Sigma),\Zl)$ by \Cref{thm:curves:cup-product} with $R=\Zl$; the factor $\frac12$
is a unit for odd $\ell$; composing with the automorphism $w_{M^\Sigma}^*$ preserves perfectness, and alternation needs $w^*$ to be an isometric involution
(\Cref{cor:curves:twisted-pairing}). Proposition prop:pullback: ``$\varphi^*(a\cup b)=\varphi^*a\cup\varphi^*b$ and $\varphi^*$ multiplies the fundamental class of $H^2$ by $\degphi$; so
$\langle\varphi^*a,\varphi^*b\rangle_{X_0(N)}=\degphi\,\langle a,b\rangle_E$, the projection formula'' is \Cref{prop:curves:transfer}(iv); ``the cup product on $H^1(E,\Z)$ is unimodular''
is \Cref{thm:curves:cup-product} for $g=1$ (\Cref{ex:curves:torus}).
\end{paperlink}

\section[The Jacobian; torsion and Tate modules]{The Jacobian: \texorpdfstring{$\Pic^0$}{Pic0}, complex torus, abelian variety; torsion and Tate modules}\label{sec:curves:jacobian}

The Jacobian $J(X)$ of a curve is $\Pic^0$, made into a complex torus by integrating holomorphic forms and into an abelian variety by a theorem. The paper uses it
three times: the modular parametrisation factors as $X_0(N)\to J_0(N)\to E$, and optimality of $E$ means that the kernel of $J_0(N)\to E$ is connected (Section 1,
Proposition prop:pullback); Remark rem:free speaks of $J_0(N)[\mf m]$ and of the Tate module of $J_0(N)$; and the paragraph after Proposition prop:pullback identifies the
lattice $M_{f,\lambda}$ with $H^1(E,\Zl)=\Hom(T_\ell E,\Zl)\simeq T_\ell E(-1)$. This section proves the facts behind these sentences for a general curve, in particular
\Cref{prop:curves:degree-cup-product}, which turns ``connected kernel'' into ``saturated image of $\varphi^*$ on $H^1$'' and computes the cup product on that image.

\begin{definition}[the Jacobian]\label{def:curves:jacobian}
For a smooth projective curve $C$ over $k$, $J(C)(\bar k)=\Pic^0(C)$, with the action of $G_k$ of \Cref{prop:curves:picard}(ii). For a compact Riemann surface $X$ of genus
$g\ge1$, the \emph{analytic Jacobian} is $J_{\mathrm{an}}(X)=H^0(X,\Omega^1)^\vee/\Lambda_X$, with $\Lambda_X$ the period lattice of \Cref{thm:curves:bilinear-relations}(e), a complex torus of
dimension $g$. Examples: $J(\PP^1)=0$ (\Cref{prop:curves:picard}(i)); for $g=1$ and $O\in C(k)$, $J(C)\simeq C$ by \Cref{prop:curves:genus-one-pic}.
\end{definition}

\begin{theorem}[Abel-Jacobi]\label{thm:curves:abel-jacobi}
Let $X$ be a compact Riemann surface of genus $g\ge1$.
\begin{enumerate}[label=(\roman*)]
\item The \emph{Abel-Jacobi map} $\operatorname{AJ}\colon\Div^0(X)\to J_{\mathrm{an}}(X)$, $\sum_i\bigl((P_i)-(Q_i)\bigr)\mapsto\bigl(\omega\mapsto\sum_i\int_{Q_i}^{P_i}\omega\bigr)\bmod\Lambda_X$, is a well-defined
homomorphism.
\item (Abel; stated.) $\ker\operatorname{AJ}=\operatorname{Prin}(X)$. (Jacobi inversion; stated.) $\operatorname{AJ}$ is surjective. Hence $\operatorname{AJ}\colon\Pic^0(X)\to J_{\mathrm{an}}(X)$ is an
isomorphism.
\item For a base point $P_0$, $\iota\colon X\to J_{\mathrm{an}}(X)$, $P\mapsto\operatorname{AJ}((P)-(P_0))$, is holomorphic and injective, and $\iota_*\colon H_1(X,\Z)\to H_1(J_{\mathrm{an}}(X),\Z)=\Lambda_X$
is the period isomorphism of \Cref{thm:curves:bilinear-relations}(e). For $g=1$, $\iota\colon X\to\C/\Lambda_X$ is an isomorphism.
\end{enumerate}
\end{theorem}

\begin{proof}
(i) Two paths from $Q_i$ to $P_i$ differ by a closed loop, so the integrals differ by a period, an element of $\Lambda_X$; the map is additive.
(ii) We do not prove this; see \cite[Chapter VIII]{Miranda}, \cite[Sections 20 and 21]{Forster81} and \cite[Chapter 2, Section 2]{GriffithsHarris78} (not verified).
The ``only if'' half of Abel's theorem ($\operatorname{AJ}(\operatorname{div}f)=0$) is a trace computation: for nonconstant $f\colon X\to\PP^1$,
$\operatorname{div}f=f^*((0)-(\infty))$, and as in the proof of \Cref{prop:curves:jacobian-functoriality}(iii), $\operatorname{AJ}(\operatorname{div}f)(\omega)\equiv\int_\infty^0\Tr_f\omega$
modulo periods; $\Tr_f\omega$ is a holomorphic form on $\PP^1$, hence $0$ (\Cref{prop:curves:differentials}(iv)).
(iii) In a chart near $P_1$, $\iota(P)=\iota(P_1)+(\omega\mapsto\int_{P_1}^P\omega)$, and $P\mapsto\int_{P_1}^P\omega$ is a primitive of a holomorphic function, so $\iota$ is holomorphic. If
$\iota(P)=\iota(Q)$ then $\operatorname{AJ}((P)-(Q))=0$, so $(P)\sim(Q)$ by Abel's theorem, and $P=Q$ since $g\ge1$ (\Cref{cor:curves:rr-consequences}(d)). For a loop $\gamma$ based at
$P_0$, the path $t\mapsto(\omega\mapsto\int_{\gamma|[0,t]}\omega)$ in $H^0(X,\Omega^1)^\vee$ lifts $\iota\circ\gamma$ and ends at the period vector of $\gamma$; and $H_1(V/\Lambda,\Z)=\Lambda$ for a
torus, by the endpoint of the lift of a loop through $0$ (as in \Cref{ex:curves:torus}). So $\iota_*$ is the period map. For $g=1$, $\iota$ is an injective holomorphic map
between compact Riemann surfaces, hence surjective and of degree one (\Cref{thm:curves:holomorphic-maps}), hence an isomorphism. This is the analytic form of
\Cref{prop:curves:genus-one-pic}; \Cref{ch:ellcurves} makes it explicit with the Weierstrass function $\wp$.
\end{proof}

\begin{theorem}[the Jacobian variety; stated]\label{thm:curves:jacobian-variety}
Let $C$ be a smooth projective curve of genus $g\ge1$ over $k$.
\begin{enumerate}[label=(\roman*)]
\item There is an abelian variety $J=\Jac(C)$ over $k$ (a smooth projective connected group variety) of dimension $g$ with $J(\bar k)=\Pic^0(C)$ compatibly with $G_k$; and
$J(L)\supseteq\Div^0_L(C)/\operatorname{div}(L(C)^\times)$, with equality when $C(L)\ne\emptyset$ (\Cref{prop:curves:picard}(iii)).
\item Over $\C$, $J(\C)\simeq J_{\mathrm{an}}(C(\C))$ as complex Lie groups, via $\operatorname{AJ}$; the torus is algebraic because the cup product, with the positivity of
\Cref{thm:curves:bilinear-relations}, is a Riemann form on $\Lambda_X$ (Riemann, Lefschetz).
\item (Universal property.) For $P_0\in C(k)$, $\iota\colon C\to J$, $P\mapsto[(P)-(P_0)]$, is a closed immersion defined over $k$, and every morphism $C\to A$ to an abelian variety
over $k$ sending $P_0$ to $0$ factors uniquely as $h\circ\iota$ with $h\colon J\to A$ a homomorphism over $k$.
\item (Mordell-Weil.) $J(K)$ is finitely generated for a number field $K$; in particular $J(K)_{\tors}$ is finite.
\end{enumerate}
\end{theorem}

We do not prove this; see \cite[Sections 1 to 6]{MilneJacobian86}, \cite{MilneAbelian86}, \cite[Section 24]{Mumford70} and \cite[Chapter 9]{BLR} (section numbers not
verified).

\begin{proposition}[functoriality of the Jacobian]\label{prop:curves:jacobian-functoriality}
Let $\varphi\colon X\to Y$ be a nonconstant morphism of smooth projective curves, $n=\deg\varphi$.
\begin{enumerate}[label=(\roman*)]
\item $\varphi_*\colon J(X)\to J(Y)$ (pushforward of divisor classes, ``Albanese'') and $\varphi^*\colon J(Y)\to J(X)$ (pullback, ``Picard'') are homomorphisms, $G_k$-equivariant if $\varphi$
is defined over $k$, with $\varphi_*\varphi^*=[n]$ on $J(Y)$.
\item Over $\C$, under $\operatorname{AJ}$, $\varphi_*$ is induced by the transpose $(\varphi^*)^\vee\colon H^0(X,\Omega^1)^\vee\to H^0(Y,\Omega^1)^\vee$ of the pullback of forms, and on lattices by
$\varphi_*\colon H_1(X,\Z)\to H_1(Y,\Z)$.
\item Over $\C$, $\varphi^*$ on $J$ is induced by the transpose of the trace $\Tr\colon H^0(X,\Omega^1)\to H^0(Y,\Omega^1)$ and on lattices by the transfer $\varphi^!$.
\item If $Y=E$ has genus one, with origin $O=\varphi(P_0)$ and $J(E)=E$ (\Cref{prop:curves:genus-one-pic}), then $\varphi=\varphi_*\circ\iota_X$ with $\iota_X(P)=[(P)-(P_0)]$, and
$\varphi_*\circ\varphi^*\colon E\to J(X)\to E$ is multiplication by $n$.
\end{enumerate}
\end{proposition}

\begin{proof}
(i) is \Cref{cor:curves:degree-three-ways}(iii) and \Cref{prop:curves:galois-action}(iii). (ii) For $D=\sum((P_i)-(Q_i))$ and paths $c_i$ from $Q_i$ to $P_i$,
$\operatorname{AJ}_Y(\varphi_*D)(\omega)=\sum\int_{\varphi\circ c_i}\omega=\sum\int_{c_i}\varphi^*\omega=\operatorname{AJ}_X(D)(\varphi^*\omega)$, by change of variables
(\Cref{prop:curves:transfer}(i)); on closed loops this is the statement about lattices. (iii) (Sketch.) For a holomorphic form $\omega$ on $X$, $\Tr\omega=\sum s_i^*\omega$ is
holomorphic on $Y'$, and it extends holomorphically over each branch point: in the normal form $w=z^e$ at a point over it, the $e$ local sections are
$z=\zeta_e^kw^{1/e}$, and for $\omega=\sum_mc_mz^m\,dz$ one finds $\sum_ks_k^*\omega=\sum_{e\mid m+1}c_mw^{(m+1)/e-1}\,dw$, with nonnegative exponents. For $D=(P)-(Q)$ and a path $c$
from $Q$ to $P$ through $Y'$ except at its ends, the $n$ lifts of $c$ run from the points of $\varphi^*(Q)$ to those of $\varphi^*(P)$, with multiplicities $e_\varphi$, so
$\operatorname{AJ}_X(\varphi^*D)(\omega)=\int_c\Tr\omega=\operatorname{AJ}_Y(D)(\Tr\omega)$. On loops this is $\varphi^!$. See \cite[Section 6]{MilneJacobian86} (not verified).
(iv) $\varphi_*\iota_X(P)=[(\varphi(P))-(O)]$, which is $\varphi(P)$ under $\kappa^{-1}$ of \Cref{prop:curves:genus-one-pic}; and $\varphi_*\varphi^*=[n]$ by (i).
\end{proof}

This is the universal property of \Cref{thm:curves:jacobian-variety}(iii) made explicit: $\varphi$ factors through $J(X)$, and for $X=X_0(N)$, $P_0=\infty$, the factor
$\varphi_*\colon J_0(N)\to E$ is ``the map $J_0(N)\to E$ induced by $\varphi$'' of the paper. $E$ is its own dual, so ``$E^\vee\to J_0(N)\to E$ is multiplication by $\deg\varphi_E$''
(\Cref{ch:moddeg}) is (iv).

\begin{lemma}[homomorphisms of complex tori]\label{lem:curves:tori}
Let $V$, $W$ be finite-dimensional complex vector spaces with lattices $\Lambda\subset V$, $M\subset W$.
\begin{enumerate}[label=(\alph*)]
\item Let $A\colon V\to W$ be $\C$-linear and surjective with $A(\Lambda)\subseteq M$, and $\pi\colon V/\Lambda\to W/M$ the induced homomorphism. Then $\pi$ is surjective, $A(\Lambda)$ has finite
index in $M$, the identity component of $\ker\pi$ is $(\ker A+\Lambda)/\Lambda$, and
\[
\pi_0(\ker\pi)=\ker\pi/(\ker\pi)^0\ \simeq\ M/A(\Lambda)=\coker\bigl(H_1(\pi)\bigr).
\]
\item Let $B\colon W\to V$ be $\C$-linear and injective with $B(M)\subseteq\Lambda$, and $\beta\colon W/M\to V/\Lambda$ the induced homomorphism. Then
$\ker\beta\simeq(\Lambda\cap B(W))/B(M)=\operatorname{sat}(B(M))/B(M)$, the saturation being taken in $\Lambda$; so $\beta$ is injective iff $B(M)$ is saturated in $\Lambda$.
\item A holomorphic map $\psi\colon\C/\Lambda\to\C/\Lambda'$ with $\psi(0)=0$ is $z\mapsto cz$ for a $c\in\C$ with $c\Lambda\subseteq\Lambda'$; if $c\ne0$, it is unramified of degree
$[\Lambda':c\Lambda]$, and $\psi_*\colon H_1=\Lambda\to H_1=\Lambda'$ is the inclusion $c\Lambda\subseteq\Lambda'$ (after multiplication by $c$).
\end{enumerate}
\end{lemma}

\begin{proof}
(a) $\pi$ is surjective because $A$ is. $\Lambda$ spans $V$ over $\R$, so $A(\Lambda)$ spans $W$ over $\R$; it is a subgroup of the lattice $M$ of full rank, so of finite index. Now
$\ker\pi=A^{-1}(M)/\Lambda$. The map $A$ induces $A^{-1}(M)/(\ker A+\Lambda)\to M/A(\Lambda)$; it is surjective (every $m\in M$ is $A(v)$ for some $v$, and $v\in A^{-1}(M)$) and injective
(if $A(v)=A(\lambda)$ with $\lambda\in\Lambda$ then $v\in\lambda+\ker A$). So $H=(\ker A+\Lambda)/\Lambda$ has finite index in $\ker\pi$. $H$ is connected, as the image of the vector space
$\ker A$. $H$ is closed: $\ker A+\Lambda=A^{-1}(A(\Lambda))$, and $A(\Lambda)$ is discrete, hence closed, in $W$. So $\ker\pi$ is a finite disjoint union of closed cosets of $H$, each
therefore open in $\ker\pi$, and $H$ is the identity component. Finally $H_1(V/\Lambda,\Z)=\Lambda$ and $H_1(\pi)=A|_\Lambda$, as in the proof of \Cref{thm:curves:abel-jacobi}(iii).
(b) $\ker\beta=B^{-1}(\Lambda)/M$, and $B$ maps it isomorphically onto $(\Lambda\cap B(W))/B(M)$. Since $B(M)$ spans the real space $B(W)$, $\Lambda\cap B(W)$ is the intersection of $\Lambda$ with the
real span of $B(M)$, which is the saturation (\Cref{ch:modules}).
(c) $\psi$ lifts to a holomorphic $\tilde\psi\colon\C\to\C$ with $\tilde\psi(0)=0$ (the plane is simply connected). For $\lambda\in\Lambda$, $\tilde\psi(z+\lambda)-\tilde\psi(z)$ is a continuous
function with values in the discrete $\Lambda'$, hence constant; so $\tilde\psi'$ is $\Lambda$-periodic and entire, hence bounded and constant (Liouville), and $\tilde\psi(z)=cz$ with
$c\Lambda\subseteq\Lambda'$. If $c\ne0$, $\psi$ is a local biholomorphism, so unramified, and its fibre over $0$ is $c^{-1}\Lambda'/\Lambda\simeq\Lambda'/c\Lambda$; its degree is the order of the
fibre (\Cref{thm:curves:holomorphic-maps}(c)). The last claim is the lift of loops.
\end{proof}

\begin{proposition}[the degree and the cup product; saturation]\label{prop:curves:degree-cup-product}
Let $\varphi\colon X\to Y$ be a nonconstant holomorphic map of compact Riemann surfaces of genus $\ge1$, $n=\deg\varphi$.
\begin{enumerate}[label=(\roman*)]
\item $\langle\varphi^*x,\varphi^*y\rangle_X=n\,\langle x,y\rangle_Y$ for $x,y\in H^1(Y,\Z)$.
\item $\varphi_*\varphi^*=n$ on $H^1(Y,\Z)$, $\varphi_*\varphi^!=n$ on $H_1(Y,\Z)$ and $\varphi_*\varphi^*=[n]$ on $J(Y)$. Moreover $\operatorname{PD}_X\circ\varphi^*=\varphi^!\circ\operatorname{PD}_Y$.
\item The following are equivalent: (a) $\varphi^*H^1(Y,\Z)$ is saturated in $H^1(X,\Z)$; (b) $\varphi_*\colon H_1(X,\Z)\to H_1(Y,\Z)$ is surjective; (c) $\ker(\varphi_*\colon J(X)\to J(Y))$ is
connected; (d) $\varphi^*\colon J(Y)\to J(X)$ is injective. More precisely
\begin{align*}
[\operatorname{sat}(\varphi^*H^1(Y,\Z)):\varphi^*H^1(Y,\Z)]&=\#\coker\bigl(\varphi_*\colon H_1(X,\Z)\to H_1(Y,\Z)\bigr)\\
&=\#\pi_0\bigl(\ker(\varphi_*\colon J(X)\to J(Y))\bigr)\\
&=\#\ker(\varphi^*\colon J(Y)\to J(X)),
\end{align*}
a divisor of $n^{2g_Y}$.
\item Let $Y=E$ have genus one and $M=\varphi^*H^1(E,\Z)$. The cup product on $M$ has ideal of values $n\Z$ and Gram determinant $n^2$. If (iii) holds,
$[H^1(X,\Z):M\oplus M^\perp]=n^2$, and for every prime $\ell$, $[H^1(X,\Zl):(M\otimes\Zl)\oplus(M\otimes\Zl)^\perp]=\ell^{2v_\ell(n)}$. In general, with $m_0=\#\pi_0(\ker(J(X)\to E))$,
$m_0$ divides $n$ and the cup product on $\operatorname{sat}(M)$ has ideal of values $(n/m_0)\Z$.
\item (Isogenies.) Let $Y=E$, $\psi\colon E\to E'$ an isogeny of degree $m$ (a nonconstant holomorphic map of genus-one surfaces with $\psi(O)=O'$) and $\varphi'=\psi\circ\varphi$. Then
$(\varphi')^*H^1(E',\Z)$ has index $m$ in $\varphi^*H^1(E,\Z)$. Hence if $\varphi^*H^1(E,\Z)$ is saturated and $m>1$, then $(\varphi')^*H^1(E',\Z)$ is not.
\end{enumerate}
\end{proposition}

\begin{proof}
(i) and the first two identities of (ii) are \Cref{prop:curves:transfer}(ii), (iv); the third is \Cref{prop:curves:jacobian-functoriality}(i). For the last, let
$\xi\in H^1(Y,\Z)$ and $\eta\in H^1(X,\Z)$. Then $\eta(\operatorname{PD}_X\varphi^*\xi)=\langle\varphi^*\xi,\eta\rangle_X$, while
$\eta(\varphi^!\operatorname{PD}_Y\xi)=(\varphi_*\eta)(\operatorname{PD}_Y\xi)=\langle\xi,\varphi_*\eta\rangle_Y=-\langle\varphi_*\eta,\xi\rangle_Y=-\langle\eta,\varphi^*\xi\rangle_X=\langle\varphi^*\xi,\eta\rangle_X$, by the projection
formula and alternation.

(iii) Put $L=H_1(X,\Z)$, $L'=H_1(Y,\Z)$, $u=\varphi_*\colon L\to L'$. Its cokernel $Q$ is finite and killed by $n$ (\Cref{prop:curves:transfer}(iv)). By the Smith normal form (\Cref{ch:modules})
there are bases of $L$ and $L'$ in which $u$ is $(D\mid0)$ with $D=\operatorname{diag}(d_1,\dots,d_r)$, $r=2g_Y$, all $d_i\ne0$, and $Q\simeq\bigoplus\Z/d_i$. In the dual bases
$\varphi^*=u^{\mathsf T}=\binom D0$, whose image $\bigoplus d_i\Z\oplus0$ has saturation $\Z^r\oplus0$, of index $\prod d_i=\#Q$. This gives the first equality and (a)$\Leftrightarrow$(b).
By \Cref{prop:curves:jacobian-functoriality}(ii), $\varphi_*$ on $J$ is the map of \Cref{lem:curves:tori}(a) with $A=(\varphi^*)^\vee$, surjective because $\varphi^*$ is injective on holomorphic forms
($\varphi$ is a local biholomorphism off a finite set), and with $H_1(\varphi_*)=u$; so $\pi_0(\ker)\simeq\coker u$, which is the second equality and (b)$\Leftrightarrow$(c). By
\Cref{prop:curves:jacobian-functoriality}(iii) (sketched), $\varphi^*$ on $J$ is the map of \Cref{lem:curves:tori}(b) with lattice map $\varphi^!$ (injective, as $\varphi_*\varphi^!=n$); so
$\#\ker(\varphi^*)$ is the index of $\varphi^!H_1(Y,\Z)$ in its saturation in $H_1(X,\Z)$. Poincare duality is an isomorphism carrying $\varphi^*H^1(Y,\Z)$ to $\varphi^!H_1(Y,\Z)$, by (ii),
so this index is the first term. Hence the last equality and (a)$\Leftrightarrow$(d). As $nL'\subseteq u(L)$, $\#Q$ divides $n^{2g_Y}$.

(iv) The cup product on $H^1(E,\Z)$ is perfect and alternating of rank two, so its ideal of values is $\Z$ (\Cref{ch:modules}: a rank-two alternating form with Gram matrix
$\sm0a{-a}0$ has ideal of values $a\Z$ and determinant $a^2$). By (i) and the injectivity of $\varphi^*$, the cup product on $M$ is $n$ times a unimodular alternating form: ideal of
values $n\Z$, Gram determinant $n^2$. If $M$ is saturated, then $M\otimes\Zl$ is saturated in $H^1(X,\Zl)$, the cup product is perfect there (\Cref{thm:curves:cup-product}), and the
theorem on Gram determinants of saturated sublattices (\Cref{ch:modules}) gives $[H^1(X,\Zl):(M\otimes\Zl)\oplus(M\otimes\Zl)^\perp]=\ell^{v_\ell(n^2)}$. The orthogonal commutes with
$\otimes\Zl$ (it is a kernel, and $\Zl$ is flat over $\Z$), and an index of lattices is the product of its $\ell$-parts, which gives $n^2$. In general $[\operatorname{sat}(M):M]=m_0$ by (iii),
so the Gram determinant of $\operatorname{sat}(M)$ is $n^2/m_0^2$, an integer square, and its ideal of values is $(n/m_0)\Z$; in particular $m_0\mid n$.

(v) By \Cref{lem:curves:tori}(c), in suitable coordinates $E=\C/\Lambda$, $E'=\C/\Lambda'$, $\psi(z)=z$ with $\Lambda\subseteq\Lambda'$ of index $m$, and $\psi_*$ on $H_1$ is this inclusion. Its transpose
$\psi^*\colon H^1(E',\Z)\to H^1(E,\Z)$ has image of index $m$ (the determinant of a matrix and of its transpose agree). As $\varphi^*$ is injective on $H^1(E,\Z)$ with image
$\varphi^*H^1(E,\Z)$, the image $\varphi^*\psi^*H^1(E',\Z)$ has index $m$ in it. If $\varphi^*H^1(E,\Z)$ is saturated and $m>1$, then $(\varphi')^*H^1(E',\Z)$ is a proper sublattice of finite index
of a saturated lattice of the same rank, so it is not saturated.
\end{proof}

All the numbers of the proposition can be seen on the least saturated example, $\varphi=[m]\colon E\to E$, $m\ge2$, of degree $n=m^2$. Here $\varphi^*=m$ on
$H^1(E,\Z)$, so $M=\varphi^*H^1(E,\Z)=mH^1(E,\Z)$, whose saturation is $H^1(E,\Z)$, of index $m^2$. On the other side, $\varphi_*=m$ on $H_1(E,\Z)\simeq\Z^2$ has cokernel
$(\Z/m)^2$, of order $m^2$; the kernel of $[m]$ on $J(E)=E$ is $E[m]$, finite of order $m^2$ (\Cref{prop:curves:torsion}(i)), so it has $m^2$ components; and
$\varphi^*=[m]$ on $J$ has the same kernel. All four quantities of (iii) are $m^2$, a divisor of $n^{2g_Y}=m^4$. In (iv), the Gram matrix on $M=mH^1$ is $m^2$ times the
unimodular one, so the ideal of values is $m^2\Z=n\Z$ and the determinant $m^4=n^2$; here $m_0=m^2$ and the ideal of values on $\operatorname{sat}(M)=H^1(E,\Z)$ is
$(n/m_0)\Z=\Z$. At the other extreme, a parametrisation of degree one ($X_0(11)\to\curve{11a1}$) is an isomorphism and every index is $1$.

\begin{proposition}[torsion of Jacobians]\label{prop:curves:torsion}
\leavevmode
\begin{enumerate}[label=(\roman*)]
\item For a complex torus $V/\Lambda$ of dimension $g$ and $n\ge1$, $(V/\Lambda)[n]=\frac1n\Lambda/\Lambda\simeq\Lambda/n\Lambda\simeq(\Z/n)^{2g}$, the middle map being $x\mapsto nx$. Hence
$J(X)[n]\simeq H_1(X,\Z/n)$ for a compact Riemann surface $X$, and $[n]\colon J(X)\to J(X)$ has $n^{2g}$ points in each fibre.
\item For a curve $C$ over $k$ with $\operatorname{char}k\nmid n$, $J[n](\bar k)\simeq(\Z/n)^{2g}$ as groups, with a continuous action of $G_k$ (stated).
\item For a commutative ring $\T$ of endomorphisms of $J$ and an ideal $\mf a\subseteq\T$, put $J[\mf a]=\{P:tP=0\text{ for all }t\in\mf a\}$. Then $J[\ell]$ is a $\T/\ell\T$-module, and for a
maximal ideal $\mf m\ni\ell$, $J[\mf m]$ is a vector space over the field $\T/\mf m$.
\end{enumerate}
\end{proposition}

\begin{proof}
(i) $x\in V$ maps to an $n$-torsion point iff $nx\in\Lambda$. Multiplication by $n$ maps $\frac1n\Lambda$ isomorphically onto $\Lambda$ and $\Lambda$ onto $n\Lambda$, so
$\frac1n\Lambda/\Lambda\simeq\Lambda/n\Lambda$; and $\Lambda\simeq\Z^{2g}$. With \Cref{thm:curves:abel-jacobi} and $\Lambda_X\simeq H_1(X,\Z)$ this gives the second statement; $[n]$ is a surjective homomorphism, so each fibre is a coset of
$J[n]$. (ii) is \cite[Section 8]{MilneAbelian86} (not verified); we do not prove it. (iii) $\T$ acts on $J[\mf a]$ because $\T$ is commutative; $\ell\T$ kills $J[\ell]$ and $\mf m$
kills $J[\mf m]$.
\end{proof}

\begin{definition}[Tate modules]\label{def:curves:tate-module}
For a prime $\ell$, $T_\ell J=\varprojlim_mJ[\ell^m]$, with transition maps multiplication by $\ell$, and $V_\ell J=T_\ell J\otimes\Ql$ (\Cref{ch:groups} for inverse limits and $\Zl$).
\end{definition}

\begin{proposition}[the Tate module and $H^1$]\label{prop:curves:tate-module}
Let $X$ be a compact Riemann surface of genus $g$ and $J=J(X)$.
\begin{enumerate}[label=(\roman*)]
\item The isomorphisms $J[\ell^m]=\ell^{-m}\Lambda_X/\Lambda_X\to\Lambda_X/\ell^m\Lambda_X$, $x\mapsto\ell^mx$, are compatible with the transition maps (multiplication by $\ell$ on the left,
reduction on the right) and give a canonical isomorphism $T_\ell J\simeq\Lambda_X\otimes\Zl=H_1(X,\Zl)$; $T_\ell J$ is free of rank $2g$ and $T_\ell J/\ell^mT_\ell J=J[\ell^m]$.
\item (Duality.) $H^1(X,\Zl)=\Hom(H_1(X,\Z),\Zl)\simeq\Hom_{\Zl}(T_\ell J,\Zl)$ canonically: $H^1$ is the $\Zl$-dual of the Tate module.
\item Composing with Poincare duality, $H^1(X,\Zl)\simeq H_1(X,\Zl)\simeq T_\ell J$ is an isomorphism of $\Zl$-modules. It depends on the orientation, it is not
Galois-equivariant when $X$ comes from a curve over a number field (\Cref{warn:curves:not-galois}), and it carries the action $\mathcal C^*$ of a correspondence on $H^1$ to the action
$(\mathcal C^t)_*$ of the transpose on $T_\ell J$ (\Cref{thm:curves:adjoint-transpose}(c)).
\item For a curve $C$ over $k$ and $\ell\ne\operatorname{char}k$, $T_\ell J$ is free of rank $2g$ over $\Zl$ with a continuous action of $G_k$ (stated). For $g=1$, $T_\ell J(E)=T_\ell E$
(\Cref{ch:ellcurves}).
\end{enumerate}
\end{proposition}

\begin{proof}
(i) If $x\in\ell^{-m-1}\Lambda_X$, then $\ell x\in\ell^{-m}\Lambda_X$ is sent to $\ell^m(\ell x)=\ell^{m+1}x$, the image of the class of $\ell^{m+1}x$ under reduction modulo $\ell^m$: the squares commute.
So $T_\ell J\simeq\varprojlim\Lambda_X/\ell^m\Lambda_X=\Lambda_X\otimes\Zl$ (for a free module of finite rank, $\varprojlim_m\Z^{2g}/\ell^m=\Zl^{2g}$), and $\Lambda_X\simeq H_1(X,\Z)$ by
\Cref{thm:curves:abel-jacobi}(iii). (ii) $\Hom(H_1(X,\Z),\Zl)=\Hom_{\Zl}(H_1(X,\Z)\otimes\Zl,\Zl)$. (iii) is \Cref{thm:curves:cup-product} tensored with $\Zl$, and the last clause is
proved in \Cref{thm:curves:adjoint-transpose}(c). (iv) is \cite[Section 15]{MilneAbelian86} and \cite[Section 18]{Mumford70} (not verified); we do not prove it.
\end{proof}

\begin{theorem}[the Weil pairing and etale cohomology; stated]\label{thm:curves:weil-pairing-etale}
Let $C$ be a smooth projective curve of genus $g$ over $k$, $J$ its Jacobian and $\ell\ne\operatorname{char}k$.
\begin{enumerate}[label=(\alph*)]
\item (Weil pairing.) There is a perfect, alternating, $G_k$-equivariant $\Zl$-bilinear pairing $e_\ell\colon T_\ell J\times T_\ell J\to\Zl(1)$ (from the canonical principal polarisation).
Over $\C$, after identifying $\Zl(1)$ with $\Zl$ by $\zeta_{\ell^m}=e^{2\pi i/\ell^m}$, it is the intersection pairing on $H_1(X,\Zl)$ up to a universal sign. For $g=1$ it is the Weil pairing
of \Cref{ch:ellcurves}, and $\det\rho_{E,\ell}=\chi_\ell$.
\item (Etale cohomology.) $H^1_{\mathrm{et}}(C_{\bar k},\Zl)$ is free of rank $2g$ with a continuous action of $G_k$, and $H^1_{\mathrm{et}}(C_{\bar k},\Zl)\simeq\Hom_{\Zl}(T_\ell J,\Zl)$ as
$G_k$-modules (contragredient action, \Cref{ch:reps}). For $k\subseteq\C$ there is a comparison isomorphism $H^1_{\mathrm{et}}(C_{\bar k},\Zl)\simeq H^1(C(\C),\Z)\otimes\Zl$, compatible with
\Cref{prop:curves:tate-module}(ii), with cup products, and with pullbacks and pushforwards along morphisms.
\item Hence $H^1_{\mathrm{et}}(C_{\bar k},\Zl)\simeq T_\ell J(-1)=T_\ell J\otimes\Zl(-1)$ as $G_k$-modules, and the cup product $H^1\times H^1\to H^2_{\mathrm{et}}=\Zl(-1)$ is perfect, alternating and
$G_k$-equivariant. For $g=1$: $H^1(E,\Zl)\simeq(T_\ell E)^\vee\simeq T_\ell E(-1)$, with determinant $\chi_\ell^{-1}$.
\item (Weil.) For $k=\F_q$: the arithmetic Frobenius acts on $T_\ell J$ as the Frobenius endomorphism $\pi_q$ of $J$; $P(T)=\det(T-\pi_q\mid T_\ell J)$ has integer coefficients independent of
$\ell$, degree $2g$ and roots $\alpha_i$ with $\abs{\alpha_i}=\sqrt q$; $\#C(\F_{q^m})=q^m+1-\sum_i\alpha_i^m$; and on $H^1=(T_\ell J)^\vee$ the geometric Frobenius has the same eigenvalues
$\alpha_i$. For $g=1$, $a_q=q+1-\#E(\F_q)=\tr(\Frob_q\mid T_\ell E)=\tr(\Frob_q^{-1}\mid H^1(E))$.
\end{enumerate}
\end{theorem}

We do not prove this; see \cite[Sections 16 and 19]{MilneAbelian86} and \cite{MilneJacobian86} for (a), (d); \cite[Section 20]{Mumford70} for the Weil pairing; \cite{MilneEtale}
for (b) and the comparison theorem (all section numbers not verified); and \cite[Chapter III, Section 8, and Chapter V, Section 2]{SilvermanAEC} for $g=1$. Part (c) follows from
(a) and (b): the Weil pairing gives $T_\ell J\simeq\Hom(T_\ell J,\Zl(1))=(T_\ell J)^\vee(1)$, so $(T_\ell J)^\vee\simeq T_\ell J(-1)$. In (d), under the contragredient action $\Frob_q$ acts on
$(T_\ell J)^\vee$ with eigenvalues $\alpha_i^{-1}$, so the geometric Frobenius $\Frob_q^{-1}$ acts with eigenvalues $\alpha_i$ (\Cref{ch:notation}, B-C3).

\begin{warning}[$H^1\simeq T_\ell J$ is not Galois-equivariant]\label{warn:curves:not-galois}
The isomorphism $H^1(X,\Zl)\simeq T_\ell J$ of \Cref{prop:curves:tate-module}(iii) uses Poincare duality, which over a number field is the Weil pairing and carries a Tate twist:
the Galois-equivariant statement is $H^1_{\mathrm{et}}\simeq T_\ell J(-1)$ (\Cref{thm:curves:weil-pairing-etale}(c)). For an elliptic curve, $T_\ell E$ has determinant $\chi_\ell$ and
$H^1(E,\Zl)$ has determinant $\chi_\ell^{-1}$; they are not isomorphic as Galois modules. They have the same adjoint: $\ad(T^\vee)\simeq\ad T=\ad(T(-1))$ (\Cref{ch:reps},
\Cref{ch:notation}, B-C12(b)). Moreover $H^1$ carries the action $\mathcal C^*$ of correspondences and $T_\ell J$ the action $\mathcal C_*$, and Poincare duality turns $\mathcal C^*$ into
$(\mathcal C^t)_*$, not $\mathcal C_*$.
\end{warning}

\begin{paperlink}\label{pl:curves:pullback}
Proposition prop:pullback of the paper is \Cref{prop:curves:degree-cup-product} applied to $\varphi=\varphi_E\colon X_0(N)\to E$, with two inputs from
\Cref{ch:moddeg}. Step by step:
\begin{enumerate}[label=(\alph*)]
\item ``By optimality the kernel $B$ of the map $J_0(N)\to E$ induced by $\varphi$ is connected; the exact sequence $H_1(J_0(N),\Z)\to H_1(E,\Z)\to\pi_0(B)=0$ shows that $\varphi_*$ is surjective,
so its dual $\varphi^*$ is injective with torsion-free cokernel'': \Cref{prop:curves:degree-cup-product}(iii), (c)$\Rightarrow$(b)$\Rightarrow$(a), with \Cref{lem:curves:tori}(a) for the exact
sequence and \Cref{prop:curves:jacobian-functoriality}(iv) for the factorisation; ``the same holds after tensoring with $\Zl$'' because a saturated sublattice is a direct summand.
\item ``$\varphi^*H^1(E,\C)$, which is spanned by $\varphi^*\omega_E$ and its conjugate, is $\C f\,dz\oplus\C\overline{f\,dz}$'': \Cref{prop:curves:cohomology}(iii) for $E$, and
$\varphi_E^*\omega_E=c_E\cdot2\pi if(z)\,dz$ (\Cref{ch:moddeg}). ``The image of $\varphi^*$ is a saturated lattice of the same rank in the same plane; so the two coincide'': two saturated
sublattices with the same rational span are equal. The book reads this step over $\Q$: the Hecke operators act on $H^1(X_0(N),\Z)$, so $H^1(X_0(N),\Q)[I_f]$ is a $\Q$-form of the plane
$H^1(X_0(N),\C)[I_f]$ (kernels commute with field extension), of dimension two; it contains $\varphi^*H^1(E,\Q)$, also of dimension two; and $H^1(X_0(N),\Zl)[I_f]=H^1(X_0(N),\Z)[I_f]\otimes\Zl$
(kernels commute with the flat extension $\Z\to\Zl$). The paper's ``intersection with $H^1(X_0(N),\Z_\ell)$'' of a complex plane should be read this way.
\item ``$\varphi^*$ multiplies the fundamental class of $H^2$ by $\degphi$; so $\langle\varphi^*a,\varphi^*b\rangle=\degphi\langle a,b\rangle_E$'': \Cref{prop:curves:transfer}(iv).
\item ``$w_N$ \dots preserves the plane, on which it acts by $\varepsilon_N(f)$'': the plane is also the eigenspace of the anemic Hecke algebra for the eigenvalues of $f$ (strong multiplicity one
for the newform $f$, \Cref{ch:hecke}), which $w_N$ preserves; then $w_N^*(f\,dz)=(f|_2W_N)\,dz=\varepsilon_N(f)f\,dz$ (\Cref{ch:modforms}), and likewise on the conjugate since
$\varepsilon_N(f)=\pm1$ is real. Hence $\langle\varphi^*a,w_N^*\varphi^*b\rangle=\varepsilon_N(f)\degphi\langle a,b\rangle_E$.
\item ``The Gram determinant of a rank-two alternating form is the square of its ideal of values'': \Cref{ch:modules}; this is \Cref{prop:curves:degree-cup-product}(iv).
\end{enumerate}
\Cref{prop:curves:degree-cup-product}(iv), (v) also explain ``minimal degree'' in the definition of the optimal curve: for a non-optimal parametrisation the image of $\varphi^*$ is not saturated, and the
ideal of values on its saturation is $(n/m_0)\Z$, with $m_0$ the order of the component group of \Cref{prop:curves:degree-cup-product}(iv). Writing $\varphi=\psi\circ\varphi_{\mathrm{opt}}$ with $\psi$ an isogeny, the kernel of $\varphi_{\mathrm{opt},*}$
is connected, so $m_0=\#\ker\psi=\deg\psi$ and $n/m_0=\deg\varphi_{\mathrm{opt}}$ is the degree of the optimal parametrisation through which $\varphi$ factors (\Cref{ch:moddeg}). Over $\Zl$ only the $\ell$-part of $m_0$ matters, which is why,
when $E[\ell]$ is irreducible, the choice of curve in the class does not change $v_\ell$ (Section 1 of the paper).
\end{paperlink}

\begin{paperlink}\label{pl:curves:free}
The paragraph after Proposition prop:pullback: ``The lattice $M_{f,\lambda}$ of DFG04 is therefore $H^1(E,\Zl)=\Hom(T_\ell E,\Zl)\simeq T_\ell E(-1)$, by Galois-equivariance of $\varphi^*$ in etale
cohomology, and its adjoint is $\ad T_\ell E$''. This is \Cref{prop:curves:tate-module}(ii) and \Cref{thm:curves:weil-pairing-etale}(b), (c): the comparison isomorphism is compatible with
$\varphi_E^*$, and $\varphi_E$ is defined over $\Q$. Remark rem:free: ``$J_0(N)[\mf m]$ is two-dimensional'' is a statement about the $\T/\mf m$-vector space of
\Cref{prop:curves:torsion}(iii); ``the Tate module is free of rank two'' concerns $T_\ell J_0(N)_{\mf m}$ with the action $t_*$; and ``the identification of $T_\ell J_0(N)$ with $H^1(X_0(N),\Zl)$ as Hecke
modules uses Poincare duality composed with $w_N$, which is not Galois equivariant'' is \Cref{prop:curves:tate-module}(iii) and \Cref{warn:curves:not-galois}: $\operatorname{PD}$ turns $t^*$ into
$(t^t)_*=(w_Ntw_N)_*$ (\Cref{thm:curves:adjoint-transpose}(c) and \Cref{ch:moddeg}), so $x\mapsto w_{N*}\operatorname{PD}(x)$ is Hecke-equivariant.
\end{paperlink}

\section{Correspondences, their actions, and adjoints}\label{sec:curves:correspondences}

Hecke operators are correspondences: for $p\nmid N$, $T_p$ on $X_0(N)$ is the pair of maps $X_0(Np)\rightrightarrows X_0(N)$ given by $z\mapsto z$ and $z\mapsto pz$ (\Cref{ch:moddeg}). This section defines
correspondences between curves, their actions on divisors, the Jacobian, $H_1$, $H^1$ and holomorphic differentials, and proves the general theorem behind the sentence of
Section 4.1 of the paper ``the Hecke operators, $U_p$ included, are self-adjoint for it, because $w_{M^\Sigma}T_nw_{M^\Sigma}^{-1}$ is the adjoint of $T_n$ for the cup product'': the
adjoint of a correspondence for the cup product is its transpose (\Cref{thm:curves:adjoint-transpose}). \Cref{cor:curves:twisted-pairing} then turns ``the transpose of every $t$ is
$wtw$'' into ``every $t$ is self-adjoint for the perfect alternating pairing $\langle x,w^*y\rangle$'', and gives the orthogonal splitting into local components used in Lemma lem:hecke
and in Step 2 of the proof of Theorem thm:quant. The specific identity $T_n^t=w_NT_nw_N$ on $X_0(N)$ is proved in \Cref{ch:moddeg}.

\begin{definition}[correspondences]\label{def:curves:correspondence}
Let $X$, $Y$ be smooth projective curves over $k$ (or compact Riemann surfaces). A \emph{correspondence} $\mathcal C$ from $X$ to $Y$ is a triple $(C,\alpha,\beta)$ with $C$ a finite disjoint
union of smooth projective curves and $\alpha\colon C\to X$, $\beta\colon C\to Y$ morphisms nonconstant on every component of $C$. Its \emph{transpose} is $\mathcal C^t=(C,\beta,\alpha)$, from $Y$
to $X$. The \emph{sum} of two correspondences is the disjoint union. For automorphisms $u$ of $Y$ and $v$ of $X$ put $u\mathcal Cv=(C,v^{-1}\circ\alpha,u\circ\beta)$.
\end{definition}

Examples. (a) The \emph{graph} $\Gamma_\varphi=(X,\mathrm{id},\varphi)$ of a nonconstant morphism $\varphi\colon X\to Y$, and its transpose $\Gamma_\varphi^t=(X,\varphi,\mathrm{id})$. (b) An automorphism $w$ of
$X$ gives $\Gamma_w$, whose transpose acts as $\Gamma_{w^{-1}}$ (\Cref{prop:curves:correspondence-actions}(c) and \Cref{prop:curves:transfer}(v)). (c) (Forward pointer,
\Cref{ch:moddeg}.) For a prime $p\nmid N$, the Hecke correspondence $T_p=(X_0(Np),\alpha_1,\alpha_p)$ on $X_0(N)$, with $\alpha_1(z)=z$ and $\alpha_p(z)=pz$ on $\HH$. For $p\mid N$ the same triple is $U_p^t=w_NU_pw_N$, not $U_p$, and $U_p$ is
$(X_0(Np),\alpha_p,\alpha_1)$ (\Cref{ch:moddeg}). The image of $(\alpha,\beta)\colon C\to X\times Y$ is a divisor on the surface $X\times Y$, and correspondences up to the obvious equivalence form a ring under composition (fibre
products); we need neither, because the Hecke algebra is defined as the ring of operators that the correspondences generate (\Cref{ch:hecke}).

\begin{definition}[actions of a correspondence]\label{def:curves:correspondence-actions}
Let $\mathcal C=(C,\alpha,\beta)$ be a correspondence from $X$ to $Y$ (over $\C$ for (ii) to (iv)).
\begin{enumerate}[label=(\roman*)]
\item On divisors and the Jacobian: $\mathcal C_*=\beta_*\alpha^*\colon\Div(X)\to\Div(Y)$; it multiplies degrees by $\deg\alpha$ and maps principal divisors to principal divisors
(\Cref{cor:curves:degree-three-ways}(iii)), so it induces $\mathcal C_*\colon J(X)\to J(Y)$.
\item On homology: $\mathcal C_*=\beta_*\alpha^!\colon H_1(X,\Z)\to H_1(Y,\Z)$, with the transfer $\alpha^!$ of \Cref{prop:curves:transfer}(ii).
\item On cohomology: $\mathcal C^*=\alpha_*\beta^*\colon H^1(Y,\Z)\to H^1(X,\Z)$, with $\alpha_*$ the pushforward of \Cref{prop:curves:transfer}(iii).
\item On holomorphic differentials: $\mathcal C^*\omega=\Tr_\alpha(\beta^*\omega)$, the trace along $\alpha$ of \Cref{prop:curves:jacobian-functoriality}(iii), which is holomorphic.
\end{enumerate}
The same definitions apply with coefficients in $\Zl$, $\Ql$ or $\C$, and componentwise when $C$ is disconnected.
\end{definition}

\begin{proposition}[compatibilities]\label{prop:curves:correspondence-actions}
Let $\mathcal C=(C,\alpha,\beta)$ be a correspondence from $X$ to $Y$.
\begin{enumerate}[label=(\alph*)]
\item $(\mathcal C^*\xi)(\gamma)=\xi(\mathcal C_*\gamma)$ for $\xi\in H^1(Y,\Z)$ and $\gamma\in H_1(X,\Z)$: \emph{$\mathcal C^*$ on $H^1$ is the transpose of $\mathcal C_*$ on $H_1$}.
\item Under Abel-Jacobi, $\mathcal C_*$ on $J(X)$ is induced by the transpose of $\mathcal C^*$ on holomorphic differentials and, on lattices, by $\mathcal C_*$ on $H_1$. Hence on
$T_\ell J(X)\simeq H_1(X,\Zl)$ a correspondence acts by $\mathcal C_*$, and on $H^1(X,\Zl)=\Hom(T_\ell J(X),\Zl)$ by the transpose.
\item For a graph: $\Gamma_\varphi^*=\varphi^*$ on $H^1$ and on differentials; $(\Gamma_\varphi)_*=\varphi_*$ on $H_1$ and on $J$ (Albanese); $(\Gamma_\varphi^t)_*=\varphi^!$ on $H_1$ and $\varphi^*$ on $J$
(Picard); $(\Gamma_\varphi^t)^*=\varphi_*$ on $H^1$.
\item The actions are additive in $\mathcal C$, and $u\mathcal Cv$ acts by $u_*\mathcal C_*v_*$ on $H_1$ and by $v^*\mathcal C^*u^*$ on $H^1$; moreover $(u\mathcal Cv)^t=v^{-1}\mathcal C^tu^{-1}$.
\end{enumerate}
\end{proposition}

\begin{proof}
(a) $(\alpha_*\beta^*\xi)(\gamma)=(\beta^*\xi)(\alpha^!\gamma)=\xi(\beta_*\alpha^!\gamma)$, by the definitions of $\alpha_*$ on $H^1$ as the transpose of $\alpha^!$ and of $\beta^*$ as the transpose of
$\beta_*$. (b) (Sketch.) Compose \Cref{prop:curves:jacobian-functoriality}(iii) for $\alpha$ (pullback along $\alpha$: transpose of the trace, transfer on lattices) with (ii) for $\beta$
(pushforward: transpose of $\beta^*$, $\beta_*$ on lattices). The statement on $T_\ell J$ follows from \Cref{prop:curves:tate-module}(i), which is natural for homomorphisms of tori, and the one
on $H^1$ from (a). (c) Substitute $\alpha=\mathrm{id}$, $\beta=\varphi$, respectively $\alpha=\varphi$, $\beta=\mathrm{id}$; on $J$, $(\Gamma_\varphi^t)_*D=\mathrm{id}_*\varphi^*D=\varphi^*D$.
(d) Additivity is clear. By \Cref{prop:curves:transfer}(v), $(u\beta)_*(v^{-1}\alpha)^!=u_*\beta_*\alpha^!(v^{-1})^!=u_*\beta_*\alpha^!v_*$, and $(v^{-1}\alpha)_*(u\beta)^*=(v^{-1})_*\alpha_*\beta^*u^*=v^*\alpha_*\beta^*u^*$
on $H^1$. The transpose of $(C,v^{-1}\alpha,u\beta)$ is $(C,u\beta,v^{-1}\alpha)=v^{-1}(C,\beta,\alpha)u^{-1}$ in the notation of \Cref{def:curves:correspondence}.
\end{proof}

\begin{theorem}[the adjoint of a correspondence is its transpose]\label{thm:curves:adjoint-transpose}
Let $\mathcal C=(C,\alpha,\beta)$ be a correspondence from $X$ to $Y$ (compact Riemann surfaces). For all $\xi\in H^1(Y,\Z)$ and $\eta\in H^1(X,\Z)$,
\[
\langle\mathcal C^*\xi,\eta\rangle_X=\langle\xi,(\mathcal C^t)^*\eta\rangle_Y,
\]
that is, for the cup products, $(\mathcal C^*)^\dagger=(\mathcal C^t)^*$. Consequently:
\begin{enumerate}[label=(\alph*)]
\item for an automorphism $w$ of $X$, $\langle w^*x,w^*y\rangle=\langle x,y\rangle$: $w^*$ is an isometry, with adjoint $(w^{-1})^*$;
\item for a nonconstant $\varphi\colon X\to Y$, $\langle\varphi^*\xi,\eta\rangle_X=\langle\xi,\varphi_*\eta\rangle_Y$, and \Cref{prop:curves:degree-cup-product}(i), (ii) follow again;
\item $\operatorname{PD}_X\circ\mathcal C^*=(\mathcal C^t)_*\circ\operatorname{PD}_Y$: under Poincare duality the action $\mathcal C^*$ on $H^1$ corresponds to the action of the transpose on $H_1$, hence on $T_\ell J$;
\item the same holds with coefficients in $\Zl$, $\Ql$ or $\C$, and for $\Zl$-linear combinations of correspondences.
\end{enumerate}
\end{theorem}

\begin{proof}
The projection formula (\Cref{prop:curves:transfer}(iv)) and alternation give, for a nonconstant $\varphi\colon X_1\to Y_1$, $\theta\in H^1(X_1)$, $\zeta\in H^1(Y_1)$, both
\[
\langle\varphi_*\theta,\zeta\rangle_{Y_1}=\langle\theta,\varphi^*\zeta\rangle_{X_1}\quad\text{and}\quad\langle\zeta,\varphi_*\theta\rangle_{Y_1}=-\langle\varphi_*\theta,\zeta\rangle_{Y_1}=-\langle\theta,\varphi^*\zeta\rangle_{X_1}=\langle\varphi^*\zeta,\theta\rangle_{X_1}.
\]
Apply the first to
$\varphi=\alpha$ (componentwise on $C$) and the second to $\varphi=\beta$:
\[
\langle\alpha_*\beta^*\xi,\eta\rangle_X=\langle\beta^*\xi,\alpha^*\eta\rangle_C=\langle\xi,\beta_*\alpha^*\eta\rangle_Y.
\]
The identity holds on integral classes, hence with any coefficients by linearity.
(a) Take $\mathcal C=\Gamma_w$: $\Gamma_w^*=w^*$ and $(\Gamma_w^t)^*=w_*=(w^{-1})^*$ (\Cref{prop:curves:transfer}(v)); so $\langle w^*x,w^*y\rangle=\langle x,(w^{-1})^*w^*y\rangle=\langle x,y\rangle$.
(b) Take $\mathcal C=\Gamma_\varphi$. (c) For $\eta\in H^1(X,\Z)$: $\eta(\operatorname{PD}_X\mathcal C^*\xi)=\langle\mathcal C^*\xi,\eta\rangle_X=\langle\xi,(\mathcal C^t)^*\eta\rangle_Y=((\mathcal C^t)^*\eta)(\operatorname{PD}_Y\xi)=\eta((\mathcal C^t)_*\operatorname{PD}_Y\xi)$,
by \Cref{prop:curves:correspondence-actions}(a) for $\mathcal C^t$. With \Cref{prop:curves:tate-module}(i) and \Cref{prop:curves:correspondence-actions}(b) this is the statement on $T_\ell J$. (d) is linearity.
\end{proof}

\begin{corollary}[twisted pairings and self-adjoint Hecke algebras]\label{cor:curves:twisted-pairing}
Let $X$ be a compact Riemann surface, $w$ an automorphism of $X$ with $w^2=\mathrm{id}$, and $\T$ a commutative $\Zl$-subalgebra of $\End_{\Zl}(H^1(X,\Zl))$ (so a $\Zl$-submodule, not merely a subring), each of whose elements is of the form $t=\sum_ic_i\mathcal C_i^*$
with $c_i\in\Zl$ and $\mathcal C_i$ correspondences from $X$ to $X$. Put $t^t=\sum_ic_i(\mathcal C_i^t)^*$ (well defined, being the adjoint $t^\dagger$ by \Cref{thm:curves:adjoint-transpose}), and
assume $t^t=w^*tw^*$ for all $t\in\T$. Put $[x,y]=\langle x,w^*y\rangle$. Then:
\begin{enumerate}[label=(\roman*)]
\item $[\ ,\ ]$ is a perfect alternating pairing on $H^1(X,\Zl)$;
\item every $t\in\T$ is self-adjoint for $[\ ,\ ]$: $[tx,y]=[x,ty]$;
\item every idempotent $e$ of $\T$ is self-adjoint, $H^1(X,\Zl)=eH^1\oplus(1-e)H^1$ is an orthogonal decomposition and $[\ ,\ ]$ is perfect on each summand; in particular, $\T$ being a finite
$\Zl$-algebra, $H^1(X,\Zl)=\bigoplus_{\mf m}H^1(X,\Zl)_{\mf m}$ over the maximal ideals of $\T$ is orthogonal for $[\ ,\ ]$, with $[\ ,\ ]$ perfect on each $H^1(X,\Zl)_{\mf m}$.
\end{enumerate}
\end{corollary}

\begin{proof}
By \Cref{thm:curves:adjoint-transpose}, $t^\dagger=t^t=w^*t(w^*)^{-1}$ for the cup product, which is perfect and antisymmetric on $H^1(X,\Zl)$ (\Cref{thm:curves:cup-product}). By
\Cref{thm:curves:adjoint-transpose}(a), $w^*$ is an isometry, and $w^*$ is an involution, so $(w^*)^\dagger=(w^*)^{-1}=w^*$. The proposition on adjoints and twisted pairings of
\Cref{ch:modules} (its parts on adjoints, idempotents and twisted pairings) now gives: $[\ ,\ ]$ is perfect, as the composite of a perfect pairing with the automorphism $w^*$; for
$t\in\T$, $[tx,y]=\langle x,t^\dagger w^*y\rangle=\langle x,w^*tw^*w^*y\rangle=[x,ty]$; and $[y,x]=\langle y,w^*x\rangle=-\langle w^*x,y\rangle=-\langle x,w^*y\rangle=-[x,y]$, so $[\ ,\ ]$ is antisymmetric, hence alternating
since $2$ is a nonzerodivisor in $\Zl$. For (iii), the same proposition of \Cref{ch:modules} shows that idempotents of a commutative ring of self-adjoint endomorphisms are self-adjoint, that the
decomposition is orthogonal and that the pairing is perfect on each summand. $\T$ is a $\Zl$-submodule of the finite free module $\End_{\Zl}(H^1(X,\Zl))$ by hypothesis, so it is finite and free over $\Zl$ (a submodule of a finite free module over the principal ideal domain $\Zl$), and
it is the product of its localisations $\T_{\mf m}$, cut out by idempotents $e_{\mf m}$ (\Cref{ch:modules}); $H^1(X,\Zl)_{\mf m}=e_{\mf m}H^1(X,\Zl)$.
\end{proof}

\begin{example}[adjoints of correspondences]\label{ex:curves:correspondence-examples}
\leavevmode
\begin{enumerate}[label=(\alph*)]
\item On $X=\C/\Lambda$ the automorphism $[-1]$ acts by $-1$ on $H^1$; it is an isometry and its own adjoint (\Cref{ex:curves:torus}).
\item For $\varphi\colon X\to E$ of degree $n$, the correspondence $\Gamma_\varphi^t$ followed by $\Gamma_\varphi$ acts on $H^1(E)$ by $\varphi_*\varphi^*=n$ and on $E=J(E)$ by $[n]$
(\Cref{prop:curves:degree-cup-product}(ii)). On $H^1(X,\Q)$ the operator $p=\frac1n\varphi^*\varphi_*$ satisfies $p^2=\frac1{n^2}\varphi^*(\varphi_*\varphi^*)\varphi_*=p$, has image $\varphi^*H^1(E,\Q)$
and is the identity there; and $p^\dagger=\frac1n\varphi_*^\dagger\varphi^{*\dagger}=\frac1n\varphi^*\varphi_*=p$ by \Cref{thm:curves:adjoint-transpose}(b) (for an antisymmetric pairing
$(A^\dagger)^\dagger=A$). So $p$ is the orthogonal projector onto $\varphi^*H^1(E,\Q)$ for the cup product; it is integral on $H^1(X,\Zl)$ for $\ell\nmid n$.
\item On a genus-one curve $E$, for an endomorphism $\psi$ (a nonconstant morphism fixing $O$) the adjoint of $\psi^*$ on $H^1(E,\Z)$ is $\psi_*=(\deg\psi)(\psi^*)^{-1}$ (on $H^1(E,\Q)$;
by \Cref{thm:curves:adjoint-transpose}(b) and \Cref{prop:curves:degree-cup-product}(ii), $\psi_*\psi^*=\deg\psi$). For $\psi=[m]$, $\psi^*=\psi_*=m$ is self-adjoint. For $E=\C/\Z[i]$ and
$\psi(z)=iz$: with $z=s+it$, $\psi^*s=-t$ and $\psi^*t=s$, so in the basis $ds$, $dt$ the matrix of $\psi^*$ is $\sm01{-1}0$ (which is the matrix $J$ of \Cref{ch:groups}),
a rotation of order $4$; its adjoint for the Gram matrix $G=\sm01{-1}0$ is $G^{-1}(\psi^*)^{\mathsf T}G=\sm0{-1}10=(\psi^*)^{-1}$, the inverse rotation. So a correspondence need not
be self-adjoint for the cup product.
\item (Forward pointer, \Cref{ch:moddeg}.) On $X_0(N)$, $T_n^t=w_NT_nw_N$ for all $n$, $U_p$ included. For $p\nmid N$, $w_N$ commutes with $T_p$, so $T_p^t$ acts as $T_p$ and $T_p^*$ is
self-adjoint for the cup product; for $p\mid N$, $U_p^t=w_NU_pw_N$ is in general not $U_p$, and $U_p^*$ is not self-adjoint for the cup product. This is why the paper composes the
cup product with $w$: by \Cref{cor:curves:twisted-pairing} every $t$ is self-adjoint for $\langle x,w_N^*y\rangle$.
\end{enumerate}
\end{example}

\begin{paperlink}\label{pl:curves:hecke}
In Section 4.1 of the paper, the pairing $\hat\delta_!$ ``is perfect and alternating on $H^1(X_0(M^\Sigma),\Zl)$, and the Hecke operators, $U_p$ included, are self-adjoint for it, because
$w_{M^\Sigma}T_nw_{M^\Sigma}^{-1}$ is the adjoint of $T_n$ for the cup product (this is what the composition with $w$ is for)''. With the input
$T_n^t=w_{M^\Sigma}T_nw_{M^\Sigma}$ from
\Cref{ch:moddeg}, \Cref{thm:curves:adjoint-transpose} says that $w^*T_n^*w^*$ is the adjoint of $T_n^*$, and \Cref{cor:curves:twisted-pairing}(i), (ii) is the sentence, with $X=X_0(M^\Sigma)$,
$w=w_{M^\Sigma}$ and $\T$ the image of the Hecke algebra tensored with $\Zl$ in $\End_{\Zl}(H^1(X_0(M^\Sigma),\Zl))$, a commutative $\Zl$-subalgebra. Every element of
$\T$ is a $\Zl$-combination of the $T_n^*$ ($\T$ is the $\Z$-span of the $T_n$, by the multiplicativity relations of \Cref{ch:hecke}), so the identity $t^t=w^*tw^*$
follows from the case $t=T_n$ by linearity; the factor $\frac12$ in $\hat\delta_!$ is a unit. The paragraph after Lemma lem:hecke and Step 2 of the proof of Theorem thm:quant:
``the idempotent $e_{\mf m}$ is self-adjoint for $\hat\delta_!$, so $H^1(X_0(N),\Zl)$ splits orthogonally into its local components and the pairing is perfect on each'' is
\Cref{cor:curves:twisted-pairing}(iii). This supplies the hypothesis of Lemma lem:ddt417 for the module $L=H^1(X_0(N),\Zl)_{\mf m}$:
a perfect $\Zl$-bilinear pairing for which $\T_{\mf m}$ is self-adjoint.
\end{paperlink}

\section{Exercises}\label{sec:curves:exercises}

\begin{exercise}\label{ex:curves:ex-singular}
(a) Show that $y^2=x^3+ax+b$ ($\operatorname{char}k\ne2,3$) is singular iff $4a^3+27b^2=0$, and that the singular point is then $(x_0,0)$ with $x_0$ the double root. (b) For the cusp
$y^2=x^3$ show that $t\mapsto(t^2,t^3)$ is a bijection $\A^1\to C$ which is not an isomorphism, by computing $\dim\mf m_P/\mf m_P^2$ at $P=(0,0)$. (c) Show that the Fermat curve
$X^d+Y^d+Z^d=0$ is singular at every point when $p=\operatorname{char}k$ divides $d$, and that for $d=2p$ it is not even irreducible.
\emph{Hints.} (a) A singular point has $f_y=2y=0$, and then $x$ is a common root of $x^3+ax+b$ and $3x^2+a$; the discriminant of $x^3+ax+b$ is $-(4a^3+27b^2)$.
(b) $f=y^2-x^3$ has no linear term at $(0,0)$. (c) $X^{2p}+Y^{2p}+Z^{2p}=(X^2+Y^2+Z^2)^p$.
\end{exercise}

\begin{exercise}\label{ex:curves:ex-smooth}
For \curve{121a1} ($y^2+xy+y=x^3+x^2-30x-76$) and \curve{20a1} ($y^2=x^3+x^2+4x+4$) show that the three partial derivatives of the homogenised equation have no common zero on the curve over
$\Qbar$. For which primes $p$ is the reduction of the model of \curve{20a1} modulo $p$ singular? (Answer: $2$ and $5$; at $2$ find the singular point.)
\emph{Hint.} The discriminant of the model of \curve{20a1} is $-2^8\cdot5^2=-6400$. Modulo $2$ the equation is $y^2=x^3+x^2$, with $f_y\equiv0$ and $f_x\equiv x^2$,
singular at $(0,0)$; modulo $5$, $x^3+x^2+4x+4\equiv(x+1)^2(x-1)$, singular at $(-1,0)$.
\end{exercise}

\begin{exercise}\label{ex:curves:ex-p1}
On $\PP^1$ show that $l(D)=\deg D+1$ for every divisor of degree $\ge-1$, by reducing to $D=n(\infty)$ with \Cref{prop:curves:picard}(i). Deduce Riemann-Roch for $\PP^1$ with $g=0$ and
$K=-2(\infty)$, including the cases $\deg D\le-2$. \emph{Hint.} $\mathcal L(n(\infty))$ is the space of polynomials in $x$ of degree at most $n$.
\end{exercise}

\begin{exercise}\label{ex:curves:ex-quartic-model}
Let $h$ be squarefree of degree $4$, $\operatorname{char}k\ne2$, and $C$ the smooth projective model of $y^2=h(x)$ (\Cref{cor:curves:rh-consequences}(c)). Show that $C$ has two points over $\infty$,
that $\operatorname{div}(dx/y)=0$, and that $g(C)=1$. If $h$ has no root in $k$ and its leading coefficient is not a square, $C$ may have no $k$-rational point at infinity: what does
\Cref{thm:curves:genus-one-weierstrass} need before it applies?
\end{exercise}

\begin{exercise}\label{ex:curves:ex-bases}
Let $C$ have genus one, $O\in C(k)$, and $x$, $y$ as in \Cref{thm:curves:genus-one-weierstrass}. Show that $\{x^iy^j:j\in\{0,1\},\ 2i+3j\le n\}$ is a basis of $\mathcal L(n(O))$ for every
$n\ge1$. Explain why the construction cannot be started on the conic $X^2+Y^2+Z^2=0$ over $\Q$ or on Selmer's cubic.
\end{exercise}

\begin{exercise}\label{ex:curves:ex-chord}
On a Weierstrass cubic with $O=[0:1:0]$, use \Cref{prop:curves:genus-one-pic} to show: (a) if a line meets $E$ in $P$, $Q$, $R$ (with multiplicity), then $\kappa(P)+\kappa(Q)+\kappa(R)=0$; (b) the
inverse of $P=(x_0,y_0)$ is $(x_0,-y_0-a_1x_0-a_3)$; (c) $P$ has order $2$ iff $2y_0+a_1x_0+a_3=0$.
\end{exercise}

\begin{exercise}\label{ex:curves:ex-xy}
For a Weierstrass cubic with $\operatorname{char}k\ne2,3$, show that $x\colon C\to\PP^1$ has degree $2$ and is ramified exactly at $O$ and at the three points with $2y+a_1x+a_3=0$, and verify
Riemann-Hurwitz. Show that $y\colon C\to\PP^1$ has degree $3$, that $e_y(O)=3$, and that the total ramification $\sum(e_y(P)-1)$ is $6$.
\end{exercise}

\begin{exercise}\label{ex:curves:ex-quartic}
(a) Show that the plane quartic $y^3=x^4-1$ (projectively $X^4-Y^3Z-Z^4=0$) is smooth in characteristic $0$, and compute its genus twice: by \Cref{thm:curves:plane-genus}, and by
Riemann-Hurwitz for $x\colon C\to\PP^1$ (degree $3$, totally ramified over the four roots of $x^4=1$ and over $\infty$). (b) For the Fermat quartic, show that $\omega_0$, $x\omega_0$, $y\omega_0$
($\omega_0=dx/f_y$) are regular, by checking the order at the four points at infinity. \emph{Answer to (a):} $g=3$ both times; Riemann-Hurwitz reads
$2g-2=3(-2)+5\cdot(3-1)=4$.
\end{exercise}

\begin{exercise}\label{ex:curves:ex-genus2}
Let $\varphi\colon X\to E$ be a nonconstant morphism over $\C$ from a curve of genus $2$ to a curve of genus $1$. Show that $\varphi$ has either two ramification points with $e=2$, or one with $e=3$,
and give the possible ramification patterns over the branch points when $\deg\varphi=2$. Show that a nonconstant morphism from a genus-one curve to $\PP^1$ has degree at least $2$.
\emph{Hint.} Riemann-Hurwitz gives $\sum(e_\varphi(P)-1)=2$. If $\deg\varphi=2$, a point with $e=2$ is its whole fibre, so there are exactly two branch points, each
with a single preimage. A morphism of degree one is an isomorphism (\Cref{cor:curves:degree-three-ways}(v)).
\end{exercise}

\begin{exercise}\label{ex:curves:ex-degrees}
The optimal curves \curve{37a1} and \curve{99a1} have modular degrees $2$ and $4$ (Cremona's tables), and $g(X_0(37))=2$, $g(X_0(99))=9$. Compute the total ramification of their modular
parametrisations. Show that for every optimal $E$ of conductor $N$, $\deg\varphi_E=1$ forces $g(X_0(N))=1$, and that $g(X_0(N))\ge2$ forces $\deg\varphi_E\ge2$.
\end{exercise}

\begin{exercise}\label{ex:curves:ex-basis-change}
For $X=\C/\Lambda$ and a new basis $(\omega_1',\omega_2')=(a\omega_1+b\omega_2,c\omega_1+d\omega_2)$ of $\Lambda$, $\gamma=\sm abcd\in\GL_2(\Z)$, compute the cup product of the new dual basis
$ds'$, $dt'$ and show $\langle ds',dt'\rangle=\det\gamma\cdot\langle ds,dt\rangle$. Deduce that $\operatorname{Im}(\bar\omega_1'\omega_2')>0$ iff $\det\gamma=1$, and that the convention ``$a\cdot b=+1$'' of
\Cref{ex:curves:torus} singles out the bases oriented by the complex structure.
\end{exercise}

\begin{exercise}\label{ex:curves:ex-weil}
Let $T$ be a free $\Zl$-module of rank two with a continuous action of a group $G$ and a perfect alternating $G$-equivariant pairing $e\colon T\times T\to\Zl(1)$. Show that $\det T=\chi_\ell$ (as a
character of $G$), that $x\mapsto e(x,-)$ is a $G$-isomorphism $T\to T^\vee(1)$, and hence $T^\vee\simeq T(-1)$ with $\det T^\vee=\chi_\ell^{-1}$. Apply this to $T=T_\ell E$ and explain why the
paper's lattice $M_{f,\lambda}=\varphi_E^*H^1(E,\Zl)$ is ``homothetic to the dual of $T_\ell E$'' and not to $T_\ell E$ as a Galois module, and why $\ad M_{f,\lambda}\simeq\ad T_\ell E$ anyway.
\end{exercise}

\begin{exercise}\label{ex:curves:ex-zeta3}
Let $E=\C/\Z[\zeta_3]$ and $\psi(z)=\zeta_3z$. Write the matrix of $\psi^*$ on $H^1(E,\Z)$ in the basis dual to $1$, $\zeta_3$, and of its adjoint for the cup product. Show that $\psi^*$ is an
isometry with $(\psi^*)^\dagger=(\psi^*)^2$ and that $\psi^*+(\psi^*)^\dagger=-1$. Which elements of $\Z[\psi^*]$ are self-adjoint?
\end{exercise}

\begin{exercise}\label{ex:curves:ex-twist-counter}
Let $L=\Zl^2$ with $\langle x,y\rangle=x_1y_2-x_2y_1$, $\ell$ odd. (a) For $w=\operatorname{diag}(2,\frac12)$, an isometry of infinite order, show that $[x,y]=\langle x,wy\rangle$ is perfect but not alternating. (b) For
$w=\operatorname{diag}(1,-1)$, an involution, show that $w$ is not an isometry and that $[\ ,\ ]$ is symmetric. (c) Conclude that both hypotheses on $w$ in \Cref{cor:curves:twisted-pairing}
(an isometric involution, which an automorphism of $X$ of order two provides) are needed for alternation.
\end{exercise}

\begin{exercise}\label{ex:curves:ex-gram}
Let $H=\Z^4$ with basis $a_1^\vee,a_2^\vee,b_1^\vee,b_2^\vee$ and the symplectic cup product of \Cref{thm:curves:cup-product}, and $M$ the span of $u=a_1^\vee$ and $v=nb_1^\vee+b_2^\vee$. Show
that $M$ is saturated, that the pairing on $M$ has ideal of values $n\Z$, compute $M^\perp$ and show $[H:M\oplus M^\perp]=n^2$, as \Cref{prop:curves:degree-cup-product}(iv) predicts for
$M=\varphi^*H^1(E,\Z)$ when $g(X)=2$. Do the same for $M'=\Z u\oplus\Z\cdot2v$ and compare with the non-saturated case.
\end{exercise}

\begin{exercise}\label{ex:curves:ex-component}
Let $V=\C^2$, $\Lambda=\Z[i]^2$, $W=\C$, $M=\Z[i]$ and $A(z_1,z_2)=2z_1+(1+i)z_2$. Show that $A$ induces a surjective homomorphism $\pi$ of complex tori, compute $A(\Lambda)$ and show that
$\ker\pi$ has two connected components. Replace $A$ by $z_1+(1+i)z_2$ and show that the kernel becomes connected.
\end{exercise}

\begin{exercise}\label{ex:curves:ex-frobenius}
Let $E\colon y^2=x^3-x$ over $\F_3$. Count $E(\F_3)$, the fixed points of $\pi_3$, and deduce $a_3$ and the roots $\alpha$, $\beta$ of $T^2-a_3T+3$. Predict $\#E(\F_9)$ from
\Cref{thm:curves:weil-pairing-etale}(d) and verify it by counting the fixed points of $\pi_3^2$ (answers: $4$, $a_3=0$, $16$). Show that $\pi_3$ is a bijection on $E(\bar\F_3)$ of degree $3$,
and explain why \Cref{thm:curves:riemann-hurwitz} does not apply to it.
\end{exercise}

\begin{exercise}\label{ex:curves:ex-residue}
(a) On $\PP^1$ compute all residues of $\omega=dz/(z^3-z)$, including at $\infty$, check \Cref{thm:curves:residue}, and deduce the partial fraction decomposition of $1/(z^3-z)$. (b) Show that on a
compact Riemann surface no meromorphic $1$-form has a single simple pole and no other pole. (c) On a Weierstrass cubic, show that $x\,dx/(2y+a_1x+a_3)$ has a double pole at $O$ and no other pole,
and deduce that its residue at $O$ is $0$.
\end{exercise}

\section{Notes and references}\label{sec:curves:notes}

The algebraic side (\Cref{sec:curves:varieties,sec:curves:plane,sec:curves:function-fields,sec:curves:divisors,sec:curves:riemann-roch,sec:curves:morphisms,sec:curves:hurwitz}) follows
Silverman \cite{SilvermanAEC}, Chapter I (Sections 1 to 3: varieties, Galois action, maps) and Chapter II (Sections 1 to 5: curves, divisors, differentials, Riemann-Roch), and Chapter III,
Section 3 for Weierstrass equations. The numbered results of Silverman cited above (Propositions II.1.1, II.2.1, II.2.6, II.2.11, II.4.2, II.4.3, II.5.8, Theorems II.2.3, II.2.4, II.5.4,
Corollary II.2.12) are cited from the second edition as we recall them; the reader should treat the numbers as pointers and check them against the printed edition. Fulton \cite{Fulton}
is the gentlest complete account of plane curves, smoothness, intersection multiplicities and Bezout (Chapters 3, 5), the smooth model (Chapter 7) and Riemann-Roch (Chapter 8); Hartshorne
\cite{Hartshorne} (Chapter I, Sections 1 to 7; Chapter II, Section 6; Chapter IV, Sections 1 and 2; Appendix B) gives the scheme-theoretic versions; Corollary I.6.12, Proposition II.6.9, Theorem
IV.1.3 and Corollary IV.2.4 are cited as recalled. The finiteness of integral closure is in Atiyah and Macdonald \cite{AtiyahMacdonald}, Chapter 5 (chapter not verified).

The analytic side (\Cref{sec:curves:riemann-surfaces,sec:curves:differentials-analytic,sec:curves:cohomology,sec:curves:jacobian}) follows Miranda \cite{Miranda} (Riemann surfaces and
holomorphic maps in Chapters I and II, differentials and residues in Chapter IV, Riemann-Roch and Riemann existence in Chapters V to VII, Abel's theorem, Jacobi inversion and the
bilinear relations in Chapter VIII) and Forster \cite{Forster81} (Sections 1, 2 and 10 for the local theory, 14 to 17 for Riemann-Roch and its consequences, 19 to 21 for de Rham, periods and
Abel-Jacobi). Griffiths and Harris \cite{GriffithsHarris78}, Chapter 2, Section 2, and Farkas and Kra \cite{FarkasKra92}, Chapter III, give the period relations and the Jacobian with the
symplectic picture; Hatcher \cite{Hatcher02} is the reference for singular (co)homology, the universal coefficient theorem, transfers of coverings (Section 3.G) and Poincare duality (Section
3.3). All chapter and section numbers in this paragraph are recalled pointers, not checked against the books.

The Jacobian as an abelian variety, its universal property, its torsion and Tate modules and the Weil pairing are in Milne's two articles in Cornell-Silverman
\cite{MilneJacobian86,MilneAbelian86} and in Mumford \cite{Mumford70} (Sections 18 to 24); the Picard functor is in \cite[Chapter 9]{BLR}. Etale cohomology and the comparison theorem are in
Milne's notes \cite{MilneEtale}. Correspondences and the Hecke correspondences on modular curves are in Shimura \cite{Shimura71}, Chapter 7, and Diamond and Shurman \cite{DiamondShurman}, Section 6.3;
the self-adjointness argument of \Cref{cor:curves:twisted-pairing} is the one used by Darmon, Diamond and Taylor \cite[Section 4.4]{DDT}; Diamond, Flach and Guo state the corresponding facts without
proof \cite[Sections 1.6.1 and 1.8.2]{DFG04} (checked: Proposition 1.3 there makes the Hecke action compatible with their duality morphisms, and Section 1.8.2 says that ``the Hecke
operators are self-adjoint with respect to the pairing of (25)''). The saturation criterion and the cup product on $\varphi^*H^1(E)$ (\Cref{prop:curves:degree-cup-product}) are elementary once
\Cref{thm:curves:cup-product,lem:curves:tori} are available; compare Agashe, Ribet and Stein \cite[Section 3]{ARS12} (section not verified) for the modular degree and optimal quotients.
Section numbers for Milne, Mumford, Shimura and Diamond-Shurman are recalled pointers. The genera and modular degrees quoted were computed with PARI/GP \cite{PARI} and Cremona's tables
\cite{Cremona97}.

Three choices differ from some sources. The cup product is $\int\theta\wedge\eta$ with the complex orientation; the sign of Poincare duality and of the Weil pairing then differ from some books, but the
paper uses only sign-invariant properties. The adjoint of an operator is $A^\dagger$, the star being reserved for pullbacks. And the transfer $\varphi^!$ and the pushforward $\varphi_*$ on $H^1$ are
defined topologically, with the trace of forms as a computational device (\Cref{prop:curves:transfer}), rather than through sheaf cohomology.
